% Part F contribution.  Local labels FF.* are to be renumbered by the editor.
% AI-produced by GPT-6 Astra (OpenAI, ChatGPT Pro mode). No priority claimed.
\section{Schemes, relative topoi, and the proposed field criteria}
\label{sec:FF}

Throughout this part, a relative geometric morphism includes its specified
commutation isomorphism with the structural morphisms to the base topos.
The adjective \emph{pinned} has a separate meaning: the induced map of
topological spaces sends closed points to closed points
\citep[Definition~2.19]{FF_CHW2024}. We retain both pieces of data.
Write \(\mathcal R(k)\) for the reconstruction property in the question.
All assertions below concern characteristic zero except where a
finite field or an arithmetic base is explicitly specified.

\subsection{Points, sections, and the qualification in scan 52}

\begin{claim}{FF.1}{Proved}\label{clm:FF1}
Let \(k\) have characteristic zero and let \(X/k\) be a finite-type
scheme. The isomorphism classes of relative geometric morphisms
\(BG_k\to X_{\mathrm{et}}\) are the disjoint union, over the points
\(x\in X\), of the continuous sections of
\[
                 G_{k(x)}\longrightarrow G_k,
\]
modulo conjugation by its kernel, with compatible separable closures
understood. A point for which the displayed map is not surjective
contributes no section. The closed-point summands are exactly
\(X(k)\). In particular, for an integral curve the other summand is
the set of birational sections at its generic point.

Consequently, for every characteristic-zero field the assertion in
cote~149, scan~52, read with \emph{all} relative topos morphisms, is
false, even with universal homeomorphisms inverted.
\end{claim}
\begin{proof}
Pull back a geometric morphism along the canonical point
\(\mathrm{Set}\to BG_k\). This gives a geometric point of
\(X_{\mathrm{et}}\), hence a separably closed point above a unique
\(x\in X\). Descent from this point to \(BG_k\) is a continuous
action of \(G_k\) on its fibre functor. The automorphism group of
that functor is \(G_{k(x)}\). The relative commutation isomorphism
makes its composite to \(G_k\) the identity. Changing the compatible
geometric point acts by the geometric kernel, which gives exactly the
stated ambiguity. Conversely, such an action supplies the descent
datum and hence the geometric morphism. This is the point description
of the small \'etale topos, applied with its continuous automorphism
group; it involves the residue field at \(x\), rather than only
\(\pi_1(X)\).

If \(x\) is closed, \(k(x)/k\) is finite separable. Its absolute
Galois group maps to an open subgroup of \(G_k\) of index
\([k(x):k]\). A right inverse forces this index to be one. Conversely
a rational point gives a relative morphism, uniquely up to the
specified relative isomorphism. This also follows from
\citet[Proposition~3.6]{FF_CHW2024}.

For the counterexample put \(X=\mathbb A^1_k\). The henselian
discrete valuation field \(H=\operatorname{Frac}(k[t]^h_{(t)})\)
has residue field \(k\). In characteristic zero its decomposition
sequence splits: in a separable closure obtained by adjoining
separable constants and compatible roots of \(t\), act on the
constants and fix those roots. Thus
\(G_k\to G_H\to G_{k(t)}\) is a section over \(G_k\).
It defines a relative topos morphism whose image in \(\mathbb A^1_k\)
is the generic point. Every scheme morphism
\(\operatorname{Spec}k\to\mathbb A^1_k\) has rational, hence closed,
image. Inverting universal homeomorphisms does not change this
distinction; in particular the source field is already absolutely
weakly normal. This is the tangential construction of AP-F.1 and
AP-F.17, also illustrated in \citet[Example~2.30]{FF_CHW2024}.
\end{proof}


\begin{claim}{FF.2}{Known, conditional as stated}\label{clm:FF2}
Let \(k/\mathbb Q\) be finitely generated and \(K/k\) a regular
one-variable function field. Under conjecture~a) in cote~149,
scans~126--127, each section \(G_k\to G_K\) factors through a
decomposition group \(D_v\) of a nontrivial valuation trivial on
\(k\), with residue field finite over \(k\). For an actual section
over \(G_k\), that residue field is \(k\) itself. Its centre lies
on the smooth proper model of \(K/k\), and need not lie on a
specified open curve with that function field.
\end{claim}
\begin{proof}
The manuscript expressly requires the valuation ring to contain
\(k\) and its residue field to be algebraic over \(k\). On the
smooth proper curve, the valuative criterion gives a centre.
The valuation is nontrivial, so its centre is a closed point;
the dominating valuation ring is its discrete valuation ring.
Its residue field is therefore a finite extension \(k(v)/k\).
The decomposition sequence is
\[
       1\longrightarrow I_v\longrightarrow D_v
         \longrightarrow G_{k(v)}\longrightarrow1.
\]
The image of \(D_v\to G_k\) is the open subgroup associated to
\(k(v)/k\). A section of \(G_K\to G_k\) contained in \(D_v\)
forces that image to be all of \(G_k\), so \(k(v)=k\).
The factorization itself is the conjectural assertion, not a theorem
proved here. On scan~126 the uncertain reading is
\(I_V\) versus \(N_V\); the subsequent discussion of its normalizer
and scan~127's explicit factorization support the inertia reading.
Our conditional statement uses the explicit factorization on
scan~127 and does not silently assume that every normalizer equals
the decomposition group.
\end{proof}

\subsection{Finite fields and an arithmetic base}

\begin{claim}{FF.3}{Refuted}\label{clm:FF3}
Closed-point-preserving relative \'etale-topos morphisms over
\(\mathbb F_q\) need not be induced by morphisms after universal
homeomorphisms have been inverted. This already fails for
\(X=Y=\mathbb G_{m,\mathbb F_q}\). The corresponding assertion
for all finite-type schemes over \(\operatorname{Spec}\mathbb Z\)
is also false.
\end{claim}
\begin{proof}
Identify \(G_{\mathbb F_q}\) with \(\widehat{\mathbb Z}\), using
the arithmetic Frobenius. Each \(z\) in this central group gives a
2-automorphism of the structural functor
\(X_{\mathrm{et}}\to BG_{\mathbb F_q}\). Twist the relative
structure of the identity of \(X_{\mathrm{et}}\) by that
2-automorphism, choosing the sign so that its geometric-point
function is
\[
          a\longmapsto\operatorname{Frob}_q^z(a)
          \qquad(a\in\overline{\mathbb F}_q^{\times}).
\]
The underlying map of spaces is the identity, so this relative
automorphism is pinned. The point function follows either directly
by transporting the relative base identification, or from the
finite-extension point description in Proposition~3.6 of CHW.

Choose \(z\in\widehat{\mathbb Z}=\prod_\ell\mathbb Z_\ell\)
whose 2-adic component is zero and whose other components are one.
This is not an integer: its 2-adic component would force that integer
to be zero, contradicting its 3-adic component. Suppose our relative
morphism came from
\(f:X^{\mathrm{perf}}\to X\). Every unit in
\(\mathbb F_q[T^{\pm1/p^\infty}]\) is
\(cT^{m/p^r}\), with \(c\in\mathbb F_q^\times\).
The induced geometric fundamental-group homomorphism must be an
isomorphism, because our topos morphism is one. On
the maximal prime-to-\(p\) quotient
\(\pi_1(X_{\overline{\mathbb F}_q})^{(p')}=\prod_{\ell\ne p}\mathbb Z_\ell\)
it is multiplication by \(m/p^r\). Hence
\(m/p^r=\pm p^s\) for an integer \(s\). Evaluation at 1 gives
\(c=1\).

The negative sign is impossible: composing with an integral power
of Frobenius would make inversion on the multiplicative group,
extended by zero, a field automorphism. Additivity at \(u+1\)
would then give \(u^2+u+1=0\) for every
\(u\ne0,-1\), which is false in the infinite algebraic closure.
For the positive sign, write \(q=p^a\). Equality of the field
automorphisms gives \(az=s\) in \(\widehat{\mathbb Z}\), hence
\(z=s/a\in\mathbb Q\cap\widehat{\mathbb Z}=\mathbb Z\), another
contradiction. This proves nonalgebraicity even after perfection.

For the last assertion take \(q=p\) and regard \(X\) as a vertical
\(\mathbb Z\)-scheme. The closed immersion
\(\operatorname{Spec}\mathbb F_p\hookrightarrow
\operatorname{Spec}\mathbb Z\) induces a geometric embedding of
small \'etale topoi. Postcomposition with a geometric embedding is
fully faithful on relative factorization data. Thus a relative
identification over \((\operatorname{Spec}\mathbb Z)_{\mathrm{et}}\)
would give the same identification over \(BG_{\mathbb F_p}\), which
we have excluded. A morphism of these vertical schemes over
\(\mathbb Z\) is automatically a morphism over \(\mathbb F_p\).
\end{proof}

The distinction in the last proof addresses the positive-characteristic
caution attached to scan~52. Absolute Frobenius and the identity have
isomorphic \emph{absolute} \'etale-topos morphisms. Their relative
commutation isomorphisms retain different data. Inverting universal
homeomorphisms makes Frobenius invertible; it does not identify it with
the identity in the category of schemes.

\begin{claim}{FF.4}{Known}\label{clm:FF4}
The finite-field counterexample does not invalidate the following
restricted positive statements. CHW Theorem~0.1 gives reconstruction
over infinite finitely generated fields, with absolute weak
normalization. Carlson--Stix Theorem~D, numbered~6.12 in the version
dated 20 September 2026, gives relative reconstruction when the base
is noetherian, Nagata and geometrically unibranch, each irreducible
component has a sub-\(p\)-adic function field, the schemes are
componentwise universally open and of finite type over the base,
and the topos morphisms send constructible subsets to constructible
subsets. Universal homeomorphisms are inverted.
\end{claim}
\begin{proof}
These are the hypotheses and conclusions of the opened statements
\citep[Theorem~0.1]{FF_CHW2024}
and \citep[Theorem~D, Theorem~6.12]{FF_CS2026}.
Carlson--Stix call their morphisms \emph{chivalrous}; their Remark~2.1
shows that these are pinned on fibres. This is a stronger and
different relative condition from merely preserving the closed
points of the total space. Their Corollary~E includes the stated
normal-source Hom conclusion. Over \(\operatorname{Spec}\mathbb Z\),
the generic field is sub-\(p\)-adic, but vertical schemes such as
those in Claim~\ref{clm:FF3} fail universal openness. Thus the
restricted theorem does not assert the statement refuted there.
\end{proof}

\subsection{The birational automorphism experiment}

\begin{claim}{FF.5}{Known}\label{clm:FF5}
Let \(K_0\) be a number field and \(X/K_0\) a smooth proper
geometrically connected curve, with \(K=K_0(X)\). Every continuous
outer automorphism of \(G_K\) over \(G_{K_0}\), modulo inner
conjugation by the geometric kernel, is induced by a unique
\(K_0\)-automorphism of \(X\). Preservation of the closed-point
inertia and decomposition groups is a consequence, and need not be
assumed. In particular the experiment in cote~149, scan~52, has a
positive answer for the genus at least two curves asked about.
\end{claim}
\begin{proof}
A number field is sub-\(p\)-adic. Mochizuki's Theorem~B, equivalently
Theorem~17.1, states the birational Hom theorem for regular function
fields over a sub-\(p\)-adic base: base-linear field embeddings
correspond bijectively, contravariantly, to continuous open
homomorphisms of their \emph{full} absolute Galois groups over the
base group, modulo inner conjugation by the geometric kernel
\citep[Theorem~17.1, author pp.~87--88]{FF_Mochizuki1999}.
The title of that paper does not replace these full groups by
pro-\(p\) quotients in this theorem.

Apply the theorem to the given automorphism and to its inverse.
Its uniqueness and functoriality imply that the resulting field
embeddings are inverse, giving a \(K_0\)-automorphism of \(K\).
It extends uniquely to the smooth proper curve: the local rings at
its closed points are precisely the discrete valuation rings of
\(K/K_0\) with finite residue field, and a birational map between
regular proper curves extends across those points. The resulting
scheme automorphism transports their valuation, inertia and
decomposition groups. This proves the asserted redundancy of the
loop hypothesis. The birational argument itself does not require
the genus bound.
\end{proof}

\subsection{Two proposed field criteria}

\begin{claim}{FF.6}{Proved}\label{clm:FF6}
If \(\mathcal R(k)\) holds in characteristic zero, then for every
finite extension \(E/k\),
\[
             C_{\operatorname{Aut}(\bar k)}(G_E)=1.
\]
In particular every open subgroup of \(G_k\) has trivial centre.
This proves the necessary direction of candidate F3(a).
\end{claim}
\begin{proof}
Finite-extension ascent of reconstruction follows by taking the
slice over the finite \'etale object \(\operatorname{Spec}E\):
an \(E\)-relative map is a \(k\)-relative map with that factorization,
and uniqueness over \(k\) preserves the factorization. Thus
\(\mathcal R(E)\) holds.

Let \(\sigma\in\operatorname{Aut}(\bar k)\) centralize \(G_E\).
It preserves the fixed field \(E\). Its action on coefficients,
with the variable fixed, yields a semilinear automorphism of
\(\mathbb G_{m,E}\), whose topos morphism is relative over
\(BG_E\) because the induced conjugation on \(G_E\) is the
identity. With this relative structure its point function is
\(a\mapsto\sigma(a)\). This includes the case where \(\sigma\)
fixes \(E\): then its nontriviality is in the relative structure.
Reconstruction would make the map \(cT^n\). Its being a topos
equivalence forces \(n=\pm1\), and its value at 1 forces \(c=1\).
Inversion cannot be a field automorphism by the additivity calculation
in Claim~\ref{clm:FF3}. Hence \(\sigma(a)=a\) for all nonzero
\(a\), proving the claim. This is the explicit ghost test of
AP-F.39--F.40.
\end{proof}

\begin{claim}{FF.7}{Refuted}\label{clm:FF7}
Candidate F3(b) is false. Reconstruction need not imply the Isom
form of the Grothendieck conjecture for hyperbolic curves of genus
zero, even for a pair of four-punctured projective lines over the
base field itself. This remains a counterexample if the group
isomorphisms are required to preserve the peripheral conjugacy
classes.
\end{claim}
\begin{proof}
Put \(k=\overline{\mathbb Q}(t)\). Reconstruction over this field
is AP-F.45; the dependent argument is checked in
Claim~\ref{clm:FFaudit} below. Consider
\[
 U_2=\mathbb P^1_k\setminus\{0,1,\infty,2\},\qquad
 U_3=\mathbb P^1_k\setminus\{0,1,\infty,3\}.
\]
An isomorphism extends to their smooth proper completions, hence
is a fractional linear transformation carrying one four-element
set to the other. The cross-ratio orbits are respectively
\[
 \{2,\tfrac12,-1\},\qquad
 \{3,\tfrac13,-2,-\tfrac12,\tfrac32,\tfrac23\}.
\]
They are disjoint. The two curves are therefore not isomorphic
over \(k\), nor over an algebraically closed characteristic-zero
extension. The exact rational computation is reproduced in
\texttt{F\_cross\_ratio.py}.

On the other hand these are constant curves over
\(\overline{\mathbb Q}\). The characteristic-zero K\"unneth
isomorphism of SGA~1, Expos\'e~XIII, Proposition~4.6, applies to
this smooth curve and to \(\operatorname{Spec}k\). Resolution in
dimension one suffices for its hypothesis. Together with Expos\'e
XIII, Corollary~2.12, it gives
\[
 \pi_1(U_2)\simeq\widehat F_3\times G_k
          \simeq\pi_1(U_3)
\]
as groups over \(G_k\) \citep{FF_SGA1}. The projection is the
arithmetic structural map. A homeomorphism of the complex
four-punctured spheres gives an isomorphism of the free profinite
factors preserving their peripheral classes; the product with
the identity of \(G_k\) supplies the required relative isomorphism.
Thus the proposed equivalence already fails in its necessary
direction. No inference from a full topos to its fundamental
group is being reversed here.
\end{proof}

\begin{claim}{FF.8}{Proved}\label{clm:FF8}
For descent along a transcendental field extension, a commuting
square of small \'etale topoi with a finitely generated model is
not by itself a scalar-descent datum for a topos morphism.
Moreover, the usual model systems of a finite-type scheme over
an arbitrary field are not generally pro-objects in the category
whose transition morphisms are required to be pinned.
\end{claim}
\begin{proof}
Take \(k=\mathbb C\), a finitely generated \(F\subset\mathbb C\),
and a nonidentity \(\sigma\in\operatorname{Aut}(\mathbb C/F)\).
Act on the coefficients of \(\mathbb G_{m,\mathbb C}\) and fix
the variable. The induced topos automorphism is pinned and, since
\(BG_{\mathbb C}=\mathrm{Set}\), relative over \(\mathbb C\).
It commutes with the projection to the model over \(F\). Yet its
point function \(a\mapsto\sigma(a)\) cannot be given by a
\(\mathbb C\)-morphism of \(\mathbb G_m\), by the preceding
unit and additivity argument. Thus this commuting square does not
mean that the automorphism is the scalar extension of the identity
on the model.

For the transition issue use \(F\subset F(u)\) with \(u\)
transcendental. The closed point \(T=u\) of
\(\mathbb A^1_{F(u)}\) maps to the generic point of
\(\mathbb A^1_F\). Hence the transition is not pinned.
Likewise, the small \'etale topos of the scalar extension is
not in general the 2-fibre product of the two small topoi over
\(BG_F\): over a separably closed extension that fibre product
only makes the algebraic constant extension of the original
geometric topos, whereas transcendental scalar extension supplies
additional closed points. Candidate F3(c) must specify a genuine
descent category; the displayed commuting-square interpretation
does not suffice.
\end{proof}

\subsection{The reconstruction chain and a normal-extension enlargement}

\begin{claim}{FF.9}{Proved}\label{clm:FFcore}
Let \(D/\mathbb Q\) be finitely generated, and let \(k/D\) be
an algebraic Galois extension containing all roots of unity.
If \(k\) is Hilbertian, then every finitely generated extension
\(L/k\) has \(\mathcal R(L)\).
\end{claim}
\begin{proof}
We give the common argument underlying AP-F.43 and AP-F.49,
including the enlargement of their hypotheses. For a field \(E\)
put \(D(E^\times)=\bigcap_{n\ge1}(E^\times)^n\).

\emph{The scalar kernel.}
For every finitely generated \(M/\mathbb Q\),
\(D(M^\times)=1\). Divisors on a normal proper model put a
divisible element and each of its roots in the number field of
constants; finite valuations and then the finitely generated unit
group of that number field prove the assertion.
If \(E/k\) is finite and \(a\in D(E^\times)\), choose a finite
Galois \(M/D\) with \(a\in M\) and \(E\subset Mk\).
For \(\sigma\in\operatorname{Gal}(Mk/k)\),
\(\tau\in\operatorname{Gal}(Mk/M)\), and \(z_n\in E\) with
\(z_n^n=a\), the automorphisms \(\sigma,\tau\) commute: check
this on \(M\) and \(k\), both of which they preserve.
Moreover \(\tau z_n/z_n\in\mu_n\subset k\) is fixed by
\(\sigma\). Therefore \(\sigma z_n/z_n\in M\), and
\(\sigma a/a\in D(M^\times)=1\). Thus \(a\in k\).
For a finitely generated \(E/k\), the proper-model divisor
argument first puts \(a\) and all its roots in the finite relative
algebraic closure of \(k\) in \(E\). We have proved
\[
             D(E^\times)\subset k^\times
             \quad(E/k\text{ finitely generated}).       \tag{FF.9.1}
\]
Equality and radical closedness of \(k\) are not required.

\emph{Detection of arithmetic finite quotients.}
Every finitely generated extension of a Hilbertian field is
Hilbertian. For the transcendental case we use
\citet[Theorem~1.2, Corollary~3.2]{FF_BarySorokerParan2013};
finite extensions preserve it by
\citet[Corollary~3.2.2]{FF_Serre2007}. Let \(Y\to U\) be a
connected finite \'etale \(G\)-torsor over a smooth geometrically
connected curve over \(L\). A rational function of arbitrarily
large available degree \(d\) on its proper completion makes
\(Y\to\mathbb P^1_L\) finite \'etale after deleting finitely
many values. Hilbert specialization of the entire connected cover
gives a residue field of degree \(d|G|\). Its quotient gives
\(x\in U\) with \([L(x):L]=d\) and
\[
                         G_{L(x)}\twoheadrightarrow G.   \tag{FF.9.2}
\]
This retains the arithmetic constant quotient; geometric
connectedness of \(Y\) is not assumed. The exact specialization
statement is \citet[Proposition~3.3.1]{FF_Serre2007}.
For higher-dimensional smooth quasiprojective \(U\), successive
hyperplane sections preserve every geometric component of the
cover and its arithmetic permutation. Poonen--Slavov Theorem~1.5
applies component by component because the covering map has
zero-dimensional nonempty fibres \citep{FF_PoonenSlavov2020}.
Combining with the smooth Bertini open gives a smooth curve,
retaining any prescribed nonempty open restriction. Thus (FF.9.2)
also gives the required finite-quotient detection in that case.

\emph{One punctured-line calculation.}
Consider
\(U=\operatorname{Spec}L[T,1/(P_1\cdots P_r)]\), with distinct
monic irreducible \(P_j\), and choose \(x_0\in U(L)\).
A compatible Kummer class with actual values at all finite-extension
points, normalized by its actual value \(b_0\) at \(x_0\), is
\[
 \kappa(b_0)+\sum_j\lambda_j
             \bigl([P_j(T)]-[P_j(x_0)]\bigr),
                  \qquad\lambda_j\in\widehat{\mathbb Z}. \tag{FF.9.3}
\]
This uses the vanishing of the Picard group of this localization
of a principal ideal domain.

Choose \(F/\mathbb Q\) finitely generated, containing \(D\)
and the finite displayed data, with \(L=Fk\); then \(L/F\)
is algebraic Galois. Hilbertianity gives a nontrivial finite
extension of \(L\). Descend it, and take a finite Galois extension
\(M/\mathbb Q(t_1,\ldots,t_s)\) containing \(F\), these
extension data and the roots of all \(P_j\). Choose
\(1\ne\sigma\in\operatorname{Gal}(Mk/L)\). It preserves \(M\)
and acts nontrivially on it. We require discrete valuations
\(v_i\) of \(M\), with mutually disjoint pairs
\(\{v_i,\sigma^{-1}v_i\}\), such that the relevant nonzero coefficients and
distinct root differences are units.

Here are the two exhaustive cases. If \(\sigma\) moves the
number field \(C_M=M\cap\overline{\mathbb Q}\), choose good
rational primes splitting completely in \(C_M\). Above each,
a Gauss valuation on the rational function field has an unramified
prolongation to a suitable normal model of \(M\), outside a finite
exceptional set. The restrictions of \(v_i\) and
\(\sigma^{-1}v_i\) to \(C_M\) differ. When \(s=0\), take
the residue prime sufficiently large to avoid finitely many roots.
When \(s>0\), the residue field is infinite. If \(\sigma\)
fixes \(C_M\), it is a nonidentity automorphism of the
positive-dimensional function field over \(C_M\). Choose a proper
normal \(\sigma\)-equivariant model, for example by normalizing
a proper model of the fixed field. Infinitely many prime divisors
are moved. Indeed, if the moved divisors formed a finite set \(S\),
choose \(f\) with \(\sigma f\ne f\), and a fractional linear
transform \(h\) of \(f\) for which both \(h\) and \(1+h\)
are units at every member of \(S\). Such a transform exists
because \(C_M\) is infinite and only finitely many residue values
are excluded. Their divisors are fixed, so properness gives
\(\sigma h=ah\) and \(\sigma(1+h)=b(1+h)\), with
\(a,b\in C_M^\times\). Comparing these identities forces
\(a=b=1\), contradicting \(\sigma f\ne f\). Avoiding the
finite bad divisors and the previously selected finite orbits
now gives the required disjoint moved pairs. These are the arithmetic and geometric divisor
arguments of AP-F.43 and AP-F.45 respectively.

Weak approximation at this finite set of distinct discrete
valuations chooses \(\alpha\in M\) equal to a root of \(P_i\)
plus a uniformizer modulo the square of the \(v_i\)-ideal, and
avoiding every root modulo the \(\sigma^{-1}v_i\)-ideal.
The unit derivatives and root differences give
\[
 v_i(P_j(\alpha))=\delta_{ij},\qquad
 v_i(P_j(\sigma\alpha))=0.                               \tag{FF.9.4}
\]
Set \(R_j=P_j(\sigma\alpha)/P_j(\alpha)\). The actual value
\(b_\alpha\in Mk\) lies in one field \(F'=MS_0\), where
\(S_0/D\) is finite inside \(k\). This field is finite over
\(M\), \(\sigma\)-stable, and fixed independently of \(n\).
If \(z^n=a\in F'\) with \(z\in Mk\), the commutation
calculation above gives \(\sigma z/z\in F'\). Apply it to
the evaluated relation (FF.9.3). For integer representatives
\(a_{j,n}\) of \(\lambda_j\bmod n\) it gives
\[
 \frac{\sigma b_\alpha}{b_\alpha}
       \prod_jR_j^{-a_{j,n}}\in(F'^\times)^n.
\]
Apply the integer-valued map \(v_i\circ N_{F'/M}\), for every
\(n\). We obtain
\[
 [F':M]\lambda_i
       =-(v_i\circ N_{F'/M})(\sigma b_\alpha/b_\alpha)
                     \in\mathbb Z\subset\widehat{\mathbb Z}.
                                                               \tag{FF.9.5}
\]
An element of \(\widehat{\mathbb Z}\) with a nonzero integer
multiple in \(\mathbb Z\) is itself in \(\mathbb Z\): the
corresponding rational number is integral at every prime. Thus
all puncture exponents are integral. No valuation on the infinite
field \(Mk\) was asserted to have discrete value group.

\emph{The Albanese step.}
Let \(X/L\) be regular, affine and connected, with
\(x_0\in X(L)\), and let \(\xi\) be the Kummer class pulled
back by a pinned relative map to \(\mathbb G_m\). It has actual
values at every finite-extension point by the point-source
description. The generalized Albanese \(A\), based at \(x_0\),
identifies normalized units with \(\operatorname{Hom}(A,\mathbb G_m)\).
The comparison of \citet[Section~3, Proposition~4.1,
Corollary~4.2, Proposition~4.3]{FF_SpiessSzamuely2003} identifies
geometric \(H^1\) with the dual of \(T_\ell A\): the finite
kernel in the abelian fundamental-group comparison contributes
no map to the torsion-free group \(\mathbb Z_\ell(1)\).
The required smooth projective compactification exists in
characteristic zero.

Descend the finite geometric data to \(F\) as above and put
\(H=G_L\triangleleft G_F\), \(W=(V_\ell A)_H\). Because
\(\mu_\infty\subset L\), the target Tate twist is trivial
on \(H\). Projected point-Kummer cocycles are homomorphisms
\(c_x:G_E\to W\), for finite \(E/L\). Their images span \(W\).
Indeed, a nonzero annihilating functional, scaled to be integral,
defines a nonzero finite quotient of \(\pi_1(X)\) which is zero
on the base-point section and every closed-point decomposition
group. This contradicts (FF.9.2).

Choose finitely many spanning points. Their coordinates, their
actual values, and the normalization descend to a single
finitely generated field \(M_1/\mathbb Q\) containing \(F\).
Restriction to \(G_{M_1k}\) preserves their rational spans,
because it is restriction to a subgroup of finite index.
The normalized geometric functional
\(\lambda:T_\ell A\to\mathbb Z_\ell(1)\) satisfies
\[
 \lambda(c_x(ghg^{-1}))
       =\chi_\ell(g)\lambda(c_x(h))
 \quad(g\in G_{M_1},\ h\in G_{M_1k}).
\]
Here the point and its value are \(M_1\)-rational,
\(G_{M_1k}\) is normal in \(G_{M_1}\), and the cocycle
coboundary disappears in \(H\)-coinvariants. Spanning proves
that \(\lambda\) is \(G_{M_1}\)-equivariant.
The integral Tate-Hom theorem for 1-motives over finitely generated
fields now puts it in
\(\operatorname{Hom}_{M_1}(A,\mathbb G_m)\otimes\mathbb Z_\ell\).
We use the exact integral statement recorded in
\citet[Section~1.4]{FF_Jossen2014}; its application is over
\(M_1\), not over the infinite field \(L\).

After a finite extension splitting the boundary, every normalized
geometric unit is defined over the same finite model: its conjugates
have equal divisors and value 1 at \(x_0\), hence are equal.
For every \(\ell\) we therefore obtain the completion of one
fixed integral unit lattice. On a curve, norm to a punctured
rational line isolating one boundary point at a time. Norm
multiplies residues by residue-field degree, not ramification
index. Equation (FF.9.5) makes each boundary coefficient integral.
In higher dimensions use a smooth curve transverse to the chosen
boundary divisor; the local finite-level expression \(uz^a\)
has the same exponent on that curve. The injective integer divisor
matrix of a unit basis, and an invertible minor, then put all unit
exponents in \(\mathbb Q\cap\widehat{\mathbb Z}=\mathbb Z\).
We obtain a unit \(h\), normalized by \(h(x_0)=b_0\), with
\(\xi=\kappa(h)\). Equality of the remaining constant classes
follows from Hochschild--Serre and evaluation at \(x_0\).

\emph{Actual point functions.}
Equation (FF.9.1) leaves \(\varphi(x)=c_xh(x)\), \(c_x\in k^\times\).
For a curve and nonconstant \(h\), choose closed points of
degree greater than \(\deg h\), using (FF.9.2). Their values
lie outside \(L\), so the point function is nonconstant.
On the inverse image of \(\mathbb G_m\setminus\{1\}\), apply
the same class argument to \(T-1\), obtaining
\(\varphi(x)-1=d_xh_1(x)\), \(d_x\in k^\times\).
Choose a closed point whose degree exceeds the common pole-degree
bound of \(c_xh-d_xh_1-1\). Its vanishing forces this function
to be identically zero. Hence \(h_1=Ah+B\), with
\(A,B\in k^\times\). At every point with \(h(x)\notin k\),
linear independence of \(1,h(x)\) over \(k\) gives
\(c_x=-A/B=c\). Pulling back \(T-b\), for each
\(b\in k^\times\), gives in the same way a unit proportional
to \(ch-b\). At a point with \(h(x)\in k\), take
\(b=ch(x)\); invertibility excludes \(\varphi(x)\ne b\).
Thus \(\varphi=ch\) everywhere, and normalization gives \(c=1\).

If \(h\) is constant, pass to a nontrivial finite extension
\(L'/L\) and choose \(\beta\in L'\setminus L\).
Every value of \(\varphi\) lies in \(hk\subset L\), so
\(T-\beta\) pulls back to a unit on the whole curve. Its
associated unit has all point values in \(L'\), and the same
unbounded-degree argument makes it constant. Normalizing at
\(x_0\) gives both
\(\varphi(x)/h\in k\) and
\((\varphi(x)-\beta)/(h-\beta)\in k\). Since \(\beta/h\notin k\),
these force \(\varphi(x)=h\). For higher-dimensional \(X\),
after a finite extension join \(x_0\) to the geometric point
in question by a smooth curve. Its normalized unit has the same
Kummer class; the ratio is a divisible unit, hence constant, and
equals 1 at \(x_0\). This proves the same value equality.

We have proved precisely the point-function condition for regular
connected affine sources with a rational point in
\citet[Theorem~4.18(5)]{FF_CHW2024}. That theorem gives
\(\mathcal R(L)\) for all finite-type schemes. All arguments
apply after every finite extension of \(L\), as required by the
intermediate reductions.
\end{proof}

\begin{claim}{FF.10}{Proved}\label{clm:FFnormal}
Let \(D/\mathbb Q\) be finitely generated and let
\(D\subset k\subsetneq\bar D\) be an algebraic Galois extension.
Then \(\mathcal R(L)\) holds for every finitely generated
extension \(L/k\). In particular the field
\(\mathbb Q^{\mathrm{tr}}\) of totally real algebraic numbers
reconstructs. The assertion does not claim closure under arbitrary
subfields of \(L\).
\end{claim}
\begin{proof}
First suppose \(\mu_\infty\subset k\). Choose a finite Galois
extension \(M/D\) not contained in \(k\). Then \(kM/k\) is
proper finite and \(kM/D\) remains Galois. The field \(D\) is
Hilbertian. Weissauer's theorem, in the opened stronger form
\citet[Corollary~2.3]{FF_JardenRazon2020}, applied with its
\(R=K=D\), \(M=N=k\) and \(M'=kM\), makes \(kM\)
Hilbertian. Claim~\ref{clm:FFcore} gives reconstruction over
\(LkM\), which is finitely generated over \(kM\). The extension
\(LkM/L\) is algebraic; algebraic descent in
Claim~\ref{clm:FFdescent} below gives \(\mathcal R(L)\).

For general \(k\), put \(D_{\mathrm{cyc}}=D(\mu_\infty)\)
and \(k'=kD_{\mathrm{cyc}}\). The cyclotomic extension is
abelian and proper in \(\bar D\): Hilbert specialization gives
a nonabelian symmetric extension of \(D\), which an abelian
extension cannot contain. Concretely \(X^3-X-T\) has regular
group \(S_3\) over \(D(T)\) by the Morse-polynomial theorem
\citep[Theorem~4.4.5]{FF_Serre2007}, and specialization over
the Hilbertian field \(D\) retains that group.
Both \(k\) and \(D_{\mathrm{cyc}}\)
are proper Galois extensions of \(D\). If one contains the
other, their compositum is plainly proper. Otherwise
\citet[Theorem~2.2]{FF_JardenRazon2020}, with \(R=D\), says
that their compositum is Hilbertian, so is again proper.
Thus \(k'/D\) is a proper Galois extension containing
\(\mu_\infty\). The preceding case applies to \(Lk'\);
algebraic descent gives the assertion for \(L\).

Finally \(\mathbb Q^{\mathrm{tr}}/\mathbb Q\) is Galois:
every conjugate of a totally real algebraic number is again
totally real. It is proper, since it is contained in \(\mathbb R\)
and omits \(i\). The theorem therefore applies directly.
This argument also proves that the Hilbertianity or smallness
assumption of AP-F.49 is a sufficient route, not a necessary
hypothesis for its normal-extension conclusion.
\end{proof}

\begin{claim}{FF.11}{Proved}\label{clm:FFdescent}
Let \(K/k\) be algebraic in characteristic zero. If
\(\mathcal R(K)\) holds, then \(\mathcal R(k)\) holds.
\end{claim}
\begin{proof}
It is enough to verify the regular-affine point-function condition
of CHW Theorem~4.18(5). Given the pinned relative map over \(k\),
slice at each finite \'etale constant object corresponding to a
finite subextension of \(K/k\). These are the actual small \'etale
topoi of the finite scalar extensions. Their compatible inverse-image
functors pass to \(K\): every quasicompact separated \'etale chart,
every map between finitely many charts and every finite covering
family descends to a finite stage. CHW Theorem~A.8 and Corollary~A.9
ensure that the inverse-image functor preserves these representable
charts, while Stacks Tags~0EYL and~01ZM give their finite-presentation
descent \citep{FF_StacksLimits}. Thus we obtain a relative map on
the small topos of \(X_K\), not an assumed arbitrary transcendental
base-change square.

Closed points of \(X_K\) map to closed points of \(X\), because
\(K/k\) is algebraic. If the latter point has residue field
finite over \(k\), its scalar-extension fibre is zero-dimensional
over \(K\); every resulting point is closed. The lifted map
is therefore pinned. Reconstruction over \(K\) gives a unit
\(h_K\) on the regular source, with its prescribed point function.
Use one common algebraic closure \(\bar k\) of \(K\) and \(k\).
That point function is the original \(G_k\)-equivariant function
on \(X(\bar k)\). For every \(\sigma\in G_k\), the two
regular functions \(\sigma h_K\) and \(h_K\) have equal values
at every geometric point. Reducedness makes them equal.
Their finitely many coefficients therefore descend to \(k\).
CHW's criterion supplies reconstruction for all schemes over \(k\).
This verifies AP-F.32 in the direction used above.
\end{proof}

\begin{claim}{FF.12}{Proved}\label{clm:FFaudit}
The dependent core of AP-F.24--F.49 used here is valid with its
stated field and morphism restrictions. Its positive conclusions
AP-F.37, F.43--F.46 and F.49 remain sufficient statements; its
negative conclusions AP-F.39--F.42 and F.47--F.48 remain obstructions.
None constitutes a necessary-and-sufficient classification.
\end{claim}
\begin{proof}
We record the checks at the places where a stronger, incorrect
replacement would break the argument.

In AP-F.24 the arithmetic wreath embedding is
\((G\times_P\cdots\times_PG)\rtimes S_n\), with one common
constant quotient \(P\). It is not the unrestricted wreath
product. The jet conditions on a high-degree pencil give simple
branch transpositions and isolated boundary branch labels.
A spanning-tree gauge puts a literal \(S_n\) in the geometric
monodromy. Remaining base labels have support at most two and
generate the geometric group in their coordinate projections.
For \(n\ge5\), \(A_n\) is contained in the stable derived
term. If \(u\) is one of those base elements and
\(\tau^{-1}(i)\) is outside its support, then
\([u,\tau]_i=u_i\). Thus the stable derived term projects onto
the whole geometric group in each coordinate. After Hilbert
specialization a prosolvable base change retains that term, and
the surviving constant quotient is still \(P\). Multiplying
by a literal even permutation returns the selected block without
changing its constant image. This proves arithmetic, rather than
merely geometric, detection. The same argument with the kernel
of the largest allowed formation quotient proves AP-F.44 and
AP-F.46. Properness of the hereditary formation excludes all
sufficiently large alternating groups, which is exactly what
this replacement needs.

AP-F.26--F.31 use the five steps exhibited in
Claim~\ref{clm:FFcore}: spanning of the coinvariant Albanese by
point cocycles, arithmetic equivariance after actual values descend,
integral Tate-Hom over a finitely generated field, integral boundary
residues, and equality of actual point functions. The norm residue
is weighted by residue degree. The principal-divisor lattice need
not be primitive; the integral Tate-Hom assertion, or equivalently
the torsion-free Picard Tate quotient, controls that denominator.
Pointwise Kummer equality leaves a divisible scalar and is not the
end of the proof. Pullbacks of \(T-b\) and closed points of
unbounded degree remove it. The final step is CHW Theorem~4.18(5),
not a claim that the fundamental group equals the full topos.

For AP-F.43--F.46, the two descent bases retain distinct roles:
the Hilbertian geometric base supplies detection, while a finitely
generated arithmetic field with normal \(G_L\) supplies
Tate-Hom equivariance. The norm field containing each actual value
is fixed before its Kummer level varies. AP-F.45 uses moved
geometric divisors when all algebraic constants are present;
AP-F.43 uses moved arithmetic primes when a number-field extension
is missing. Claim~\ref{clm:FFcore} verifies both alternatives
without imposing an unnecessary condition on which occurs.
AP-F.49 uses only the inclusion (FF.9.1). Its small-group argument
uses open subgroups and quotients, not arbitrary closed subgroups.
Small extensions preserve Hilbertianity by
\citet[Theorem~4.4]{FF_BarySoroker2008}. The bounded-degree
compositum proof retains the arithmetic constant field and the
degree of the detecting point.

The algebraic descent step is Claim~\ref{clm:FFdescent}.
The subfield assertions in AP-F.37 and F.43--F.46 additionally
take relative algebraic constants in a finite-type descent field
and their normal closures; normality and regular base change
give precisely the claimed pro-formation quotient. These arguments
must retain their stated positive-transcendence restriction when
\(\overline{\mathbb Q}\) is among the constants.

For the negative statements, AP-F.39--F.40 are the centralizer
test of Claim~\ref{clm:FF6}. In AP-F.41--F.42 the Puiseux
closure uses series whose coefficients lie in a single finite
constant extension. The finite trace-dual expansion proves the
algebraic-core identity; arbitrary infinite coefficient fields
are not substituted. In AP-F.47--F.48 the compatible roots
\((1+t_1)^{1/n}\) are in the henselization by Hensel uniqueness.
The substitution \(t_1\mapsto t_1+t_1^2\) centralizes inertia
and the coefficient action on an explicitly generated algebraic
closure. A coefficient field and rationally independent values
reduce the Abhyankar equality case to an algebraic extension of
that monomial henselization. Transcendence defect destroys the
last algebraicity implication, so that proof does not cover it.

The external inputs just used were opened with their numbered
statements: SGA~1 XII~5.1 and XIII~2.12,~4.6; Serre~3.3.1,
3.4.1; Bary-Soroker--Paran~1.2,~3.2; Poonen--Slavov~1.5;
Spie\ss--Szamuely~4.1--4.3; Jossen~1.4; CHW~3.6,~4.18,
A.8--A.9; and the valuation statements listed in the bibliography.
The original Jannsen ICM article could not be reopened through
the oversized proceedings file; the exact integral theorem is
instead verified in Jossen's opened primary article. We do not
mark that unavailable original as opened.
\end{proof}

\subsection{The pro-system, with its morphisms specified}

We use the following precise relative interpretation of the
pro-form in cote~149, scans~6--8. Let \(I\) be the filtered poset
of finitely generated subfields \(F\subset k\) containing a fixed
field of definition \(F_0\) for \(X,Y\). Put
\(X_F=X_0\times_{F_0}F\), and similarly for \(Y_F\).
The transition functors are arbitrary geometric morphisms; they
are not required to be pinned. A \emph{coherent pinned relative
model morphism} is a germ, under cofinal restriction, of relative
pinned morphisms
\[
       \varphi_F:(X_F)_{\mathrm{et}}\longrightarrow
                         (Y_F)_{\mathrm{et}}
                  \quad\text{over }BG_F,
\]
together with commutation isomorphisms for every \(F\subset F'\),
over \(BG_F\), satisfying the composition coherence. Isomorphic
germs are identified. This definition includes the relative base
identifications; a family of absolute morphisms is insufficient.

\begin{claim}{FF.13}{Proved}\label{clm:FFpro}
For any characteristic-zero field \(k\) and absolutely weakly
normal finite-type \(k\)-schemes \(X,Y\), their \(k\)-morphisms
are in natural bijection with coherent pinned relative model
morphisms as just defined. The assertion is independent of the
initial field and the chosen models, under cofinal refinement.
This proves a precise coherent relative interpretation of F1(iii).
It does not identify this category with an unspecified ordinary
pro-category equipped only with a condition on closed points.
\end{claim}
\begin{proof}
Finite-presentation descent gives models and descends each
\(k\)-morphism to some finitely generated field, by Stacks
Tag~01ZM \citep{FF_StacksLimits}. These models can be taken
absolutely weakly normal. In fact in characteristic zero this is
seminormality, and it descends under a faithfully flat field
extension: in an affine reduced noetherian ring, the criterion
\(b^2,b^3\in A\Rightarrow b\in A\) in its total quotient
ring descends by flatness and contraction of ideals. Apply this
to a model whose scalar extension is the given seminormal ring.
The same argument applies to every intermediate model. Thus the
algebraic morphism gives the specified coherent system.

Conversely, CHW Theorem~0.1 reconstructs each \(\varphi_F\) as
a unique morphism \(f_F:X_F\to Y_F\). We must prove that
compatibility of topoi really implies scalar compatibility,
because a transcendental transition need not be pinned.
Fix \(F\subset F'\). For a point \(x\in X_F(\bar F)\),
put \(E=F'(x)\), finite over \(F'\), in a common algebraic
closure. The two composites to \(Y_F\) have the same relative
point morphism over \(BG_F\), by coherence. Their image is the
closed point \(y=f_F(x)\). Since \(F(y)/F\) is finite
separable, relative morphisms \(BG_E\to BG_{F(y)}\) over
\(BG_F\) are exactly the \(F\)-embeddings \(F(y)\to E\).
This assertion needs no algebraicity of \(E/F\): the finite
\(G_F\)-set of embeddings is pulled back to \(G_E\), whose
fixed elements are precisely embeddings into \(E\).
Consequently \(f_{F'}(x)=(f_F)_{F'}(x)\) with their residue
embeddings, for every such \(x\).

For completeness this point argument also handles nonseparated
targets. Cover \(Y_F\) by affine opens \(V\). Coherence on
the corresponding subterminal objects says that the inverse
images of \(V_{F'}\) under the two morphisms are the same
open, namely \((f_F^{-1}V)_{F'}\). On this open the target
is affine. The \(\bar F\)-points of the model are dense after
any scalar extension: expand the finitely many coefficients of
a function in a linearly independent set over \(\bar F\).
Vanishing at all \(\bar F\)-points then makes every coefficient
function zero. Reducedness of the source therefore gives equality
of the two morphisms on that open. They glue, proving
\(f_{F'}=(f_F)_{F'}\).

Choose one index and restrict to the cofinal upper cone above it.
The entire family is now the scalar extension of its morphism
at that index; it gives a \(k\)-morphism. Uniqueness and
independence of models are again finite-presentation descent.
All comparisons retained the structural 2-cells, so this is a
relative statement.
\end{proof}

\begin{claim}{FF.14}{Proved}\label{clm:FFc}
With coherent relative scalar descent understood as in
Claim~\ref{clm:FFpro}, candidate F3(c) is correct:
\(\mathcal R(k)\) holds if and only if every pinned relative
map between finite-type \(k\)-scheme topoi is the realization
of such a system of finitely generated models, after absolute
weak normalization of its source. The single commuting-square
interpretation is excluded by Claim~\ref{clm:FF8}.
\end{claim}
\begin{proof}
A finite-type target may also be replaced by its absolute weak
normalization: its small \'etale topos is unchanged, and the
right-adjoint property gives
\(\operatorname{Hom}(X^{\mathrm{awn}},Y^{\mathrm{awn}})
 =\operatorname{Hom}(X^{\mathrm{awn}},Y)\)
\citep[Theorem~1.13]{FF_CHW2024}. Thus the source and target
hypotheses of Claim~\ref{clm:FFpro} lose no maps in this criterion.
A coherent system realizes on the small \'etale sites by the
finite-presentation colimit of charts and covering families
(Stacks Tag~0EYL). Claim~\ref{clm:FFpro} shows that it is the
system of one algebraic morphism, so its realization is precisely
that morphism's small-topos map. Conversely every algebraic
morphism has such a system. Their essential image is therefore
exactly the algebraic maps. If reconstruction holds, every pinned
relative map is in that image. If every pinned relative map has
such a coherent descent, it is algebraic by the same claim.
Faithfulness holds over every characteristic-zero field: equal
relative topos maps give equal geometric-point maps, including
residue embeddings, and equality on affine opens of a reduced
source gives equality of the scheme maps. This proves both
directions. This characterization is a descent criterion; it
does not supply an intrinsic classification in terms of familiar
field invariants.
\end{proof}

\begin{claim}{FF.15}{Proved}\label{clm:FFcomplex}
For \((\mathbb C,\mathbb C)\), the algebraic pro-profinite
group of cote~149, scan~8, is the system
\[
 \left\{
   G_F=\varprojlim_{N/F\ {
m finite\ Galois},\ N\subset\mathbb C}
             \operatorname{Gal}(N/F)
 \right\}_{F\subset\mathbb C, F/\mathbb Q\ {
m finitely\ generated}},
                                                               \tag{FF.15.1}
\]
with restriction maps defined using the algebraic closures
\(\bar F=\{z\in\mathbb C:z\text{ is algebraic over }F\}\).
Its inverse limit is trivial, but the system is not the trivial
pro-object in profinite groups. The relative model theorem
\ref{clm:FFpro} applies to every absolutely weakly normal
finite-type complex scheme over this system of base topoi.
\end{claim}
\begin{proof}
An abstract algebraic description dispenses even with a chosen
embedding of the individual number fields in \(\mathbb C\).
Choose a set \(B\) of cardinality \(2^{\aleph_0}\) and an
algebraic closure \(\Omega\) of \(\mathbb Q(B)\).
The classification of algebraically closed fields identifies
\(\Omega\) with \(\mathbb C\), noncanonically.
A cofinal system of index fields consists of the finite Galois
extensions of \(\mathbb Q(B_0)\) inside \(\Omega\), where
\(B_0\) runs through the finite subsets of \(B\).
Indeed finitely many elements of \(\Omega\) are algebraic over
one such rational function field, and a finite normal closure
contains all of them. This description uses only independent
indeterminates, polynomial algebra, finite extensions and
automorphism groups. Different choices give isomorphic pro-objects,
with no claimed canonical identification.

Every finitely generated subfield is specified algebraically by
a finite tuple \(a\in\mathbb C^n\):
\(F=\operatorname{Frac}(\mathbb Q[T_1,\ldots,T_n]/I_a)\),
where \(I_a\) is the ideal of polynomial relations of that tuple.
Finite Galois extensions inside \(\mathbb C\), their finite
automorphism groups, composita, and restrictions define both
levels in (FF.15.1); no analytic topology enters. For
\(F\subset F'\), \(\bar F\subset\bar F'\), and restriction
is a continuous group homomorphism \(G_{F'}\to G_F\).
Its image is the open subgroup corresponding to
\(F'\cap\bar F\), finite over \(F\) by Stacks
Lemma~9.26.11 \citep{FF_StacksFields}.

A compatible family of automorphisms at every \(F\) fixes each
element of \(\mathbb C\): place that element in the index field
\(\mathbb Q(a)\), where it must be fixed. Thus the inverse
limit is \(G_{\mathbb C}=1\). At the finite quotient level,
the same fact is visible because the index field may be enlarged
to contain any chosen finite Galois extension.
Nevertheless no transition to the stage \(G_{\mathbb Q}\)
is the trivial map. Its image is
\(G_{F'\cap\overline{\mathbb Q}}\), the nontrivial absolute
Galois group of a number field. The identity of a pro-object
can equal the zero morphism only if, for every stage, some later
transition to that stage is zero. This fails here. Keeping the
two levels of the pro-system, rather than replacing them by
their inverse limit or flattening the finite quotients, is
therefore essential. The final reconstruction statement is
Claim~\ref{clm:FFpro} with \(k=\mathbb C\).
\end{proof}


\subsection{Valuations: positive and negative cases, and the remaining test}

\begin{claim}{FF.16}{Proved}\label{clm:FFmixed}
Residue characteristic \(p\) does not, even together with any
specified finite transcendence degree over \(\mathbb Q\), decide
reconstruction. For each prime \(p\) and integer \(d\geq0\),
there exist henselian valued fields of transcendence degree \(d\)
over \(\mathbb Q\), with residue characteristic \(p\), which
reconstruct, and others which fail reconstruction.
\end{claim}
\begin{proof}
Choose algebraically independent elements
\(a_1,\ldots,a_d\in\mathbb Q_p\), and give
\(D=\mathbb Q(a_1,\ldots,a_d)\) the restricted \(p\)-adic
valuation. Such a tuple exists: \(\mathbb Q_p\) is uncountable,
whereas a field of countable transcendence degree over the
countable field \(\mathbb Q\), and its algebraic closure, are
countable. Henselization universality embeds \(D^h\) as a
valued subfield of the henselian field \(\mathbb Q_p\).
Henselization is algebraic and immediate
\citep[Proposition~3.3.25, Corollary~3.3.26]{FF_ADH2017}.
Thus \(\operatorname{trdeg}_{\mathbb Q}D^h=d\), its value
group is \(\mathbb Z\), and its residue field is
\(\mathbb F_p\): these assertions hold for \(D\) because
\(\mathbb Q\subset D\subset\mathbb Q_p\).
The field \(D^h\) is sub-\(p\)-adic, so Carlson--Stix
Theorem~A gives \(\mathcal R(D^h)\)
\citep[Theorem~A, Theorem~4.1]{FF_CS2026}.
For \(d=0\) this is already the ordinary henselization of
\(\mathbb Q\) at \(p\), embedded in \(\mathbb Q_p\).

For the opposite verdict take an algebraic closure
\(K=\overline{\mathbb Q(t_1,\ldots,t_d)}\).
Extend the \(p\)-adic Gauss valuation on the displayed rational
function field to \(K\). Its residue characteristic remains
\(p\), and an algebraically closed valued field is henselian.
The transcendence degree is \(d\). The group \(G_K\) is
trivial while \(\operatorname{Aut}(K)\) is nontrivial; a
nonidentity automorphism of \(\overline{\mathbb Q}\) extends
after choosing a transcendence basis. The centralizer obstruction
of Claim~\ref{clm:FF6} therefore gives \(\neg\mathcal R(K)\).
This supplies both verdicts without inferring a uniform answer
from the residue characteristic.
\end{proof}

\begin{claim}{FF.17}{Proved}\label{clm:FFdefect}
There is an explicit henselian field \(K/\mathbb Q\) of
transcendence degree two, residue characteristic zero and
transcendence defect one which fails reconstruction. Namely, in
\(\overline{\mathbb Q}((t))\) put
\[
 u=\sum_{n\geq1}t^{n!},\qquad
 K=\overline{\mathbb Q}(t,u)^h
\]
for the induced \(t\)-adic valuation and its henselization
inside that complete field. Then
\[
 \kappa(K)=\overline{\mathbb Q},\qquad
 v(K^\times)=\mathbb Z,\qquad G_K\simeq\widehat{\mathbb Z}.
\tag{FF.17.1}
\]
The transcendence degree two is the smallest possible finite
transcendence degree for a nontrivial residue-characteristic-zero
valuation with positive transcendence defect.
\end{claim}
\begin{proof}
The power series \(u\) is transcendental over
\(\overline{\mathbb Q}(t)\). Indeed let
\(P(t,U)\in\overline{\mathbb Q}[t,U]\) be nonzero, of
\(U\)-degree \(e\) and \(t\)-degree at most \(m\), and
write \(u_N=\sum_{n=1}^N t^{n!}\). If \(P(t,u)=0\), then
\[
 v\bigl(P(t,u_N)\bigr)\geq(N+1)!,\qquad
 \deg_tP(t,u_N)\leq m+eN!
\]
whenever the polynomial on the left is nonzero. For all
sufficiently large \(N\) these inequalities are incompatible.
Consequently infinitely many distinct \(u_N\) would be roots
of the nonzero polynomial \(P(t,U)\) over
\(\overline{\mathbb Q}(t)\), a contradiction.
Henselization is algebraic and immediate, proving the asserted
transcendence degree, residue field and value group.

For compatible roots of \(t\), let
\(E=\bigcup_n K(t^{1/n})\). Each extension has degree \(n\)
by the distinct valuation classes of
\(1,t^{1/n},\ldots,t^{(n-1)/n}\); since all roots of unity
are in \(K\), it is cyclic Galois. The union is henselian,
with residue field \(\overline{\mathbb Q}\) and divisible
value group \(\mathbb Q\). In residue characteristic zero,
these two conditions make a henselian field algebraically closed
\citep[Corollary~3.3.22]{FF_ADH2017}. Thus \(E=\bar K\)
and \(G_K=\varprojlim_n\mu_n\simeq\widehat{\mathbb Z}\).
Its nontrivial centre contradicts the necessary condition in
Claim~\ref{clm:FF6}. Notice that this obstruction may fix \(K\)
pointwise: it is a nontrivial central automorphism of its algebraic
closure, detected by the retained relative structural 2-cell.

Finally the rational rank of a nontrivial ordered value group
is at least one. The residue transcendence degree is nonnegative.
Positive transcendence defect therefore requires
\(\operatorname{trdeg}_{\mathbb Q}K\geq2\). In our example
\(2-0-1=1\), as claimed. The attached finite script checks
the displayed factorial-gap inequalities for sample bounds; the
argument above proves them for all bounds and proves transcendence.
\end{proof}

\begin{claim}{FF.18}{Open}\label{clm:FFremaining}
The following sufficient direction of candidate F3(a) remains
unproved and unrefuted in this work:
\[
 \bigl[C_{\operatorname{Aut}(\bar k)}(G_E)=1
          \text{ for every finite }E/k\bigr]
                  \quad\Longrightarrow\quad\mathcal R(k).
\tag{FF.18.1}
\]
Nor have we decided every finite-transcendence-degree henselian
field with residue characteristic zero and transcendence defect.
In particular our arguments do not decide
\[
 K_* = \mathbb Q\bigl(t,\sum_{n\geq1}t^{n!}\bigr)^h
           \subset\mathbb Q((t)),
\tag{FF.18.2}
\]
with its induced valuation. This is a specified unresolved test,
not a claim that it is the least field under inclusion or that
its centralizers satisfy (FF.18.1).
The general mixed-residue-characteristic case is also not
classified. A previously unresolved degree-zero test was
\[
                  K_p=\operatorname{Frac}(\mathbb Z_{(p)}^{sh}),
                                                        \tag{FF.18.3}
\]
with a chosen strict henselization inside
\(\overline{\mathbb Q}\). Its reconstruction is now
proved in Claim~\ref{clm:FFstrictreconstruct}, and
Claim~\ref{clm:FFdiscretealgebraic} treats every algebraic
extension of \(\mathbb Q\) admitting a mixed-characteristic
discrete valuation. These conclusions do not classify
arbitrary nondiscrete valuations or all finite-transcendence-degree
mixed-characteristic henselian fields.
Claims~\ref{clm:FFtamereconstruct}--\ref{clm:FFpadicsubfields}
further settle bounded wild ramification. The next
specified degree-zero test
\(U_5=T_5(\mu_{5^\infty})\) from
Claim~\ref{clm:FFcyclotomictest} is also settled
positively by Claim~\ref{clm:FFcyclotomicreconstruct},
despite the detector failure in
Claim~\ref{clm:FFcyclotomicdetector}.
Claims~\ref{clm:FFmaxabelianreconstruct} and
\ref{clm:FFnilpotentreconstruct} further settle the
maximal abelian and each fixed bounded-nilpotency
extension. A remaining specified degree-zero
test is the maximal metabelian field
\(B_2\) of Claim~\ref{clm:FFmetabeliantest},
which has the centralizer condition over every
finite extension. Its reconstruction is not decided.
Claims~\ref{clm:FFmetabelianaffine}--\ref{clm:FFhighergenusnorm}
subsequently determine affine automorphism point functions,
all abelian generic actions on rational and elliptic
Galois slices and on the specified genus-two slice,
and the portion of the etale site
generated by genus-zero covers. This includes all
abelian covers of arbitrary genus, all Galois genus-one
covers, and all degree-three genus-one covers.
Arbitrary covers and noninvertible pinned maps
remain unresolved.
\end{claim}
\begin{proof}[State of the attempted proof and its exact boundary]
Claim~\ref{clm:FF6} proves the necessary direction. The known
substitution ghosts of AP-F.41, F.42, F.47 and F.48 fail its
hypothesis, so none refutes the converse. In particular one must
not infer sufficiency from having checked all the previous
counterexamples. Attempts to replace the reconstruction argument
by this rigidity condition have not supplied the missing step:
centralizers control invertible semilinear ghosts, whereas a
general pinned map to \(\mathbb G_m\) has Kummer cocycles and
actual point values which have not been forced to descend by
centralizer triviality alone. The finite-quotient detection and
normal arithmetic action in Claim~\ref{clm:FFcore} are genuine
extra inputs; we have not deduced them from (FF.18.1).

For valuations, AP-F.48 supplies a coefficient field \(F\) and
parameters with rationally independent values. Under the
Abhyankar equality the ambient field is algebraic over the
corresponding monomial henselization, so algebraic descent
transfers its failure. Positive transcendence defect removes
that algebraicity statement. It does not remove all failures,
as Claim~\ref{clm:FFdefect} shows, but that example uses
algebraically closed residue field and a central Galois group.
For (FF.18.2) the residue field is \(\mathbb Q\), and the
same radical-and-constant argument gives
\(G_{K_*}\simeq\widehat{\mathbb Z}(1)\rtimes G_{\mathbb Q}\).
Its centre does not supply the previous obstruction.

There is a concrete obstruction to copying the old substitution
proof. Put \(u=\sum t^{n!}\). Then
\(v(u-t-t^2)=6\), whereas substitution
\(t\mapsto t+t^2\), \(u\mapsto u\), gives
\[
 u-(t+t^2)-(t+t^2)^2
       =-t^2-2t^3-t^4+t^6+\cdots,
\]
of valuation two. Thus that field embedding of
\(\mathbb Q(t,u)\) is not valuation preserving and does not
extend by the henselization universal property used in AP-F.48.
Allowing \(u\) to move would require an actual algebraic
self-correspondence preserving its formal branch; none has been
constructed here. Attempts to use rigid henselian fields would
also require proving the centralizer condition for every finite
extension and independently constructing a nonalgebraic pinned
map. Neither assertion follows from field rigidity alone.

Similarly, Claim~\ref{clm:FFpro} proves the pro-form for the
explicit coherent relative model category. The transition maps
of the models need not be pinned, so simply writing an ordinary
\(\operatorname{Pro}(\operatorname{Top}^{\mathrm{pin}})\)
is invalid. We have not proved that any other proposed definition
of a closed-point-preserving ordinary pro-morphism has the same
morphisms as that coherent category. The literal unspecified
pro-category reading of F1(iii), and that part of F6, is therefore
not silently marked proved. The correct coherent-descent version
of F3(c) is Claim~\ref{clm:FFc}; the bare commuting-square
version is refuted by Claim~\ref{clm:FF8}.
\end{proof}

\subsection{Function fields over unrestricted characteristic-zero constants}

\begin{claim}{FF.19}{Proved}\label{clm:FFfunction}
Let \(F\) be any characteristic-zero field. Every finitely
 generated transcendental extension \(L/F\) satisfies
\(\mathcal R(L)\). In particular this applies to
\(\mathbb C(t)\), without a finite-transcendence-degree
restriction on the constant field.
\end{claim}
\begin{proof}
First suppose that the constant field \(C\) is algebraically
closed, and that \(L/C\) is finitely generated of positive
transcendence degree. The field \(L\), and every finite extension
of it, is Hilbertian by
\citet[Theorem~1.2, Corollary~3.2]{FF_BarySorokerParan2013}.
The divisor argument gives
\[
                     D(E^\times)=C^\times
        \quad(E/L\text{ finite}),                         \tag{FF.19.1}
\]
since the relative constants in \(E/C\) are \(C\) itself.
Consequently arithmetic finite-quotient detection and unbounded
closed-point degrees hold exactly as in Claim~\ref{clm:FFcore}.
We verify separately the two steps of that proof which previously
used an arithmetic normal algebraic extension.

\emph{Integral puncture exponents.}
Use the punctured line, polynomials \(P_j\), actual value at
\(x_0\), and coefficients \(\lambda_j\) of (FF.9.3).
Choose a nontrivial finite Galois extension \(M/L\) containing
all roots of the \(P_j\), and
\(1\ne\sigma\in\operatorname{Gal}(M/L)\).
The finitely generated function field \(M/C\) has a proper
normal model, and the automorphism \(\sigma\) fixes its
algebraically closed constant field. The moved-prime-divisor
argument in Claim~\ref{clm:FFcore} supplies the disjoint good
valuation pairs. Weak approximation supplies \(\alpha\in M\)
with the valuation pattern (FF.9.4).
This time the actual value \(b_\alpha\) belongs to \(M\)
itself, because \(M/L\) is finite. At every Kummer level,
evaluate the class and apply \(\sigma\). Since the constant
normalization is fixed by \(\sigma\), we obtain
\[
 \frac{\sigma b_\alpha}{b_\alpha}
        \prod_j\left(\frac{P_j(\sigma\alpha)}
                            {P_j(\alpha)}\right)^{-a_{j,n}}
                         \in(M^\times)^n.
\]
Apply the discrete valuation \(v_i\) of \(M/C\), for every
\(n\). It gives
\(\lambda_i=-v_i(\sigma b_\alpha/b_\alpha)\in\mathbb Z\).
No norm to a finitely generated \(\mathbb Q\)-field and no
valuation on an infinite algebraic compositum are needed.

\emph{Arithmetic equivariance through a normal image.}
Let \(X/L\) be regular, affine and connected, with rational
base point \(x_0\), and let \(A\) be its generalized Albanese.
Choose a finitely generated subfield \(C_0/\mathbb Q\) of
\(C\) and a regular finitely generated function field
\(F_0/C_0\), with compatible models for this finite list of
geometric data, such that
\[
 L=F_0C,\qquad F_0\text{ and }C\text{ are linearly disjoint
 over }C_0.                                                \tag{FF.19.2}
\]
These are ordinary models of the geometrically integral
\(C\)-variety with function field \(L\); enlarge \(C_0\)
to include all coefficients. Geometric integrality is checked
after the algebraically closed extension \(C/C_0\), so the
model has the required regularity. Put
\(\bar C_0=\{c\in C:c\text{ is algebraic over }C_0\}\).
The algebraic closure of \(F_0\) inside \(L\) is
\(F_0\bar C_0\). To see this after adjoining \(\bar C_0\),
a nontrivial finite extension of the function field of a normal
geometrically integral variety cannot acquire a rational generic
section merely by extending its algebraically closed constant
field: a connected normal cover over that field remains
geometrically connected. This also proves linear disjointness
from every finite algebraic extension not already in the constant
extension. Thus restriction has image
\[
 \operatorname{im}(G_L\longrightarrow G_{F_0})
           =H:=G_{F_0\bar C_0}\triangleleft G_{F_0}.       \tag{FF.19.3}
\]
This is a quotient image, not an identification of \(G_L\)
with a subgroup of \(G_{F_0}\). Surjectivity onto the indicated
closed subgroup follows equivalently from the fixed field and
the Galois correspondence, or from the finite linear-disjointness
statement just proved.

Fix \(\ell\). The geometric Kummer functional
\(\lambda:T_\ell A\to\mathbb Z_\ell(1)\) of the pinned map
is invariant under \(G_L\); its action factors through (FF.19.3).
Since \(H\) fixes all roots of unity, \(\lambda\) factors
through \(W=(V_\ell A)_H\). Finite-quotient detection shows,
as in Claim~\ref{clm:FFcore}, that the projected point-Kummer
cocycles span \(W\). Choose finitely many spanning points and
put their finite fields of definition into a single finite
extension \(E/L\). Restriction to \(G_E\) keeps their spans,
since their projected cocycles are homomorphisms into a rational
\(\ell\)-adic vector space and this subgroup has finite index.

Their coordinates and actual multiplicative values are finitely
many elements of \(E\). Descend them by enlarging \(C_0\)
to a finitely generated \(C_1\subset C\) and taking a
regular finitely generated model field \(M_1/C_1\) inside
\(E\), containing \(F_0\), such that
\(E=M_1C\). In particular \(M_1/\mathbb Q\) is finitely
 generated. The same argument gives
\[
 \operatorname{im}(G_E\to G_{M_1})
       =H_1:=G_{M_1\bar C_1}\triangleleft G_{M_1}.         \tag{FF.19.4}
\]
Under restriction \(r:G_{M_1}\to G_{F_0}\), the image
\(r(H_1)\) is the open subgroup
\(\operatorname{im}(G_E\to H)\) of \(H\).
In particular \(H_1\) acts trivially on \(W\).

Each chosen Albanese point and its multiplicative value is
\(M_1\)-rational. Its Kummer cocycle is therefore defined on
\(G_{M_1}\). Equality of the two evaluated cocycles on
\(G_E\) descends to \(H_1\), since \(G_E\to H_1\) is
surjective and both cocycles are restrictions of cocycles over
\(M_1\). For \(g\in G_{M_1}\), \(h\in H_1\), the
cocycle conjugation identity, projected to \(W\), now gives
\[
 \lambda\bigl(r(g)c_x(h)\bigr)
             =\chi_\ell(g)\lambda\bigl(c_x(h)\bigr).
\]
The error term is of the form
\((1-r(ghg^{-1}))c_x(g)\), which vanishes in \(W\)
because \(H_1\) is normal and its image lies in \(H\).
The finitely many images span \(W\), proving
\(G_{M_1}\)-equivariance of \(\lambda\).
The integral Tate-Hom statement of
\citet[Section~1.4]{FF_Jossen2014}, now over the genuinely
finitely generated field \(M_1/\mathbb Q\), puts this
functional in the completed normalized-unit lattice.

The rest of Claim~\ref{clm:FFcore} transfers without a new
field hypothesis: norms to punctured lines and transverse curves
make boundary coefficients integral; the integer divisor matrix
makes unit exponents integral; base-point evaluation removes the
constant cohomology class. Equation (FF.19.1), high-degree closed
points and pullbacks of \(T-b\) remove the divisible pointwise
scalars. In the constant case choose a proper finite extension of
\(L\), which Hilbertianity supplies. This proves the precise
point-function condition CHW~4.18(5), hence \(\mathcal R(L)\).

Finally let \(L/F\) be as in the statement, without assuming
\(F\) algebraically closed. In an algebraic closure of \(L\)
choose an algebraic closure \(\bar F\) of \(F\). The field
\(L\bar F\) is finitely generated of the same positive
transcendence degree over \(\bar F\), so reconstructs by the
case just proved. Since \(L\bar F/L\) is algebraic,
Claim~\ref{clm:FFdescent} gives \(\mathcal R(L)\).
\end{proof}

\begin{claim}{FF.20}{Known}\label{clm:FFmixedcentralizer}
The necessary centralizer condition of candidate F3(a) holds for
every characteristic-zero field admitting a discrete valuation
of residue characteristic \(p>0\). It also holds for every
field finitely generated and transcendental over any subfield.
These statements do not by themselves prove reconstruction.
\end{claim}
\begin{proof}
Hoshi--Tsujimura Corollary~4.1 says that
\(\operatorname{Aut}(E)\to\operatorname{Out}(G_E)\) is
injective if \(E\) embeds in the fraction field of a noetherian
local domain of mixed characteristic. Corollary~4.5(3) makes
\(G_E\) slim under the same hypothesis
\citep{FF_HoshiTsujimura2023}. A mixed-characteristic discrete
valuation ring is such a domain. Every finite extension of its
fraction field has a prolongation with value group contained
in \(\frac1e\mathbb Z\), hence again has a discrete
mixed-characteristic valuation. Thus both conclusions hold
for every finite extension \(E/k\).

An automorphism of \(\bar k\) centralizing \(G_E\) preserves
its fixed field \(E\); its restriction is in the kernel of the
displayed outer-action map. Injectivity makes that restriction
the identity, leaving an element of \(Z(G_E)=1\). This proves
the desired centralizer statement. For a finitely generated
transcendental field use instead Hoshi--Tsujimura Corollary~1.6
and Corollary~4.5(1); that hypothesis is also preserved by finite
extensions. It gives an independent necessary-condition check
on Claim~\ref{clm:FFfunction}.
\end{proof}

\subsection{The sufficient classes in R9, with their quantifiers}

\begin{claim}{FF.21}{Proved}\label{clm:FFclasses}
The following are sufficient conditions for \(\mathcal R(k_0)\),
always in characteristic zero.
\begin{enumerate}
\item There is a finitely generated \(D/\mathbb Q\) with
\(D\subset k_0\), \(k_0/D\) algebraic, and proper normal
closure \(N\ne\bar D\) inside \(\bar D\).
\item For some finitely generated \(D/\mathbb Q\), the field
\(k_0\) is a subfield of a finitely generated extension of
\(D^{\mathcal C}\). Here \(\mathcal C\) is a proper
isomorphism-invariant class of finite groups, containing all
cyclic groups and closed under subgroups, quotients and
extensions; \(D^{\mathcal C}\) is the maximal Galois
pro-\(\mathcal C\) extension. Equivalently, for some finitely
generated \(D_0\subset k_0\), the extension \(k_0/D_0\)
is algebraic and its normal-closure group is virtually
pro-\(\mathcal C\). This includes AP-F.37 and AP-F.43
by taking \(\mathcal C\) to be the solvable groups, and
AP-F.44 in general.
\item The field \(k_0\) has positive transcendence degree
over \(\mathbb Q\) and is a subfield of a finitely generated
extension of \(K_0^{\mathcal C}\), where
\(K_0/\overline{\mathbb Q}\) is finitely generated of
positive transcendence degree. Equivalently, it is algebraic
over some finitely generated \(D_0/\mathbb Q\) of positive
transcendence degree whose normal closure \(N/D_0\) has
\[
 \operatorname{Gal}
       (N\overline{\mathbb Q}/D_0\overline{\mathbb Q})
       \quad\hbox{virtually pro-}\mathcal C.
\]
This is AP-F.46. In particular every positive-transcendence
subfield of a finitely generated extension of
\(\overline{\mathbb Q}\) reconstructs, as in AP-F.45.
\item In the setting of AP-F.49, let \(D/\mathbb Q\)
be finitely generated of positive transcendence degree, and
let \(k/D\) be algebraic Galois with
\(\overline{\mathbb Q}\subset k\). If \(k\) is
Hilbertian, every finitely generated \(L/k\) reconstructs.
If instead
\(\operatorname{Gal}(k/D\overline{\mathbb Q})\) is
small, then also every positive-transcendence subfield
\(k_0\subset L\) reconstructs. Small means finitely many
open subgroups of each prescribed finite index.
\end{enumerate}
In (2) and (3), taking the union over \(\mathcal C\) is
equivalent to requiring that an open subgroup of the respective
normal-closure group have no finite quotient containing some
fixed \(A_m\), \(m\ge5\), as a subquotient.
These are sufficient classes and equivalences between their
descriptions, not a necessary classification of all reconstruction
fields. In particular (1) does not assert that arbitrary subfields
of a reconstruction field reconstruct.
\end{claim}
\begin{proof}
For (1), Claim~\ref{clm:FFnormal} gives \(\mathcal R(N)\),
and Claim~\ref{clm:FFdescent} descends along the algebraic
extension \(N/k_0\).

We give the subfield argument for (2), since subfield passage
cannot be inferred just from (1). Put \(k=D^{\mathcal C}\),
and choose a finitely generated \(F/\mathbb Q\) inside
\(L\), containing \(D\), with \(L=Fk\). Given
\(k_0\subset L\), choose a finite transcendence basis \(T\)
of \(k_0/\mathbb Q\), enlarge \(F\) to contain it, and
put \(A=\mathbb Q(T)\). The relative algebraic closure
\(R\) of \(A\) in \(F\) is finite over \(A\)
\citep[Lemma~9.26.11]{FF_StacksFields}, and \(F/R\) is
regular. Let \(N/R\) be the normal closure of \(Rk_0/R\).
Regularity identifies its conjugates over \(R\) with those
after scalar extension to \(F\). Since \(L/F\) is Galois
pro-\(\mathcal C\), all these conjugates lie in \(L\),
and restriction gives
\[
 \operatorname{Gal}(N/R)
   \simeq\operatorname{Gal}(FN/F),
\]
a quotient of \(\operatorname{Gal}(L/F)\). Thus it is
pro-\(\mathcal C\). This normal extension is proper.
Indeed a proper hereditary extension-closed class containing
the cyclic groups omits \(A_n\) for all sufficiently large
\(n\): if it contained arbitrarily large alternating groups,
the embeddings \(G\hookrightarrow S_{|G|}\hookrightarrow
A_{|G|+2}\) would put every finite group in the class.
Every finitely generated characteristic-zero field is
Hilbertian, and the Morse-polynomial construction supplies
\(S_n\)-extensions for every \(n\)
\citep[Theorems~3.4.1, 4.4.5]{FF_Serre2007}; for positive
transcendence degree use also
\citep[Theorem~1.2]{FF_BarySorokerParan2013}.
Consequently its absolute Galois group cannot be pro-
\(\mathcal C\). Apply (1) over \(R\), and then descend
from \(N\) to \(k_0\), since \(R/A\) and \(k_0/A\)
are algebraic.

For the stated equivalence, take the finite Galois closure
\(R'/A\) of \(R/A\). The normal closure of \(k_0/A\)
is contained in the compositum of the finitely many
\(A\)-conjugates of \(NR'\). Over \(R'\) this is
pro-\(\mathcal C\): its group is a closed subgroup of a
finite product of pro-\(\mathcal C\) groups. Thus its
group over \(A\) is virtually pro-\(\mathcal C\).
Conversely an open pro-\(\mathcal C\) subgroup of a normal
closure fixes a finite extension \(D_1/D_0\); the whole
normal closure is contained in \(D_1^{\mathcal C}\).
This supplies exactly the embedding condition in (2).

For (3), choose \(D/\mathbb Q\) finitely generated with
\(K_0=D\overline{\mathbb Q}\). The field
\(k=K_0^{\mathcal C}\) is normal over \(D\): every
\(D\)-automorphism preserves \(K_0\) and permutes its
Galois \(\mathcal C\)-extensions. We can choose \(F\)
as above with \(L=Fk\) normal over \(F\), while
\(L/F\overline{\mathbb Q}\) is pro-\(\mathcal C\).
The same \(A,R,N\) argument now gives
\[
 \operatorname{Gal}
   (N\overline{\mathbb Q}/R\overline{\mathbb Q})
 \quad\hbox{pro-}\mathcal C.
\]
Here regularity survives algebraic extension of constants,
so \(F\overline{\mathbb Q}/R\overline{\mathbb Q}\)
is regular and the indicated group is the corresponding
quotient of \(\operatorname{Gal}
(L/F\overline{\mathbb Q})\). The field
\(R\overline{\mathbb Q}\) has positive transcendence
degree over \(\overline{\mathbb Q}\); it is Hilbertian
and has the same arbitrarily large symmetric-group quotients.
Therefore \(N\overline{\mathbb Q}\ne\bar R\).
This is a proper normal extension of the finitely generated
field \(R/\mathbb Q\); (1) and algebraic descent apply.
The equivalence follows by the same finite normal-closure
argument, now after adjoining \(\overline{\mathbb Q}\).
Conversely, a finite extension of
\(D_0\overline{\mathbb Q}\) fixed by the specified open
subgroup is a finitely generated positive-transcendence
\(\overline{\mathbb Q}\)-field and supplies the required
geometric pro-\(\mathcal C\) envelope. The special case
with no geometric algebraic enlargement gives AP-F.45;
it also follows directly from Claim~\ref{clm:FFfunction}
and algebraic descent.

For (4), its Hilbertian assertion is exactly
Claim~\ref{clm:FFcore}. In the small case the original
Hilbertian proof uses
\citet[Theorem~4.4]{FF_BarySoroker2008} over the Hilbertian
field \(D\overline{\mathbb Q}\). We can also derive
the subfield conclusion directly from (1).
In the notation just used, \(k\cap F\overline{\mathbb Q}\)
is finite over \(D\overline{\mathbb Q}\), because it
is algebraic inside a finitely generated extension.
Thus \(\operatorname{Gal}(L/F\overline{\mathbb Q})\)
is an open subgroup of the given small group, and is small.
Its quotient
\(\operatorname{Gal}(N\overline{\mathbb Q}/
R\overline{\mathbb Q})\) is small as well. It cannot
be the entire absolute Galois group of
\(R\overline{\mathbb Q}\). To check this without a
smallness assertion for arbitrary closed subgroups, choose
a separating transcendence basis containing \(x\).
Outside finitely many \(a\in\overline{\mathbb Q}\),
the extension of the corresponding rational-function field
is unramified at \(x-a\), so a prime over this divisor
has \(\operatorname{ord}(x-a)=1\). The distinct
\(x-a\) therefore give infinitely many distinct square
classes, hence infinitely many quadratic extensions.
This absolute Galois group is not small. Once again
\(N\overline{\mathbb Q}/R\) is proper normal, and (1)
and algebraic descent prove the assertion.

Finally, groups with no \(A_m\) subquotient form a proper
class closed under subgroups, quotients and extensions and
containing all cyclic groups. For extension closure, a simple
nonabelian subquotient in an extension must be a subquotient
of either its normal kernel or its quotient: intersect the
chosen subgroup with the kernel and use simplicity.
Conversely the alternating-group exclusion for any proper
\(\mathcal C\) was proved above. This proves the final
equivalence and retains all the finite-extension and
positive-transcendence restrictions.
\end{proof}

\subsection{The negative classes and the two required ghost tests}

\begin{claim}{FF.22}{Proved}\label{clm:FFghosts}
The negative statements AP-F.39--F.42 and AP-F.47--F.48
have the following precise content.
\begin{enumerate}
\item Every characteristic-zero field with finite absolute
Galois group fails \(\mathcal R\); in particular all
algebraically closed and real closed fields fail.
\item Let \(F\) have characteristic zero and equip
\(F(t_1,\ldots,t_r)\), \(r\ge1\), with a monomial
valuation whose positive parameter values are rationally
independent. Its henselization \(H\), and every algebraic
extension of \(H\), fails \(\mathcal R\). This includes
every finite-transcendence-degree field over \(\mathbb Q\)
with a nontrivial henselian valuation of residue
characteristic zero satisfying the Abhyankar equality.
\item For any characteristic-zero \(F\), put
\[
 P(F)=\bigcup_{n\ge1}F((t^{1/n})),\qquad
 A(F)=\{a\in P(F):a\text{ is algebraic over }F(t)\},
\]
with compatible roots throughout. The fields \(F((t))\),
\(P(F)\), \(A(F)\), and all their algebraic extensions
fail \(\mathcal R\). There are topological isomorphisms
\[
 (G_{P(F)},\chi_{P(F)})\simeq(G_F,\chi_F),\qquad
 (G_{A(F)},\chi_{A(F)})\simeq(G_F,\chi_F),
\]
with the full cyclotomic characters, and
\(\lvert A(F)\rvert=\max(\lvert F\rvert,\aleph_0)\).
\end{enumerate}
The explicit obstruction in (2) and (3) is the lift of
\(t_1\mapsto t_1+t_1^2\), respectively
\(t\mapsto t+t^2\), centralizing the full absolute Galois
group. It gives an actual pinned relative automorphism of
the full \(\mathbb G_m\) topos. These ghosts also violate
the genus-zero Isom condition in candidate F3(b), already
on the four-punctured lines with parameters \(t\) and
\(t+t^2\); and they cannot satisfy coherent scalar descent
in candidate F3(c). Thus all three candidates correctly
exclude these particular ghosts, although candidate (b)
is false by its opposite-direction failure in
Claim~\ref{clm:FF7}.
\end{claim}
\begin{proof}
An algebraically closed characteristic-zero field has a
nonidentity automorphism. Choose a nontrivial automorphism
of \(\overline{\mathbb Q}\), fix a transcendence basis,
and extend to the algebraic closure of the resulting
rational-function field. The resulting embedding is onto:
its image is algebraically closed and the whole field is
algebraic over it. Its Galois centralizer condition fails,
so Claim~\ref{clm:FF6} applies. Reconstruction ascends
finite extensions by slicing at their finite \(\acute etale\)
constant objects, as in that claim. Hence a finite absolute
Galois group would force the finite algebraic closure to
reconstruct, a contradiction. In the real closed case one
can also use its nontrivial central element in the group
of order two; triviality of the field's own automorphism
group does not remove that relative structural-cell ghost.

For (2), write \(\Gamma_0=\sum_i\mathbb Z\gamma_i\)
for the monomial value group. Distinct exponent vectors
have distinct values, so least monomials prove both
multiplicativity and residue field \(F\). Henselization
is immediate \citep[Corollary~3.3.26]{FF_ADH2017}, giving
\(v(H^\times)=\Gamma_0\), \(\kappa(H)=F\).
Inside an algebraic closure choose \(\bar F\) and
compatible roots \(u_{i,n}=t_i^{1/n}\). For every finite
\(F'/F\) in \(\bar F\), residue independence of an
\(F\)-basis of \(F'\) proves
\[
 [HF':H]=[F':F],\quad
 v((HF')^\times)=\Gamma_0,\quad \kappa(HF')=F'.
\]
Indeed, dividing a nonzero linear combination by one of
its least-valued coefficients gives a nonzero residue.
The \(n^r\) monomials
\(u_{1,n}^{a_1}\cdots u_{r,n}^{a_r}\),
\(0\le a_i<n\), have distinct value classes modulo
\(\Gamma_0\). They are independent over \(HF'\) and
form a basis of
\(E(F',n)=HF'(u_{1,n},\ldots,u_{r,n})\), whose degree,
value group and residue field are therefore
\[
 [E(F',n):H]=[F':F]n^r,\quad
 v(E(F',n)^\times)=\tfrac1n\Gamma_0,\quad
 \kappa(E(F',n))=F'.
\]
The directed union \(\Omega=\bigcup_{F',n}E(F',n)\)
is henselian with divisible value group and algebraically
closed residue field. Since the residue characteristic
is zero, \citet[Corollary~3.3.22]{FF_ADH2017} makes
\(\Omega\) algebraically closed. This constructs the
entire algebraic closure and gives
\[
 G_H\simeq\widehat{\mathbb Z}(1)^r\rtimes G_F.
                                                        \tag{FF.22.1}
\]
The section acts on algebraic constants and fixes all
chosen parameter roots; the kernel multiplies the latter
by compatible roots of unity. Thus the action is the
diagonal cyclotomic action.

The substitution \(t_1\mapsto t_1+t_1^2\), fixing the
other variables and \(F\), preserves the monomial
valuation: each monomial acquires a unit of residue one
and the least value remains unique. It extends uniquely
to an embedding \(\sigma:H\to H\). Hensel's lemma
gives a unique positive-valued root \(r_1\) of
\(r_1+r_1^2=t_1\), with \(v(r_1)=\gamma_1\) and
\(r_1/t_1\) of residue one. Substitution by \(r_1\)
is likewise valuation preserving. The two induced
embeddings of \(H\) are inverse, first on the rational
function field and then by the uniqueness in henselization
universality. Thus \(\sigma\) is an automorphism.

For each \(n\), let \(\beta_n\in1+\mathfrak m_H\)
be the unique root of \(X^n=1+t_1\). These exist because
\(n\) is a unit in the residue field, and uniqueness
gives \(\beta_{mn}^{m}=\beta_n\). On \(E(F',n)\)
define
\[
 \tau|_H=\sigma,\quad \tau|_{F'}=1,\quad
 \tau(u_{1,n})=u_{1,n}\beta_n,\quad
 \tau(u_{i,n})=u_{i,n}\ (i>1).
\]
The displayed basis proves this is a well-defined
automorphism: the images satisfy the transformed equations,
and the image contains the base and all original generators.
Compatibility gives \(\tau\in\operatorname{Aut}(\Omega)\).
Every \(g\in G_H\) preserves \(\bar F\), multiplies
\(u_{i,n}\) by a root of unity, and fixes \(\beta_n\).
Consequently \(g\tau=\tau g\) on the field generators
of every \(E(F',n)\), hence on \(\Omega\). This is a
full-group centralizer calculation, not a claim about
the tame quotient alone.

Claim~\ref{clm:FF6} supplies the pinned relative
\(\mathbb G_m\)-topos automorphism. Its rational-point
function is \(a\mapsto\sigma^{-1}(a)\). A
\(H\)-morphism of \(\mathbb G_m\) is \(cT^m\),
and a topos automorphism forces \(m=\pm1\). The values
at \(1\) and \(2\) force \(c=1,m=1\), while the
value at \(t_1\) is \(r_1\ne t_1\). Thus the map
is not algebraic. Claim~\ref{clm:FFdescent}, used
contrapositively, excludes every algebraic extension.
Equivalently, the same \(\tau\) normalizes the fixed field
of every closed subgroup of \(G_H\), because it centralizes
that subgroup. It therefore restricts to every intermediate
algebraic extension and still moves \(t_1\).

For the intrinsic Abhyankar assertion, a henselian
equicharacteristic-zero valued field has a coefficient
subfield \(F\subset\mathcal O_v\) mapping isomorphically
to its residue field \citep[Proposition~3.3.8]{FF_ADH2017}.
Concretely, lift a residue transcendence basis and enlarge
the resulting coefficient field maximally; any missing
algebraic residue element has a separable minimal polynomial,
whose simple root lifts by Hensel's lemma, contradicting
maximality. Choose elements \(t_1,\ldots,t_r\) whose
positive values form a rational basis of the value group.
Their values imply algebraic independence over \(F\)
and the monomial valuation above. The henselization
embeds in the given henselian field. The Abhyankar equality
is exactly what makes the given field algebraic over
\(F(t_1,\ldots,t_r)^h\), so (2) follows. With
transcendence defect this last algebraicity conclusion
is unavailable, as Claim~\ref{clm:FFremaining} records.

For (3) the exact bounded-coefficient Puiseux closure is
\[
 \Omega=\bigcup_{\substack{F'/F\ \mathrm{finite}\\n\ge1}}
                       F'((t^{1/n})),\qquad F'\subset\bar F.
                                                        \tag{FF.22.2}
\]
It is an algebraic closure of \(F((t))\)
\citep[Proposition, p.~314]{FF_vanDenDriesRibenboim1984}.
The finite coefficient restriction matters. A series with
arbitrary coefficients in \(\bar F\) need not be
algebraic over \(F((t))\). Expanding the coefficients
in a finite \(F\)-basis of \(F'\) proves
\(\Omega=P(F)\bar F\). Also \(P(F)\cap\bar F=F\):
an element algebraic over \(F\) cannot have nonzero
valuation, by comparing the unique least term in its
minimal polynomial; reduction then supplies a root in
\(F\) of that irreducible polynomial. Coefficientwise
action therefore identifies \(G_{P(F)}\) with \(G_F\).
It and its inverse are continuous, since every element
of \(\Omega\) uses one finite coefficient extension.
Roots of unity are constants, proving cyclotomic compatibility.

Let \(\Omega_0\) be the elements of \(\Omega\)
algebraic over \(F(t)\). It is algebraically closed and
is an algebraic closure of \(A(F)\). We verify
\[
                         \Omega_0=A(F)\bar F.             \tag{FF.22.3}
\]
Take \(a\in\Omega_0\cap F'((v))\),
\(v=t^{1/n}\), enlarging \(F'/F\) to be finite
Galois. If \((e_i)\) and \((e_i^*)\) are trace-dual
\(F\)-bases of \(F'\), write
\[
 a=\sum_i e_i a_i,\qquad
 a_i=\operatorname{Tr}_{F'((v))/F((v))}(e_i^*a).
\]
Every coefficientwise conjugate of \(e_i^*a\) is
algebraic over \(F(t)\), so \(a_i\in A(F)\).
This proves (FF.22.3); its converse inclusion is immediate.
Restriction now gives the asserted topological and
cyclotomic identification \(G_{A(F)}\simeq G_F\).
For a finite \(E/F\), the same coefficient expansion
gives \(P(F)E=P(E)\), and the same trace argument gives
\(A(F)E=A(E)\). Hence all the corresponding open-subgroup
and finite-extension data agree as well.

Substitution \(t\mapsto t+t^2\) acts on every finite
Puiseux level by
\[
 t^{1/n}\longmapsto t^{1/n}(1+t)^{1/n},
\]
using the rational binomial series of constant term one.
Its inverse uses
\(r(t)=(-1+\sqrt{1+4t})/2\) and the corresponding
constant-term-one root of \(r(t)/t\). These substitutions
preserve (FF.22.2), are compatible as the denominator and
the coefficient field increase, and fix \(\bar F\).
They centralize coefficientwise \(G_F\). They also
preserve \(A(F)\) and \(\Omega_0\): one direction is
algebraic over \(F(t+t^2)\subset F(t)\), the other over
\(F(r(t))\), which is algebraic over \(F(t)\).
For \(F((t))\), an element of its Galois group acts on
the generators of (FF.22.2) by coefficient action and
root multipliers. The same generator calculation as in
(2) proves centralization there too. Thus the displayed
substitutions give actual nonalgebraic pinned relative
maps on all three full topoi, and algebraic descent
excludes all their algebraic extensions. Finally \(A(F)\)
is algebraic over and contains \(F(t)\), which proves
its cardinality assertion. Taking \(F=\mathbb Q\)
compares the reconstructing field \(\mathbb Q\) with
the nonreconstructing countable field \(A(\mathbb Q)\)
having the same full cyclotomic Galois data.

To test the candidates on these ghosts, put
\(U_t=\mathbb P^1\setminus\{0,1,\infty,t\}\)
and \(U_{t+t^2}=\mathbb P^1\setminus
\{0,1,\infty,t+t^2\}\). The centralizing lift makes
the semilinear isomorphism a relative fundamental-group
isomorphism over the identity of \(G_k\), with only
geometric inner ambiguity. A scheme isomorphism would
require \(t+t^2\) in
\[
 \left\{t,t^{-1},1-t,(1-t)^{-1},
                    (t-1)/t,t/(t-1)\right\}.
\]
Every such equality is a nonzero polynomial relation on
the transcendental element \(t\); none holds. Thus
candidate (b) excludes these fields even at four punctures.
Candidate (a) fails by the centralizing lift itself.
Candidate (c), with actual coherent scalar descent,
would contradict Claim~\ref{clm:FFc}, since its descended
map reconstructs over the finitely generated field.
The program \texttt{F\_puiseux\_checks.py} verifies both
inverse substitutions modulo \(t^{13}\), the compatible
roots for denominators dividing 12, and the six exact
cross-ratio inequalities. The infinite identities and
the topos maps have been proved above.
\end{proof}

\subsection{A countable mixed-characteristic test of transcendence degree zero}

\begin{claim}{FF.23}{Proved}\label{clm:FFstricttest}
Fix a prime \(p\) and a prolongation of its valuation to
\(\overline{\mathbb Q}\), and let
\(K_p=\operatorname{Frac}(\mathbb Z_{(p)}^{sh})\) as
in (FF.18.3). It is countable, algebraic over \(\mathbb Q\),
henselian discretely valued, with value group \(\mathbb Z\)
and residue field \(\overline{\mathbb F}_p\).
It satisfies candidate F3(a)'s centralizer condition over
every finite extension and candidate F3(b)'s Isom condition
for all hyperbolic curves over every finite extension.

For every number field \(D\subset K_p\), its normal
closure over \(D\) is \(\overline{\mathbb Q}\).
It is not sub-\(q\)-adic for any prime \(q\), and is
not a finite extension of a proper normal algebraic extension
of a number field. Consequently it lies in none of the
sufficient field classes stated in Claims~\ref{clm:FF4},
\ref{clm:FFnormal}, \ref{clm:FFfunction}, and
\ref{clm:FFclasses}. Its centralizer condition excludes
all the negative ghost classes in
Claim~\ref{clm:FFghosts}.
Thus it is a specified field on neither earlier list with the
smallest possible cardinality and transcendence degree
in characteristic zero. We claim no least field under
inclusion. Claim~\ref{clm:FFstrictreconstruct} subsequently
proves its reconstruction by a different argument.
\end{claim}
\begin{proof}
The strict henselization is a strictly henselian discrete
valuation ring with maximal ideal generated by \(p\)
and residue field \(\overline{\mathbb F}_p\)
\citep[Lemmas~10.155.2, 15.46.11]{FF_StacksHensel}.
It is a filtered colimit of pointed \(\acute etale\)
\(\mathbb Z_{(p)}\)-algebras. Their generic fields are
finite separable extensions of \(\mathbb Q\), so its
fraction field is algebraic over \(\mathbb Q\), and
therefore countable. The generator \(p\) gives the
specified normalization of its value group. The roots of
unity of every order prime to \(p\) lie in \(K_p\),
because their residue roots lift uniquely from finite
extensions of \(\mathbb F_p\). Their union is a
nontrivial divisible subgroup of \(K_p^\times\).
Sub-\(q\)-adic fields are Kummer-faithful
\citep[Remark~3.9(2)]{FF_CS2026}; hence \(K_p\) is not
sub-\(q\)-adic for any \(q\).

Claim~\ref{clm:FFmixedcentralizer} applies to this mixed
characteristic discrete valuation and to every finite
extension, proving candidate (a)'s condition.
The field embeds in
\(\operatorname{Frac}W(\overline{\mathbb F}_p)\),
the completion of the maximal unramified extension of
\(\mathbb Q_p\): the strict henselization universal
property gives a local map into the strictly henselian
Witt ring, and its kernel is zero because \(p\) remains
nonzero. Thus it is generalized sub-\(p\)-adic
in \citet[Definition~4.11]{FF_Mochizuki2003}.
Every finite extension remains a subfield of a finite
extension of that complete field, hence is still generalized
sub-\(p\)-adic. Theorem~4.12 of that source, with all
geometric primes retained, proves the relative hyperbolic
Isom assertion, modulo geometric inner conjugation.

We prove the two-valuation fact needed for the normal
closure, rather than use it without its comparison.
A field with two distinct rank-one henselian valuations
is separably closed. Indeed the valuation rings are
incomparable. There is an \(x\) with
\(v_1(x)>0>v_2(x)\): choose an element integral for the
first valuation and negative for the second, and if its
first value is zero multiply a sufficiently large power
by an element of positive first value. Put
\(e_n=x^n/(1+x^n)\). Then \(e_n\to0\) in \(v_1\)
and \(e_n\to1\) in \(v_2\). Thus
\(a+(b-a)e_n\) approximates any prescribed pair \(a,b\)
simultaneously; this proves the required two-valuation
approximation coefficient by coefficient.

If there were a monic irreducible separable polynomial
\(f\) of degree \(d>1\), irreducibility would hold on
a sufficiently small coefficient neighborhood for \(v_1\).
Here are the details of this often-used assertion. Choose
a root \(\alpha\), and choose \(\delta\) in the base
field of value larger than all the finitely many values
\(v_1(\alpha-\alpha')\), with \(\alpha'\ne\alpha\)
a conjugate, and sufficiently large to control the higher
Taylor coefficients. For a sufficiently close polynomial
\(g\), the polynomial
\[
 \frac{g(\alpha+\delta Y)}{\delta g'(\alpha)}
\]
has integral coefficients, constant term in the maximal
ideal, linear coefficient one and all higher coefficients
in the maximal ideal. Hensel's lemma in the finite
henselian extension generated by \(\alpha\) supplies
a root \(\beta\) of \(g\) with
\(v_1(\beta-\alpha)>v_1(\delta)\).
If an automorphism over the base fixes \(\beta\),
uniqueness of the extended henselian valuation implies
that it cannot move \(\alpha\): otherwise both terms
in \(\alpha-\alpha'=(\alpha-\beta)+(\beta-\alpha')\)
would have value strictly greater than its left-hand side.
Thus the field generated by \(\alpha\) is contained
in that generated by \(\beta\). Their degrees force
\(g\) to be irreducible. This proves the asserted open
irreducibility neighborhood.

On the other hand, choose a monic polynomial with \(d\)
distinct roots in the base field. A sufficiently small
coefficient neighborhood for \(v_2\) consists of split
polynomials, by Hensel's lemma at the distinct simple
roots, with disjoint root neighborhoods. Simultaneous
approximation produces a polynomial in both neighborhoods,
an impossibility. This proves the two-valuation fact.

Now fix a number field \(D\subset K_p\), let
\(\widetilde v\) be the chosen valuation on
\(\overline{\mathbb Q}\), and write
\(D_{\widetilde v}\subset G_D\) for its decomposition
subgroup. It is a proper subgroup. To see this concretely,
choose a rational prime \(q\ne p\) unramified in
\(D\), and an integer \(a\) satisfying
\[
 a\equiv1\pmod{p^3},\qquad a\equiv q\pmod{q^2}.
\]
Its \(q\)-valuations in \(D\) are one, so it is not
a square in \(D\); it is a square in \(\mathbb Q_p\)
by Hensel's lemma (also for \(p=2\)). Thus the chosen
local decomposition group fixes a root of this nontrivial
quadratic extension, and is proper in \(G_D\).

Let \(N/D\) be the normal closure of \(K_p/D\).
Since \(K_p\) is henselian, every element of \(G_{K_p}\)
preserves \(\widetilde v\); consequently
\[
 G_N\subset G_{K_p}\subset D_{\widetilde v}
                       \subsetneq G_D.
\]
The normality of \(G_N\) makes the distinct cosets of
\(D_{\widetilde v}\) give distinct conjugate valuations
on \(N\). Each is henselian, since \(N\) is algebraic
over the corresponding conjugate of \(K_p\). They have
rank one because their value groups are subgroups of
\(\mathbb Q\). The two-valuation fact makes \(N\)
algebraically closed, proving \(N=\overline{\mathbb Q}\).

The same proof gives a slightly stronger exclusion needed
for Claim~\ref{clm:FFnormal}. Choose a conjugate
\(gK_p\) for which the prolongation differs. Their
compositum has the two distinct henselian valuations,
so is \(\overline{\mathbb Q}\), and therefore
\[
                   G_{K_p}\cap gG_{K_p}g^{-1}=1.          \tag{FF.23.1}
\]
If \(K_p\) were finite over a normal algebraic extension
\(k/D\), both groups on the left would be open in the
normal group \(G_k\). Their intersection would be open,
forcing \(G_{K_p}\) to be finite. This is impossible:
the valuation of a root of \(X^n-p\) is \(1/n\),
so this polynomial is irreducible over \(K_p\) for every
\(n\), and its finite extension degrees are unbounded.
If \(K_p\) were finitely generated over such a \(k\),
it would be finite because it is algebraic over
\(\mathbb Q\). This excludes that enlarged positive
class too. The arithmetic formation classes are excluded
by the full normal closure and their arbitrarily large
symmetric-group obstruction; the geometric and function-field
classes require positive transcendence degree.

Finally, every negative field family in
Claim~\ref{clm:FFghosts} has a nontrivial centralizer over
it or a finite extension, whereas \(K_p\) satisfies
the opposite condition. It therefore lies on neither
of those earlier lists. Countability and transcendence
degree zero are the least possible numerical values
for an infinite characteristic-zero field; this says
nothing about a least field under inclusion. The failure
of Kummer-faithfulness also explains the exact point at
which Carlson--Stix Proposition~3.18 cannot be used to
upgrade the known Isom assertion to full topos reconstruction.
It is a failure of that attempted proof, not a counterexample
to reconstruction or to centralizer sufficiency. The
\(p\)-primary argument below proves reconstruction
without Kummer-faithfulness.
\end{proof}

\subsection{What point detection retains over the strict henselian test}

\begin{claim}{FF.24}{Proved}\label{clm:FFstrictdetection}
Fix a prime \(p\), and let \(K_p\) be the field of
Claim~\ref{clm:FFstricttest}.
\begin{enumerate}
\item For every finite extension \(L/K_p\),
\[
 \bigcap_{n\geq1}(L^\times)^{p^n}=\mu_{p'}\subset K_p^\times,
                                                        \tag{FF.24.1}
\]
where \(\mu_{p'}\) denotes all roots of unity of order
prime to \(p\).
\item There is a connected nontrivial finite \'etale Galois
cover of an elliptic curve over \(K_p\) whose pullback
at every closed point is completely split. Thus the
finite-quotient detection statement used in
Claim~\ref{clm:FFcore} is false over \(K_p\).
\item If \(X/K_p\) is of finite type and a pinned topos
morphism \(X_{\mathrm{\acute et}}\to
\mathbb G_{m,K_p,\mathrm{\acute et}}\) has constant
underlying closed-point function with value
\(c\in\mathbb G_m(K_p)\), it factors through the
closed subtopos \(c_{\mathrm{\acute et}}\).
For an integral finite-type curve \(X\), every
nonconstant underlying continuous point function has
finite fibres over closed target points.
\item Every field embedding \(K_p\hookrightarrow K_p\)
is an automorphism.
\end{enumerate}
Consequently neither the invisible prime-to-\(p\)
quotients in (2), a residue-tube point function, nor
a proper self-embedding supplies a reconstruction
counterexample.
\end{claim}

\begin{proof}
The unique valuation on a finite extension \(L\) is
henselian, discrete, and has residue field
\(\overline{\mathbb F}_p\). Normalize it to have value
group \(\mathbb Z\), and write \(e=v(p)>0\).
Hensel's lemma identifies \(\mu_{p'}\) with the
nonzero residue elements; these roots already lie
in \(K_p\). Suppose \(a\) belongs to the intersection
in (FF.24.1). Its valuation is divisible by every
\(p^n\), hence is zero. Let \(\zeta\in\mu_{p'}\)
be its residue lift. For each \(n\), choose
\(a_n\in L^\times\) with \(a_n^{p^n}=a\), and
let \(\zeta_n\) be its prime-to-\(p\) residue lift.
Uniqueness gives \(\zeta_n^{p^n}=\zeta\).
Thus \(b=a/\zeta=(a_n/\zeta_n)^{p^n}\), with
\(a_n/\zeta_n\in1+\mathfrak m_L\).
For \(s=v(u)>0\), the binomial formula gives
\[
 v((1+u)^p-1)\geq\min\{ps,e+s\}\geq s+1.
\]
Iteration yields \(v(b-1)\geq n+1\) for every
\(n\), so \(b=1\). Conversely, raising to \(p\)
is an automorphism of \(\mu_{p'}\), proving (1).

Choose an elliptic scheme \(\mathcal E\) over
\(\mathbb Z_{(p)}\): for odd \(p\) one may use
\(y^2=x^3-x\), whose discriminant is \(64\),
and for \(p=2\) use \(y^2+y=x^3\), whose
discriminant is \(-27\). Take \(\ell=2\) in
the first case and \(\ell=3\) in the second.
Over the strict henselization, \(\mathcal E[\ell]\)
is the constant group \((\mathbb Z/\ell)^2\).
On the generic fibre \(E/K_p\),
\([\ell]:E\to E\) is therefore a connected
finite \'etale Galois cover with that nontrivial
deck group. For a closed point with residue field
\(L\), properness extends the point to
\(\mathcal E(\mathcal O_L)\). Pulling back
\([\ell]\) gives a finite \'etale scheme over
the strictly henselian ring \(\mathcal O_L\),
so the pullback is a disjoint union of sections.
Using the origin to split the arithmetic fundamental
group, the resulting quotient
\[
 \pi_1(E)\longrightarrow (\mathbb Z/\ell)^2
\]
is nonzero and kills every point-section image
\(G_L\to\pi_1(E)\), including after every finite
base extension. This proves the precise failure in (2).

For (3), the inverse image of
\(\mathbb G_m\setminus\{c\}\) is an open subset
of \(X\) with no closed points. A finite-type scheme
over a field is Jacobson, so this open is empty.
The open--closed decomposition of the \'etale topos
then factors the morphism through the complementary
closed subtopos. If \(X\) is an integral curve,
the inverse image of any closed target point is
a closed subset of \(X\). Unless it is all of
\(X\), it is zero-dimensional and finite. This
also explains why a point function that collapses
an infinite residue tube cannot be the underlying
function of a nonconstant topos morphism.

Finally let \(\iota:K_p\hookrightarrow K_p\).
Since \(K_p/\mathbb Q\) is algebraic, extend
\(\iota\) to \(g\in G_{\mathbb Q}\), so
\(gK_p\subset K_p\). Transport the chosen
henselian valuation to \(gK_p\), and extend it
to the algebraic extension \(K_p/gK_p\).
It is again henselian of rank one. If it differs
from the original valuation of \(K_p\), the
two-valuation argument in
Claim~\ref{clm:FFstricttest} makes \(K_p\)
algebraically closed, a contradiction. They must
therefore coincide. Their unique prolongations
to \(\overline{\mathbb Q}\) coincide as well,
so \(g\) lies in the decomposition group of
the chosen \(p\)-adic valuation. This group
normalizes its inertia subgroup, whose fixed
field is exactly \(K_p\). Hence \(gK_p=K_p\),
and \(\iota\) is surjective.
\end{proof}

\begin{claim}{FF.25}{Proved}\label{clm:FFpdetector}
Let \(K/K_p\) be finite, let \(A/K\) be a semiabelian
variety, and put \(V_pA=T_pA\otimes_{\mathbb Z_p}\mathbb Q_p\).
Every nonzero continuous \(G_K\)-equivariant linear map
\[
                  \lambda:V_pA\longrightarrow\mathbb Q_p(1)
\]
is detected by a point: for some finite \(L/K\) and
\(a\in A(L)\),
\[
             \lambda_*\kappa_A(a)\ne0
                 \quad\hbox{in }H^1(G_L,\mathbb Q_p(1)). \tag{FF.25.1}
\]
If \(j:X\to A\) is a morphism based at \(x_0\in X(K)\),
with \(j(x_0)=0\), from a geometrically
integral finite-type \(K\)-scheme and \(j(X_{\bar K})\)
generates \(A_{\bar K}\), the detecting points may be
chosen in \(j(X(L))\). Finitely many such points,
over one finite extension, jointly detect every member of
\(\operatorname{Hom}_{G_K}(V_pA,\mathbb Q_p(1))\).
\end{claim}

\begin{proof}
We first explain the comparisons with complete fields.
A finite extension of \(K_p\) is an excellent henselian
discretely valued field of characteristic zero. Completion
induces an equivalence on finite separable extensions.
Indeed, approximate the coefficients of a separable primitive
polynomial over the completion by coefficients in the field;
the Krasner comparison proved in
Claim~\ref{clm:FFstricttest} identifies the resulting completed
extensions. For full faithfulness, a henselian field of
characteristic zero is relatively algebraically closed in its
completion: approximate a putative algebraic element closely
enough for the same simple-root Hensel argument to put it in
the field. A map between completed finite extensions therefore
sends their algebraic elements back into the corresponding
finite extensions. This proves the equivalence and identifies
the absolute Galois groups. Smooth points over the original
field are dense in the points over the completion: use an
\'etale coordinate chart and Hensel's lemma, then the density
of the field in its completion.

Suppose first that \(A=B\) is abelian. After a finite extension
it has semistable reduction. The exact source is
\citet[IX, Theorem~3.6]{FF_SGA7I1972}: over a regular connected
noetherian scheme of dimension one, an abelian generic fibre
acquires semistable reduction after a finite Galois extension.
Apply it to the valuation ring, and then complete. Write
\(R=\mathcal O_{\widehat L}\), with residue field
\(\overline{\mathbb F}_p\). The finite part of the N\'eron
model supplies a \(p\)-divisible group \(G/R\) such that
\[
 0\longrightarrow V_pG\longrightarrow V_pB
       \longrightarrow W\longrightarrow0,             \tag{FF.25.2}
\]
where \(W\) is unramified and hence trivial over this
strictly henselian base. Moreover, the connected part
\(G^0\) is the \(p\)-divisible group of the formal
completion of the connected N\'eron model at the identity.
These are precisely
IX (5.1.4)--(5.1.5), (5.5.8), and Proposition~5.6 of
\cite{FF_SGA7I1972}; Proposition~5.6 includes the residue
prime, not just primes different from it.

The \(p\)-adic cyclotomic character on \(G_{\widehat L}\)
has infinite image. For example, a primitive
\(p^n\)-th root of unity has
\(v(\zeta_{p^n}-1)=v(p)/(p^{n-1}(p-1))\), which is
incompatible with a fixed discrete value group as
\(n\to\infty\). Thus
\(\operatorname{Hom}_{G_{\widehat L}}(W,\mathbb Q_p(1))=0\),
and \(\lambda\) is nonzero on \(V_pG\).
Multiply \(\lambda\) by a power of \(p\) to make this
restriction integral. Tate's generic-fibre theorem
\citep[Theorem~4 and Corollary~1]{FF_Tate1967} extends it
uniquely to a morphism
\[
                       f:G\longrightarrow\mu_{p^\infty}
\]
over \(R\). The stated theorem applies to an integrally
closed noetherian domain with characteristic-zero fraction
field, so this complete DVR satisfies its hypotheses.
The restriction to \(G^0\) is nonzero. Otherwise \(f\)
would factor through the \'etale quotient, which is constant
over \(R\), whereas its generic Tate module has no nonzero
map to \(\mathbb Z_p(1)\).

The equivalence of connected \(p\)-divisible groups with
divisible formal Lie groups
\citep[Section~2.2, Proposition~1]{FF_Tate1967} therefore
gives a nonzero formal group homomorphism
\(\widehat f:\widehat B\to\widehat{\mathbb G}_m\).
In characteristic zero, formal logarithms turn a formal
group homomorphism into its linear differential. Hence
\(d\widehat f\ne0\). On a sufficiently small ball,
\(\log\widehat f\) is a nonzero continuous linear form
in logarithmic coordinates. There is a point of this ball
with nonzero image, and density supplies such a point
\(a\in B(L)\). Its image \(u=\widehat f(a)\) is a
principal unit in \(\widehat L\) with \(\log u\ne0\).
Its rational \(p\)-Kummer class is nonzero: if a power
\(u^{p^s}\) were infinitely \(p\)-divisible, the binomial
valuation proof of (FF.24.1), which applies equally to
the complete DVR, would give \(u^{p^s}=1\), contradicting
its logarithm.

For clarity, the Kummer comparison here is a comparison of
actual cocycles. Choose a compatible sequence of division
points of \(a\) in the formal group over the algebraic
closure. These points exist because multiplication by
\(p\) on a divisible formal group is an isogeny. Their
successive Galois differences give the Kummer cocycle
in \(T_pG\), whose image in \(T_pB\) is the usual
point-Kummer cocycle. Applying \(\widehat f\) gives
the compatible roots of \(u\), hence its multiplicative
Kummer cocycle. Thus the nonzero class just obtained is
\(p^s\lambda_*\kappa_B(a)\) for the fixed scaling power.
The completion comparison proves (FF.25.1) over \(L\).

For a general semiabelian variety, pass to a finite
extension splitting its torus in
\(1\to T\to A\to B\to0\). If \(\lambda|_{V_pT}\ne0\),
one coordinate of \(T\simeq\mathbb G_m^r\), evaluated
at a sufficiently small nontrivial principal unit, detects
it by (FF.24.1). If this restriction is zero, \(\lambda\)
factors through a nonzero functional of \(V_pB\).
The abelian case gives a detecting \(b\in B(L)\).
Hilbert's theorem 90 for the split torus gives
\(A(L)\twoheadrightarrow B(L)\), so a lift detects
\(\lambda\) by functoriality of Kummer cocycles.

Now assume that \(j(X_{\bar K})\) generates \(A_{\bar K}\).
Every point of \(A(\bar K)\) is a finite sum of signed
points of \(j(X(\bar K))\). To justify this, the dimensions
of the images of successive signed sum maps eventually
stabilize; their closure is the algebraic subgroup
generated by the image, hence is \(A\). A dominant sum
map has constructible image containing a nonempty open
\(U\). Since \(U\cap(a+U)\ne\varnothing\), every
\(a\) is the difference of two such sums. All finitely
many points involved are defined over a finite extension.
If every \(j(x)\) annihilated \(\lambda\), additivity
of the Kummer map would therefore annihilate every
point, contrary to the preceding proof.

Finally the space of functionals is finite-dimensional
over \(\mathbb Q_p\). Its intersections with the
kernels of point evaluations decrease to zero, so a
finite subcollection already has zero intersection.
Put these points over one finite compositum. Restriction
on \(H^1(-,\mathbb Q_p(1))\) to a subgroup of finite
index is injective, because restriction followed by
corestriction is multiplication by that index. Thus
the detecting property persists over this compositum.
\end{proof}

\subsection{Reconstruction over strictly henselian algebraic fields}

\begin{claim}{FF.26}{Proved}\label{clm:FFstrictreconstruct}
For every prime \(p\) and every finite extension
\(K/\operatorname{Frac}(\mathbb Z_{(p)}^{sh})\),
the field \(K\) satisfies \(\mathcal R(K)\).
Consequently every subfield of such a \(K\) reconstructs.
\end{claim}

\begin{proof}
Throughout, all topoi and their structural isomorphisms are
relative to the stated field. Finite scalar extension of a
pinned map is formed by the corresponding finite \'etale
slice. We may make finite extensions during the proof:
if a unit on \(X_L\) realizes the original point function
over a finite extension \(L/K\), pass to the finite normal
closure of \(L/K\). Every conjugate unit has the same
geometric point function. Equality on geometric points of
the reduced scheme makes the units equal, so the unit
descends to \(K\). This reduction does not assume
reconstruction over any of the fields in advance.

Let \(X/K\) be regular, quasiprojective and connected, with
\(x_0\in X(K)\), and let
\[
 \varphi:X_{\mathrm{\acute et}}\longrightarrow
                    \mathbb G_{m,K,\mathrm{\acute et}}
\]
be pinned. Regularity and the rational point make \(X\)
geometrically integral; in dimension zero it is \(\operatorname{Spec}K\),
where Claim~\ref{clm:FF1} applies directly. We henceforth
assume positive dimension. Its rational points are Zariski
 dense. Indeed, an \'etale coordinate chart at \(x_0\)
and Hensel's lemma lift every point of a sufficiently
small valuation polydisc in \(K^n\) to \(X(K)\).
The image of a proper closed subset under the quasi-finite
coordinate map has dimension less than \(n\). A product
of infinite valuation balls cannot lie in a proper
algebraic subset of \(\mathbb A^n_K\), by induction
on the variables of a nonzero polynomial. Thus every
nonempty open of \(X\) has a rational point. This
also holds after each finite extension and permits the
same argument on all inverse-image opens below. Write \(b(x)\) for its actual
value at a finite-extension point, and \(b_0=b(x_0)\).
Its pulled-back \(p\)-Kummer class is
\[
                   \xi=\varphi^*\kappa_p(T)
                          \in H^1(X,\mathbb Z_p(1)).
\]
In particular, evaluating \(\xi\) at \(x\in X(L)\)
gives \(\kappa_p(b(x))\), for every finite \(L/K\).

\emph{1. Integral boundary residues.}
We first prove a calculation requiring only a compatible
\(p\)-Kummer class with actual values at all finite-extension
points, rather than an entire topos map. On a punctured
rational line, after splitting its punctures over a finite
extension, normalization at a rational point writes the
class as
\[
 \kappa_p(c)+\sum_{i=1}^r a_i
       \kappa_p\left(\frac{T-r_i}{t_0-r_i}\right),
                  \qquad a_i\in\mathbb Z_p.             \tag{FF.26.1}
\]
This follows from the Kummer sequence and the vanishing of
the Picard group of a localization of \(K[T]\).
Let \(\pi\) be a uniformizer of this finite extension,
and evaluate at \(r_i+\pi^N\) and \(r_i+\pi^{N+1}\),
where \(N\) is larger than the valuations of all nonzero
root differences. Applying the valuation map on
\(H^1(K,\mathbb Z_p(1))\), the difference between these
two evaluations is \(a_i\): the other factors have
constant valuation. Both actual values have ordinary
integer valuation. Thus \(a_i\in\mathbb Z\). The residue
at infinity, the negative sum of the finite residues,
is integral as well.

For a smooth curve and a chosen rational boundary point
\(q\), Riemann--Roch gives a function whose only pole
is \(q\). Remove its branch values and the images of
all boundary points. The resulting map to a punctured
line is finite \'etale. Corestriction of the class has
actual values: at each point of the line these are the
products of the field norms of the actual values in its
fibre. The residue of a norm is the sum of the residue
degrees times the source residues. It has no factor
equal to the ramification index. Since \(q\) is the
only point above infinity, that residue is exactly
the residue at \(q\). Hence it is integral by (FF.26.1).
This argument is valid after making a boundary point
rational by a finite extension.

Choose a smooth projective compactification \(\bar X\)
with normal-crossing boundary, available in characteristic
zero. At a smooth point of a chosen boundary component,
away from the other components, take \'etale coordinates
in which that component is the first coordinate divisor.
Fix the other coordinates and take the curve component
through the point. Its local intersection multiplicity
with the chosen divisor is one. Normalize its closure
and apply the preceding curve argument. Thus each residue
\(m_i\) of \(\xi\) at a geometric boundary component
\(D_i\) belongs to \(\mathbb Z\). After another finite
extension all components and their residue data are defined
over the base. Put \(D=\sum_i m_iD_i\).

\emph{2. Forcing arithmetic equivariance over a finite local field.}
Let \(j:X\to A\) be the generalized Albanese map based
at \(x_0\), so \(j(x_0)=0\). The geometric part of
\(\xi-\kappa_p(b_0)\) is a continuous integral functional
\[
                  \lambda:T_pA\longrightarrow\mathbb Z_p(1)
\]
equivariant for \(H_0=G_K\). The Albanese comparison and
its finite obstruction are exactly those in
\citet[Proposition~4.1 and Corollary~4.2]{FF_SpiessSzamuely2003};
the finite obstruction disappears on taking the inverse
limit, since the Tate module of a finite group is zero.
The map \(j(X_{\bar K})\) generates \(A_{\bar K}\).

Choose a number field \(D_0\subset K\) containing all
the finite displayed geometric data and \(b_0\), and
such that \(D_0K_p=K\). Such a field exists because
\(K/K_p\) is finite. The chosen prolongation of the
\(p\)-adic valuation gives a finite extension
\(M_0=(D_0)_v\) of \(\mathbb Q_p\). The field \(K\)
is the fraction field of the corresponding strict
henselization of \(D_0\): each of these two strict
henselian fields contains the other by the universal
property. The completion/Krasner comparison therefore
identifies
\[
                  H_0=G_K=I_{M_0}\triangleleft G_{M_0}.
                                                               \tag{FF.26.2}
\]
These are compatible identifications for the Tate modules
of the varieties defined over \(D_0\); their torsion
points and all relevant point-Kummer cocycles are algebraic
over the number field.

Apply Claim~\ref{clm:FFpdetector} to the finite-dimensional
space
\[
                  S=\operatorname{Hom}_{H_0}
                       (V_pA,\mathbb Q_p(1)).
\]
Choose finitely many points \(x_i\), over one finite
extension \(E/K\), whose Kummer evaluations jointly
detect \(S\). Include their actual values \(b(x_i)\)
in the finite data, and choose a number field
\(D_1\subset E\) containing \(D_0\), all these data,
and a generator for \(E/K_p\). Then
\(E=D_1K_p\), and with \(M_1=(D_1)_v\) we have
\[
        H_1=G_E=I_{M_1}\triangleleft G_{M_1},
                  \qquad G_{M_1}\subset G_{M_0}.         \tag{FF.26.3}
\]
Restriction to \(H_1\) preserves joint detection,
by restriction--corestriction over \(\mathbb Q_p\).

For \(g\in G_{M_1}\), let
\(g\lambda=\chi_p(g)\lambda\rho_A(g)^{-1}\).
The normality in (FF.26.2) shows that
\(\delta_g=g\lambda-\lambda\) still belongs to \(S\).
Put
\[
 c_i=\kappa_A(j(x_i))\in H^1(G_{M_1},V_pA),\qquad
 d_i=\kappa_p(b(x_i)/b_0)\in H^1(G_{M_1},\mathbb Q_p(1)).
\]
The actual point-value equality is
\[
                \lambda_*\operatorname{res}_{H_1}(c_i)
                        =\operatorname{res}_{H_1}(d_i). \tag{FF.26.4}
\]
An \(H^1\)-class restricted from a group to a normal
subgroup is invariant under that group's conjugation
action, including its action on the coefficient module.
Apply this elementary cocycle identity to both \(c_i\)
and \(d_i\) in (FF.26.4). It gives
\[
               (\delta_g)_*\operatorname{res}_{H_1}(c_i)=0
                                      \quad\hbox{for all }i.
\]
Joint detection proves \(\delta_g=0\). Thus \(\lambda\)
is \(G_{M_1}\)-equivariant. This argument uses cohomology
classes, so the conjugation coboundaries cause no problem;
it does not require a trivial cyclotomic character and
does not assert that arbitrary topos maps descend to
\(M_1\).

\emph{3. A power of the class is the class of a unit.}
We give the integral Picard calculation to avoid an
incorrect local Tate-Hom assertion for general 1-motives.
Let \(P=\operatorname{Pic}^0_{\bar X}\), and let
\(B_X\subset\bigoplus_i\mathbb ZD_i\) be the group
of boundary divisors algebraically equivalent to zero.
The Picard 1-motive is \([B_X\to P]\), dual to
\([0\to A]\). Its realization gives
\[
 0\longrightarrow T_pP\longrightarrow
 H^1(X_{\bar K},\mathbb Z_p(1))\longrightarrow
 B_X\otimes\mathbb Z_p\longrightarrow0.                \tag{FF.26.5}
\]
Here the last map sends the geometric class to its
boundary residues. The construction, including the
finite obstruction before inverse limit, is given in
the proof of \citet[Proposition~4.1]{FF_SpiessSzamuely2003}.
The field has characteristic zero, so its prime-to-characteristic
statement includes the present prime \(p\).

By Step~1, the residue vector \(D\) is integral, while
(FF.26.5) puts it in \(B_X\otimes\mathbb Z_p\).
Consequently its class in the finitely generated group
\((\bigoplus_i\mathbb ZD_i)/B_X\) is torsion of order
prime to \(p\). Choose an integer \(N>0\), prime
to \(p\), such that \(ND\in B_X\).
Step~2 says that the geometric class is
\(G_{M_1}\)-invariant. The existence of its invariant
multiple lifting \(ND\) in (FF.26.5) means
\[
                 [ND]\in\bigcap_n p^nP(M_1).           \tag{FF.26.6}
\]
One can check this directly in the 1-motive, without
any Hom theorem. At level \(p^n\), a lift of \(ND\)
is represented by a \(p^n\)-division point of
\([ND]\in P\). The fixed residue vector allows us
to choose its lattice component equal to \(ND\).
Galois invariance then makes that division point
\(M_1\)-rational: a change of representative would
change the lattice component by \(p^n b\), and
\(p^n b=0\) implies \(b=0\) in the torsion-free
lattice. This holds for every \(n\).

For an abelian variety over a finite extension \(M_1\)
of \(\mathbb Q_p\), the compact group \(P(M_1)\)
contains an open subgroup isomorphic, under its formal
logarithm, to a finite free \(\mathbb Z_p\)-module.
Its infinitely \(p\)-divisible subgroup is therefore
exactly its finite prime-to-\(p\) torsion subgroup.
For completeness, modulo the open subgroup there is
only a finite group. A fixed power of \(p\) kills
its \(p\)-primary part; after multiplying by its
prime-to-\(p\) exponent, an infinitely \(p\)-divisible
element lies in every \(p^n\) multiple of the free
\(\mathbb Z_p\)-module and is zero. The converse for
prime-to-\(p\) torsion is immediate.
Thus (FF.26.6) is killed by an integer \(d_0>0\)
prime to \(p\), and \(dD\) is principal, where
\(d=Nd_0\).

All divisor and Picard data were defined over the
number field \(D_0\). Their equality after extension
to \(M_1\) implies the equality over \(D_0\), since
extension of scalars is injective on points of a scheme.
The principal divisor has a rational function over
\(D_0\), uniquely after fixing its value at \(x_0\):
the one-dimensional space of functions with that divisor
descends, or equivalently one applies Hilbert 90 to its
scalar ambiguity. Choose
\(h\in\mathcal O(X)^\times\) with
\(\operatorname{div}_{\bar X}(h)=dD\) and
\(h(x_0)=b_0^d\).

The difference between \(d\lambda\) and the geometric
Kummer class of \(h\) has zero residues. It belongs,
by (FF.26.5), to \((T_pP)^{G_{M_1}}\), which is zero:
the torsion subgroup of \(P(M_1)\) is finite, by
the same compact Lie-group argument. Hochschild--Serre
and evaluation at \(x_0\) then give the integral identity
\[
                         d\xi=\kappa_p(h).              \tag{FF.26.7}
\]
In particular, for every point over a finite extension,
\[
                         b(x)^d=\zeta_xh(x),
                         \qquad\zeta_x\in\mu_{p'}\subset K_p.
                                                               \tag{FF.26.8}
\]
This uses the integral kernel (FF.24.1), not merely
the kernel after tensoring cohomology with \(\mathbb Q_p\).

\emph{4. Curves with large-degree points and a resultant.}
We record explicitly two elementary tools. Every nonempty
open of a geometrically integral curve over a finite
extension of \(K_p\) has closed points of unbounded
degree. Choose a nonconstant map from its proper
normalization to \(\mathbb P^1\). The polynomials
\(T^n-\pi\) give closed points of degree \(n\) on
the line, by Eisenstein's criterion. For arbitrarily
large \(n\) their support avoids the finite image
of the omitted points; any point in the fibre has
degree at least \(n\).

Second, if \(V\subset K(X)\) is a finite-dimensional
space of rational functions regular on a nonempty open
\(U\), after a finite extension there is a smooth
geometrically integral curve mapping into \(U\) on
which restriction is injective on \(V\). Choose finitely
many geometric points of \(U\) whose evaluations are
injective on \(V\); this is possible by successively
testing a nonzero function in the remaining kernel.
Put them over a finite extension. A geometrically
irreducible curve through them exists by
\citet[Theorem~1.8, Corollary~1.9]{FF_CharlesPoonen2016},
whose hypotheses allow an arbitrary base field and
a finite set of closed points. Normalize its reduced
structure. In characteristic zero the result is a
smooth curve. Evaluation at its points above the
chosen set proves the asserted injectivity.

Suppose now that the geometric class \(\lambda\)
is nonzero. Then \(h\) in (FF.26.7) is nonconstant.
The actual point function is nonconstant as well:
if all values were equal to a base-field constant,
(FF.26.8) would put all values of \(h\) in the
base field. A nonconstant regular function cannot
have this property, by restricting to a curve and
using closed points of arbitrarily large degree
on its image.
Thus \(U=\varphi^{-1}(\mathbb G_m\setminus\{1\})\)
is a nonempty open of \(X\). Apply Steps~1--3 on
\(U\) to the pullback of \(T-1\), using a
rational point of that open, which exists by the preceding
density argument.
They give a positive integer \(e\), a rational
function \(g\) that is a unit on \(U\), and
\[
                     (b(x)-1)^e=\eta_xg(x),
                           \qquad\eta_x\in\mu_{p'}.
                                                               \tag{FF.26.9}
\]
Consider the finite-dimensional space generated by
\(h^ig^j\), for \(0\leq i\leq e\) and
\(0\leq j\leq d\). Choose a curve as in the preceding
paragraph, so restriction to that curve is injective
on this space. Write
\[
              B=e\deg\operatorname{div}_\infty(h|_C)
                   +d\deg\operatorname{div}_\infty(g|_C).
\]
Choose a closed point \(x\in C\) of degree greater
than \(B\). The rational function
\[
 R=\operatorname{Res}_Z
       \bigl(Z^d-\zeta_xh,\ (Z-1)^e-\eta_xg\bigr)
                                                               \tag{FF.26.10}
\]
vanishes at \(x\), by (FF.26.8)--(FF.26.9), and
its pole divisor on \(C\) has degree at most \(B\).
It must vanish identically on \(C\). All its terms
belong to the chosen finite-dimensional space, with
coefficients in the base field: here it matters that
\(\zeta_x,\eta_x\in K_p\). Injectivity of restriction
therefore gives \(R=0\) in \(K(X)\).

There is a common root \(z\) in an algebraic closure
of \(K(X)\), satisfying
\(z^d=\zeta_xh\) and \((z-1)^e=\eta_xg\).
It is nonconstant because \(h\) is nonconstant.
Every conjugate \(z'\) has the form
\(z'=\alpha z\), \(z'-1=\beta(z-1)\), for
roots of unity \(\alpha,\beta\). Consequently
\((\alpha-\beta)z=1-\beta\). As \(z\) is
transcendental over the algebraic closure of the
constant field, this forces \(\alpha=\beta=1\).
Characteristic zero makes the algebraic extension
separable, so \(z\in K(X)\). Its power is a unit
on the normal scheme \(X\), hence \(z\) is a
unit on \(X\).

The geometric group \(H^1(X_{\bar K},\mathbb Z_p(1))\)
is torsion-free. From \(d\lambda=d\kappa_p(z)\)
we obtain \(\lambda=\kappa_p(z)\) geometrically.
Normalize \(z\) to have value \(b_0\) at \(x_0\).
Hochschild--Serre then yields a unit, now denoted \(h\),
with
\[
                             \xi=\kappa_p(h).            \tag{FF.26.11}
\]
This conclusion is proved whenever the geometric class
is nonzero, without assuming in advance that the
class of \(T-1\) is nonzero. The normalized unit
also descends through every finite extension made in
this construction, before its actual point function
has been identified. For a Galois conjugate, the ratio
of the two units has zero geometric \(p\)-Kummer class,
so every coefficient of its divisor on \(\bar X\)
is divisible by every \(p^n\), and hence zero.
Properness makes this ratio constant; its value at
\(x_0\) is one. Galois descent and then evaluation
at \(x_0\) in Hochschild--Serre establish
(FF.26.11) over the original base field.

\emph{5. The zero geometric class forces a constant point function.}
If \(\lambda=0\), normalization gives
\(\xi=\kappa_p(b_0)\). By (FF.24.1), every
actual value is in \(b_0\mu_{p'}\subset K\),
including values at points over arbitrary finite
extensions. Choose a proper finite extension
\(K'/K\) and \(\beta\in K'\setminus K\).
Since no value equals \(\beta\), the pullback of
\(T-\beta\) is defined on all of \(X_{K'}\).
If its geometric class were nonzero, Step~4 would
give a nonconstant unit \(g\) with
\(b(x)-\beta=\eta_xg(x)\),
\(\eta_x\in\mu_{p'}\subset K\). All values of
\(g\) would then lie in \(K'\), contradicting
nonconstancy by the large-degree-point argument.
Thus this geometric class is zero too. Normalization
at \(x_0\) gives
\[
              \frac{b(x)}{b_0}\in\mu_{p'},\qquad
              \frac{b(x)-\beta}{b_0-\beta}\in\mu_{p'}.
\]
The two scalar coefficients belong to \(K\),
and \(1,\beta\) are linearly independent over
\(K\). These equations force \(b(x)=b_0\) at
every geometric point.

\emph{6. Removing the remaining pointwise scalars.}
We now have a unit \(h\), normalized by
\(h(x_0)=b_0\), realizing \(\xi\) whenever
the point function is nonconstant, and a constant
function otherwise. On a smooth curve with
nonconstant \(h\), (FF.24.1) gives
\(b(x)=c_xh(x)\), \(c_x\in\mu_{p'}\subset K\).
The point function of \(T-1\) on its inverse-image
open is nonconstant: otherwise the values of the
nonconstant unit \(h\) on that open would all be
in the base field, since \(c_x\in\mu_{p'}\subset K\).
The same argument applies to every \(T-a\). Thus its
class is nonzero by
Step~5. It has a unit representative \(h_1\)
and \(b(x)-1=d_xh_1(x)\), with \(d_x\in\mu_{p'}\).
Choose a point of degree greater than both a common
pole-degree bound for \(c_xh-d_xh_1-1\) and the
degree of \(h:\bar C\to\mathbb P^1\), so that
\(h(x)\notin K\). The
same divisor argument gives an identity
\[
                           h_1=Ah+B,
                              \qquad A,B\in K^\times.
\]
At every point with \(h(x)\notin K\), comparison
of the coefficients of \(1,h(x)\) forces
\(c_x=c=-A/B\).

For each \(a\in K^\times\), repeat the argument
with \(T-a\). On points with \(h(x)\notin K\),
the known value \(b(x)=ch(x)\) shows that its
unit representative is proportional to \(ch-a\).
To justify the identity rather than merely pointwise
equality, choose a point of degree above both the common
pole bound and the degree of \(h\), hence with
\(h(x)\notin K\); it exists because
\(h\) is nonconstant. At a point with
\(h(x)\in K\), set \(a=ch(x)\). If
\(b(x)\ne a\), this point belongs to the open
where \(T-a\) pulls back to a unit, but that
unit is proportional to \(ch-a\) and vanishes
there, a contradiction. Hence \(b=ch\) at
every point. Normalization at \(x_0\) gives
\(c=1\).

For higher-dimensional \(X\), pass to a finite
extension and join \(x_0\) to an arbitrary geometric
point by the normalization of a geometrically
irreducible curve through them, using
\citep[Corollary~1.9]{FF_CharlesPoonen2016} again.
The two points lift after another finite extension.
The restricted class is the Kummer class of
\(h|_C\). If this is constant, Step~5 applies;
if nonconstant, the curve argument applies. Its
normalized unit agrees with \(h|_C\): a ratio
with zero \(p\)-Kummer class has every divisor
coefficient divisible by all \(p^n\), hence is
constant, and its value at the lifted \(x_0\)
is one. Thus \(b(x)=h(x)\) at every geometric
point of \(X\).

\emph{7. Reconstruction and descent.}
The finite-extension descent in the first paragraph
puts the resulting unit over the original field.
We have proved exactly the point-function condition
for every regular connected affine source with a
rational point in
\citet[Theorem~4.18(5)]{FF_CHW2024}.
That theorem gives \(\mathcal R(K)\) for all
finite-type schemes and all pinned relative topos
morphisms, with absolute weak normalization as
in the question. Finally \(K\) is algebraic over
each of its subfields, so
Claim~\ref{clm:FFdescent} gives the asserted
subfield consequence.
\end{proof}

\begin{claim}{FF.27}{Proved}\label{clm:FFdiscretealgebraic}
Let \(F/\mathbb Q\) be algebraic. If \(F\) admits a
nontrivial discrete valuation of residue characteristic
\(p>0\), then \(\mathcal R(F)\) holds. More precisely,
for any prolongation to \(\overline{\mathbb Q}\), the
compositum \(FK_p\) is finite over the strict
henselization field \(K_p\) attached to that prolongation.
Thus the fields in this statement are exactly the
subfields of finite extensions of some \(K_p\).
\end{claim}
\begin{proof}
Embed \(F\) in \(\overline{\mathbb Q}\) and fix a
prolongation of its valuation. Normalize it by
\(v(p)=1\). Since \(v(F^\times)\) is discrete and
contains \(\mathbb Z\), it is \(e^{-1}\mathbb Z\)
for an integer \(e\geq1\). The residue field is
algebraic over \(\mathbb F_p\): every element is
algebraic over \(\mathbb Q\), and restriction to
the number field it generates has finite residue field.

The fraction field \(F^{sh}\) of the chosen strict
henselization of the valuation ring embeds in
\(\overline{\mathbb Q}\). Henselization is immediate,
and strict unramified extensions preserve the value
group. Equivalently, the valuation ring of \(F\)
is a DVR, its strict henselization is a DVR with the
same uniformizer, and its residue field is
\(\overline{\mathbb F}_p\); see
\citep[Corollary~3.3.26]{FF_ADH2017} and
\citep[Tags 04GP, 0AP3]{FF_StacksHensel}.
The chosen local map \(\mathbb Z_{(p)}\to
\mathcal O_{F^{sh}}\) extends to the chosen strict
henselization by its universal property. Consequently
\[
 K_p\subset FK_p\subset F^{sh},
       \qquad v((FK_p)^\times)\subset e^{-1}\mathbb Z.
                                                        \tag{FF.27.1}
\]
If \(L/K_p\) is finite inside \(FK_p\), then
\[
 [L:K_p]=e(L/K_p)\leq e.                                \tag{FF.27.2}
\]
Here the equality has no defect term: the
henselian-completion equivalence proved in
Claim~\ref{clm:FFpdetector} preserves the degree;
the resulting finite extension of the complete DVR
has algebraically closed residue field. Its integral
closure is finite free, and comparison of lengths
modulo the base uniformizer gives degree equal to
ramification index times residue degree, the latter
being one. The inequality follows from (FF.27.1).

The compositum \(FK_p\) is the directed union of
these finite extensions. Choose one of maximal
degree, which exists because their degrees are
positive integers bounded by \(e\). Its compositum
with any other member still satisfies (FF.27.2),
so maximality forces that other member to be
contained in it. Hence \(FK_p/K_p\) is finite.
Claim~\ref{clm:FFstrictreconstruct} proves
\(\mathcal R(FK_p)\), and the algebraic descent
of Claim~\ref{clm:FFdescent} proves
\(\mathcal R(F)\).

Conversely a subfield of a finite extension of
\(K_p\) inherits a discrete mixed-characteristic
valuation: its value group is a subgroup of a
cyclic group and is nonzero because it contains
\(v(p)>0\). This proves the final characterization.
\end{proof}

\subsection{Bounded wild ramification}

\begin{claim}{FF.28}{Proved}\label{clm:FFtamereconstruct}
Fix the prolongation of the \(p\)-adic valuation used
to define \(K_p\), and set
\[
 T_p=K_p\bigl(p^{1/m}:(m,p)=1\bigr)
        \subset\overline{\mathbb Q},                    \tag{FF.28.1}
\]
with compatible roots. Every finite extension \(K/T_p\),
and every subfield of such a \(K\), satisfies
\(\mathcal R\). In particular the positive conclusion
includes a henselian field with nondiscrete value group
\(\mathbb Z_{(p)}\).
\end{claim}
\begin{proof}
We give all modifications of Claim~\ref{clm:FFstrictreconstruct};
the scheme-theoretic criterion and its resultant argument
will be applied exactly as proved there.

\emph{1. Field structure and the exact multiplicative kernel.}
The field \(T_p\) is the maximal tame extension of
\(K_p\). Indeed a finite tame extension of a strictly
henselian DVR is generated by a root of its uniformizer
of order prime to the residue characteristic. One can
see this directly: if \(\pi\) is a uniformizer upstairs
and \(p=u\pi^m\), the unit \(u\) has an \(m\)-th
root by Hensel's lemma, so the extension contains an
\(m\)-th root of \(p\), and the degrees agree.
See also \citep[Lemmas~15.116.2 and 15.116.7, Tags
09EV and 0EXW]{FF_StacksRamification}.
Its Galois group over \(K_p\) is
\(\prod_{\ell\ne p}\mathbb Z_\ell(1)\).

Every finite \(K/T_p\) can be written \(L_0T_p\)
with \(L_0/K_p\) finite. Write
\([L_0:K_p]=e=p^a m\), with \((m,p)=1\).
The residue field is algebraically closed, so
\([L_0:K_p]=e(L_0/K_p)\). If
\(p=u\pi^e\) for a uniformizer \(\pi\) of \(L_0\),
extracting an \(m\)-th root of the unit \(u\)
shows that \(L_0\) contains an \(m\)-th root of
\(p\). Consequently \(L_0\cap T_p\) has degree
exactly \(m\): its degree is prime to \(p\) and
divides \(e\), and the displayed root gives the
opposite lower bound. The Galois degree formula gives
\[
 [K:T_p]=p^a,
 \quad v(K^\times)=p^{-a}\mathbb Z_{(p)},
 \quad\overline K=\overline{\mathbb F}_p.              \tag{FF.28.2}
\]
For the value group, adjoining roots of \(p\) of
orders prime to \(p\) gives all of
\(e^{-1}\mathbb Z_{(p)}=p^{-a}\mathbb Z_{(p)}\),
while each finite step over \(L_0\) has degree and
ramification index prime to \(p\) and introduces
no further \(p\)-denominator.
Moreover \(K\) is the maximal tame extension of
\(L_0\). All its finite steps over \(L_0\) are tame.
Conversely its value group is divisible by every
integer prime to \(p\), and its residue field is
algebraically closed. Given such an integer \(n\),
any element \(x\in K^\times\) has an \(n\)-th root:
choose \(y\) with \(nv(y)=v(x)\) and extract an
\(n\)-th root of the unit \(x/y^n\) by Hensel.
This contains every tame uniformizer extension of
\(L_0\).

We also have the exact formula
\[
       \bigcap_{n\geq1}(K^\times)^{p^n}=\mu_{p'}.
                                                        \tag{FF.28.3}
\]
For a given element \(b\) in the left side, enlarge
\(L_0\) inside \(K\) to contain \(b\). The extension
\(K/L_0\) remains Galois pro-prime-to-\(p\), being
obtained by base change from \(T_p/K_p\). If
\(c^{p^n}=b\) with \(c\in K\), put
\(r=[L_0(c):L_0]\), so \((r,p)=1\). The norm
identity gives \(N(c)^{p^n}=b^r\). Choose integers
\(u,v\) with \(ur+vp^n=1\); then
\(b=(N(c)^u b^v)^{p^n}\). Thus \(b\) belongs to
the \(p\)-divisible kernel of \(L_0^\times\),
which is \(\mu_{p'}\) by Claim~\ref{clm:FFstrictdetection}.
The reverse inclusion is immediate. This proof
avoids using a discrete valuation estimate on the
dense value group in (FF.28.2).

\emph{2. Functional detection survives.}
Let \(A/K\) be semiabelian and let
\(0\ne\lambda\in\operatorname{Hom}_{G_K}
(V_pA,\mathbb Q_p(1))\). Descend \(A\) to a
finite extension \(L_0/K_p\) inside \(K\),
as in Step~1. The finite-dimensional vector space
\[
 S=\operatorname{Hom}_{G_K}(V_pA,\mathbb Q_p(1))
\]
is stable under \(G_{L_0}\), and its action factors
through the pro-prime-to-\(p\) group
\(G_{L_0}/G_K\). Its image in
\(\operatorname{GL}(S,\mathbb Q_p)\) is finite:
a compact linear group preserves a \(\mathbb Z_p\)-lattice;
the kernel of reduction modulo \(p\) (modulo \(4\)
if \(p=2\)) is pro-\(p\), whose intersection with
a pro-prime-to-\(p\) group is trivial. Hence after
a finite extension \(L_1/L_0\) inside \(K\),
\(\lambda\) is \(G_{L_1}\)-equivariant.
Claim~\ref{clm:FFpdetector} gives a finite extension
\(M/L_1\) and \(x\in A(M)\) with
\(\lambda_*\kappa_A(x)\ne0\).

This class stays nonzero over \(MK\). Indeed
\(G_{MK}\) is normal in \(G_M\) with
pro-prime-to-\(p\) quotient. Its invariant subspace
in \(\mathbb Q_p(1)\) is zero: a finite extension
of \(K\) still has bounded \(p\)-denominators in
its value group by (FF.28.2), whereas the valuations
of \(\zeta_{p^n}-1\) have unbounded
\(p\)-denominators. Thus the cyclotomic character
has infinite image there. Inflation--restriction
therefore makes
\[
 H^1(G_M,\mathbb Q_p(1))\longrightarrow
 H^1(G_{MK},\mathbb Q_p(1))
\]
injective. Alternatively continuous cohomology of
a pro-prime-to-\(p\) group with a finite-dimensional
\(p\)-adic representation vanishes in degree one,
by its finite image and averaging. The signed-sum
and finite-dimensionality arguments of
Claim~\ref{clm:FFpdetector} now give the same finite
joint detector on a pointed generating variety.

\emph{3. Integral residues become rational with prime-to-\(p\)
denominators.}
Normalize the valuation of \(K\) to have image
\(\mathbb Z_{(p)}\), choosing \(\pi\in K\)
of value one. The punctured-line calculation in
Step~1 of Claim~\ref{clm:FFstrictreconstruct}
uses the rational points \(r_i+\pi^N\) and
\(r_i+\pi^{N+1}\). For sufficiently large \(N\),
all other factor valuations agree at these two
points. The difference of the valuations of their
actual values is now in \(\mathbb Z_{(p)}\),
while its image in the \(p\)-adic completion is
exactly the coefficient \(a_i\in\mathbb Z_p\).
Thus \(a_i\in\mathbb Z_{(p)}\).
The finite-map norm argument for curves and the
transverse-curve argument for higher dimensions
are unchanged; every boundary coefficient is
in \(\mathbb Z_{(p)}\). There are finitely many,
so multiplying the class and its boundary by
one integer \(N\) prime to \(p\) makes all
coefficients integral. The normalized actual
values of this multiplied class are \(b(x)^N\).

\emph{4. Arithmetic equivariance and the unit power.}
Descend all source data to a number field
\(D_0\subset K\) with \(D_0T_p=K\), and let
\(M_0\) be its completion at the chosen place.
Step~1 identifies \(K\) with the maximal tame
extension of the chosen henselization of \(D_0\), so
\[
            G_K=P_{M_0}\triangleleft G_{M_0},          \tag{FF.28.4}
\]
where \(P\) denotes wild inertia. After selecting
the finite jointly detecting points and their
actual values, enlarge to a number field \(D_1\)
with \(D_1T_p=E\), as in Step~2 of
Claim~\ref{clm:FFstrictreconstruct}. Then
\(G_E=P_{M_1}\triangleleft G_{M_1}\), where
\(M_1=(D_1)_v\). The same conjugation identity
for the point and value cocycles, followed by
the detector in Step~2 above, makes the geometric
functional \(G_{M_1}\)-equivariant. No topos
map is asserted to descend to \(M_1\).

Apply the Picard 1-motive argument of Step~3 of
Claim~\ref{clm:FFstrictreconstruct} to the
multiplied boundary \(ND\). Its integral lattice
coordinate, fixed at every finite Tate level,
gives rational division points over the finite
\(p\)-adic field \(M_1\). The compact group
\(\operatorname{Pic}^0(M_1)\) has only prime-to-\(p\)
torsion in its infinitely \(p\)-divisible subgroup,
and its Tate invariants vanish. Consequently, for
some integer \(d>0\) prime to \(p\),
\[
              d\xi=\kappa_p(h),\qquad
              b(x)^d=\zeta_x h(x),\quad
              \zeta_x\in\mu_{p'}\subset K.            \tag{FF.28.5}
\]
Here \(d\) includes the initial multiplier \(N\).
The last equality is integral and follows from
(FF.28.3).

\emph{5. The remaining reconstruction steps.}
The only numerical input to Steps~4--6 of
Claim~\ref{clm:FFstrictreconstruct}, beyond
(FF.28.3)--(FF.28.5), is the existence of closed
points of arbitrarily large degree on a smooth
curve open. The polynomials
\(Z^{p^n}-\pi\) now supply base points of degree
\(p^n\): adjoining a root enlarges
\(\mathbb Z_{(p)}\) to
\(p^{-n}\mathbb Z_{(p)}\), of index \(p^n\),
so the polynomial has its full degree. A finite
map from the proper curve to \(\mathbb P^1\)
and avoidance of the finitely many omitted images
then supply the required points upstairs.
All the other inputs hold as before: \(K\) is
henselian of rank one, so smooth pointed varieties
have Zariski-dense rational points; it contains
\(\mu_{p'}\); every finite extension has the same
properties by (FF.28.2); and Charles--Poonen's
curve-through-points theorem applies in
characteristic zero. Thus the two-function
resultant kills the power \(d\), the proper
finite extension and \(T-\beta\) argument settles
zero geometric functionals, and the \(T-a\)
argument removes the pointwise roots of unity.
Their bounds and the finite-extension descent are
exactly those proved in
Claim~\ref{clm:FFstrictreconstruct}. This proves
its point-function conclusion here. CHW
Theorem~4.18(5) gives \(\mathcal R(K)\), and
algebraic descent gives the assertion for subfields.
\end{proof}

\begin{claim}{FF.29}{Proved}\label{clm:FFboundedwild}
Let \(F/\mathbb Q\) be algebraic, and choose a
\(p\)-adic place, normalized by \(v(p)=1\).
Suppose that for some integer \(a\geq0\),
\[
                 v(F^\times)\subset
                       p^{-a}\mathbb Z_{(p)}.          \tag{FF.29.1}
\]
Then \(FT_p/T_p\) is finite, and \(\mathcal R(F)\)
holds. Equivalently this condition says that the
\(p\)-parts of the ramification indices of all
finite number subfields, at the chosen place,
are bounded. These fields are exactly the subfields
of finite extensions of \(T_p\).
\end{claim}
\begin{proof}
Take the chosen henselization and strict residue
extension \(F^{sh}\). Henselization is immediate,
and the unramified extensions used in strict
henselization preserve its value group. The
chosen \(K_p\) is contained in \(F^{sh}\).
After adjoining roots of \(p\) of all orders
prime to \(p\), we obtain an extension containing
\(FT_p\) whose value group is still contained
in \(p^{-a}\mathbb Z_{(p)}\). To check the last
assertion even for this nondiscrete valuation,
note that its residue field is algebraically
closed and its units have all roots of order
prime to \(p\) by Hensel. For \((m,p)=1\),
the Kummer class of \(p\) modulo \(m\)-th powers
therefore has exactly the order of \(v(p)\)
in the value group modulo \(m\). Adjoining an
\(m\)-th root adds only the corresponding
prime-to-\(p\) denominators to that value group.

For every finite \(L/T_p\) inside \(FT_p\),
(FF.28.2) says
\([L:T_p]=[v(L^\times):\mathbb Z_{(p)}]\).
Thus \([L:T_p]\leq p^a\). The maximal-degree
argument of Claim~\ref{clm:FFdiscretealgebraic}
shows that \(FT_p/T_p\) is finite.
Claim~\ref{clm:FFtamereconstruct} and algebraic
descent prove \(\mathcal R(F)\). Conversely,
(FF.28.2) shows that every subfield of a finite
extension of \(T_p\) satisfies (FF.29.1).
\end{proof}

\begin{claim}{FF.30}{Proved}\label{clm:FFpadicsubfields}
Let \(F\subset\overline{\mathbb Q}_p\) be a field
of finite transcendence degree over \(\mathbb Q\).
Suppose that, for the induced valuation normalized
by \(v(p)=1\),
\[
                  v(F^\times)\subset
                         p^{-a}\mathbb Z_{(p)}
                    \quad\text{for some }a\geq0.       \tag{FF.30.1}
\]
Then \(\mathcal R(F)\) holds. In particular this
applies when the induced valuation on \(F\) is
discrete. No henselianity assumption is required.
\end{claim}
\begin{proof}
Choose a finitely generated field
\(D\subset F\) containing a transcendence basis
of \(F/\mathbb Q\). Its closure inside
\(\overline{\mathbb Q}_p\) is a finite extension
\(M/\mathbb Q_p\): it is generated topologically
by \(\mathbb Q_p\) and the finitely many
\(\mathbb Q_p\)-algebraic generators of \(D\).
Let \(K_D\) be the chosen strict henselization
field of the valuation ring of \(D\), inside
its algebraic closure in
\(\overline{\mathbb Q}_p\), and let
\(T_D\) be the maximal tame extension of \(K_D\)
there. Completion and the chosen prolongation
identify
\[
               G_{K_D}=I_M,\qquad G_{T_D}=P_M.         \tag{FF.30.2}
\]
Indeed the henselization of \(D\) is dense in
\(M\) and has the same finite separable extensions
by the Krasner--Hensel comparison proved in
Claim~\ref{clm:FFpdetector}. Strict henselization
then fixes inertia, and adjoining the prime-to-\(p\)
roots of a uniformizer fixes wild inertia.

We verify that Claims~\ref{clm:FFpdetector},
\ref{clm:FFstrictreconstruct} and
\ref{clm:FFtamereconstruct} apply to finite
extensions of \(K_D\) and \(T_D\), respectively.
The first proof uses only a henselian discrete
valuation of characteristic zero, algebraically
closed residue field of characteristic \(p\),
its completion, semistable finite-part
\(p\)-divisible groups, and Tate's theorem over
the completed DVR. These hypotheses hold here.
The exact multiplicative kernel follows from
the same discrete binomial estimate, with
\(\mu_{p'}\) supplied by strict henselization.

In the arithmetic-equivariance step of
Claim~\ref{clm:FFstrictreconstruct}, replace
its number fields \(D_i\) by finite extensions
of the fixed \(D\) containing the finitely
many model coefficients, detector coordinates,
and actual point values. Such fields suffice
because all the fields under consideration are
algebraic over \(D\). Their completions
\(M_i\) are finite \(p\)-adic fields, and
(FF.30.2) gives exactly the required normal
inertia subgroups. The finite-point conjugation
argument is therefore unchanged. The Picard
1-motive step uses compactness only for
\(\operatorname{Pic}^0(M_i)\), not for the
uncompleted model field. Its resulting torsion
or zero equality for a \(D_i\)-rational Picard
point descends from \(M_i\) to \(D_i\), since
field extension is faithfully flat. A principal
boundary divisor is represented by a normalized
function over \(D_i\) by the same Hilbert~90
argument. Thus the unit power, resultant,
zero-functional argument, scalar removal and
finite-extension descent all apply.

For \(T_D\), the norm proof of the multiplicative
kernel and finite-image descent of the functional
space use only the pro-prime-to-\(p\) quotient
of inertia by wild inertia. The valuation and
tame-root calculations in
Claim~\ref{clm:FFtamereconstruct} likewise use
only the DVR \(K_D\). The argument gives
\(\mathcal R(K)\) for every finite extension
\(K/T_D\).

Finally \(F/D\) is algebraic. The bounded
\(p\)-denominator argument of
Claim~\ref{clm:FFboundedwild}, applied with
\(K_D\) in place of \(K_p\), shows that
\(FT_D/T_D\) is finite. The normalization of
its value group may include the fixed
ramification index of \(M/\mathbb Q_p\);
this changes the finite bound by a fixed factor
and does not affect the argument. Therefore
\(\mathcal R(FT_D)\) holds. The extension
\(FT_D/F\) is algebraic because \(D\subset F\)
and \(T_D/D\) is algebraic. Algebraic descent
proves \(\mathcal R(F)\).
\end{proof}

\subsection{The next cyclotomic test and its undetected functionals}

\begin{claim}{FF.31}{Proved}\label{clm:FFcyclotomictest}
For every prime \(p\), put
\[
                   U_p=T_p(\mu_{p^\infty})
                       \subset\overline{\mathbb Q}.     \tag{FF.31.1}
\]
This countable, transcendence-degree-zero field has
henselian value group \(\mathbb Q\), residue field
\(\overline{\mathbb F}_p\), and contains every root
of unity. For every finite extension \(E/U_p\),
\[
           C_{\operatorname{Aut}(\overline{\mathbb Q})}
                         (G_E)=1.                       \tag{FF.31.2}
\]
It lies on none of the sufficient field lists
preceding this claim, including the bounded-wild class,
and none of the negative centralizer-ghost lists.
Thus \(U_5\) was a specified next test with the
smallest possible cardinality and transcendence degree
in characteristic zero; no inclusion-minimality is
claimed. Claim~\ref{clm:FFcyclotomicreconstruct}
subsequently proves reconstruction for every finite
extension of \(U_p\).
\end{claim}
\begin{proof}
Henselianity and the residue statement follow by
algebraic extension from \(K_p\). The tame
uniformizer roots supply all prime-to-\(p\)
denominators in the value group, while
\(v(\zeta_{p^n}-1)=1/(p^{n-1}(p-1))\)
supplies arbitrarily large \(p\)-denominators.
Their sums generate \(\mathbb Q\).

We first record the group input precisely.
For a finite \(p\)-adic field \(M\), its wild
inertia \(P_M\) is free pro-\(p\) of countably
infinite rank by
\citet[Theorem~2(ii), p.~463]{FF_Iwasawa1955}.
Closed subgroups of a free pro-\(p\) group are
free pro-\(p\); in a finite-rank free pro-\(p\)
group the centralizer of a nontrivial element
is procyclic. The latter statement and the
subgroup property are explicitly used and
proved/cited in
\citet[Theorem~5.1, proof, Case~I]{FF_KochloukovaZalesskii2011}.
Consequently, if \(H\subset P_M\) is nonabelian,
then \(C_{P_M}(H)=1\). Indeed if a nontrivial
\(c\) commuted with noncommuting
\(h_1,h_2\in H\), the closed subgroup
\(\langle c,h_1,h_2\rangle\) would be a
finite-rank free pro-\(p\) group in which both
\(h_i\) belong to the procyclic centralizer
of \(c\), a contradiction.

For
\(H=\ker(\chi_p:P_M\to\mathbb Z_p^\times)\),
nonabelianity is explicit. Choose four free
basis elements and map them to
\(1+e_{12},1+e_{23},1+e_{34},1+e_{45}\)
in the upper unitriangular group
\(\operatorname{UT}_5(\mathbb F_p)\).
Their first and second pair commutators map
to \(1+e_{13}\) and \(1+e_{35}\), whose
commutator is \(1+e_{15}\ne1\).
Both pair commutators belong to \(H\), since
\(\chi_p\) is abelian. This works also for
\(p=2\).

Now let \(E/U_p\) be finite. Choose a finite
extension \(L_0/K_p\) inside \(E\) such that
\(E=L_0U_p\), and then a number field model
with completion \(M\) for \(L_0\), as in
Claim~\ref{clm:FFstrictreconstruct}. We have
\[
 G_{L_0T_p}=P_M,
 \qquad G_E=\ker(\chi_p|_{P_M})=:H,
 \qquad H\triangleleft G_M.                            \tag{FF.31.3}
\]
In particular \(E\) is not algebraically closed.
Let \(C=C_{G_M}(H)\). Since \(H\) is normal,
\(C\) is normal in \(G_M\); the preceding
paragraph gives \(C\cap P_M=1\).
As \(P_M\) is normal too,
\([C,P_M]\subset C\cap P_M=1\).
But \(L_0T_p/T_p\) is finite, and
Claim~\ref{clm:FFtamereconstruct}, followed by
the necessary centralizer condition of
Claim~\ref{clm:FF6}, gives
\(C_{G_{\mathbb Q}}(P_M)=1\).
Therefore \(C=1\).

Suppose finally that
\(\sigma\in\operatorname{Aut}(\overline{\mathbb Q})
=G_{\mathbb Q}\) centralizes \(G_E\).
It preserves \(E\). Its transported rank-one
henselian valuation must equal the chosen one:
otherwise the two-valuation argument of
Claim~\ref{clm:FFstricttest} would make \(E\)
algebraically closed. Uniqueness of prolongation
from a henselian field then puts \(\sigma\)
in the chosen decomposition group
\(G_{\mathbb Q_p}\subset G_{\mathbb Q}\).
Since \(G_M\) is open, some positive power
\(\sigma^r\) belongs to \(G_M\), hence to
\(C=1\). The group \(G_{\mathbb Q_p}\) has
no torsion: its wild inertia is torsion-free
free pro-\(p\), and its tame quotient is an
extension of \(\widehat{\mathbb Z}\) by
\(\prod_{\ell\ne p}\mathbb Z_\ell\), both
torsion-free. Thus \(\sigma=1\), proving
(FF.31.2).

For completeness the list exclusions do not
rely only on failure of a sufficient hypothesis.
The two-valuation normal-closure proof of
Claim~\ref{clm:FFstricttest} applies to every
finite extension of \(U_p\). Thus its normal
closure over any number field of definition
is \(\overline{\mathbb Q}\), and a distinct
valuation-conjugate has compositum
\(\overline{\mathbb Q}\). Since its absolute
Galois group is infinite by (FF.31.3), the same
open-subgroup argument excludes being finite
over a proper normal algebraic field. It also
cannot be contained in a proper normal algebraic
extension of a number field: first adjoin that
number field to \(U_p\), a finite extension,
and apply the normal-closure assertion.
Positive-transcendence-degree function-field
classes are excluded by transcendence degree
zero. The group \(\mu_\infty\subset U_p^\times\)
is nontrivial divisible, excluding the
Kummer-faithful sub-\(q\)-adic classes. For
any valuation of residue characteristic \(q\),
its elements \(\zeta_{q^n}-1\) have unbounded
\(q\)-denominators, excluding the bounded-wild
classes at every prime, under any field embedding.
Finally (FF.31.2) excludes all negative
centralizer-ghost families. The proof of
reconstruction in Claim~\ref{clm:FFcyclotomicreconstruct}
requires the separate treatment of invisible
functionals developed next.
\end{proof}

\begin{claim}{FF.32}{Proved}\label{clm:FFcyclotomicdetector}
The functional-detection assertion of
Claim~\ref{clm:FFpdetector} fails over the
explicit field \(U_5\). More precisely there
is an elliptic curve \(A/U_5\) and a nonzero
\(G_{U_5}\)-equivariant functional
\[
               V_5A\longrightarrow\mathbb Q_5(1)
\]
whose push-forward of every point Kummer class,
over every finite extension of \(U_5\), is zero.
This is not a counterexample to
\(\mathcal R(U_5)\).
\end{claim}
\begin{proof}
Take
\(A:y^2=x^3-x\), whose discriminant \(64\)
is a unit at \(5\), and choose
\(i\in K_5\) with \(i^2=-1\) and
\(i\equiv2\pmod5\). The automorphism
\([i](x,y)=(-x,iy)\) defines the usual
\(\mathbb Z[i]\)-action. Put
\(\pi=2+i\) and \(\bar\pi=2-i\).
Then \([\pi]\) is an isogeny of degree
\(N(2+i)=5\), and its tangent map is
multiplication by \(2+i\), a unit at \(5\).
Hence \([\pi^n]\) is a finite etale isogeny
of the good elliptic model over
\(\mathcal O_{K_5}\), of degree \(5^n\).
Its kernel is constant over this strictly
henselian ring. The idempotents of
\(\mathbb Z[i]\otimes\mathbb Z_5
\simeq\mathbb Z_5\times\mathbb Z_5\)
split \(T_5A\); the \(\pi\)-primary factor
is the constant etale Tate module
\(T_\pi A\simeq\mathbb Z_5\).
Projection to this factor is nonzero.

For any finite \(L/U_5\), properness extends
\(a\in A(L)\) to the valuation ring
\(\mathcal O_L\). This ring is strictly
henselian, even though it is not discrete:
it is henselian and its residue field remains
\(\overline{\mathbb F}_5\). Pulling back
\([\pi^n]\) along the section gives a finite
etale scheme over \(\mathcal O_L\), hence
it splits. Thus \(a\) has a \(\pi^n\)-division
point in \(A(L)\) for every \(n\). Applying
the same splitting successively to each chosen
division point gives a compatible division
sequence in \(A(L)\). Its full Tate-module
Kummer cocycle is therefore zero, so the
Kummer class projected to \(T_\pi A\) is zero.
Since \(5=\pi\bar\pi\), the other factor
acts invertibly on the \(\pi\)-primary part,
so this is exactly the projection of the
ordinary \(5\)-adic point Kummer class.
But \(U_5\) contains all \(5\)-power roots
of unity. Its cyclotomic representation is
therefore trivial, and any chosen nonzero
identification
\(T_\pi A\otimes\mathbb Q_5\simeq
\mathbb Q_5(1)\) is \(G_{U_5}\)-equivariant.
This is the required undetected functional.

No pinned topos map realizing that functional
has been constructed. In particular a pinned
map with constant closed-point function must
factor through that closed point by the
Jacobson argument of
Claim~\ref{clm:FFstrictdetection}; such a map
cannot realize this nonzero geometric
functional. Thus the detector counterexample
does not decide reconstruction or refute
centralizer sufficiency.
\end{proof}

\subsection{Cyclotomic reconstruction after removing invisible functionals}

\begin{claim}{FF.33}{Proved}\label{clm:FFcyclotomiccohomology}
Let \(K/U_p\) be finite. Then the integral and
rational \(p\)-Kummer kernels of \(K^\times\)
are both \(\mu_\infty\). If \(M/\mathbb Q_p\)
is finite and
\(H=\ker(\chi_p|_{P_M})\), restriction
\[
 H^1(G_M,\mathbb Q_p(1))\longrightarrow
                  H^1(H,\mathbb Q_p(1))
                                                        \tag{FF.33.1}
\]
is injective. The same is true with inertia
\(I_M\) in place of \(G_M\).
\end{claim}
\begin{proof}
Write \(K=L_0U_p\) with \(L_0/K_p\) finite,
and enlarge \(L_0\) to contain a specified
\(a\in K^\times\). Both \(T_p/K_p\) and
\(K_p(\mu_{p^\infty})/K_p\) are abelian, so
\(K/L_0\) is Galois abelian. Its cyclotomic
character has infinite image.
If \(a\in\bigcap_n(K^\times)^{p^n}\), choose
a compatible sequence of roots in \(K\).
Such a sequence exists because \(K\) contains
all roots of unity and each finite root fibre
is nonempty; the transition maps can equally
be made surjective by adjusting by roots of
unity. Its integral Kummer cocycle factors
through \(\Gamma=\operatorname{Gal}(K/L_0)\).
For a \(\mathbb Q_p(1)\)-valued cocycle \(c\)
of an abelian group, commutativity gives
\[
 (\chi_p(s)-1)c(t)=(\chi_p(t)-1)c(s).
                                                        \tag{FF.33.2}
\]
Choose \(s\) with \(\chi_p(s)\ne1\). Equation
(FF.33.2) makes \(c\) a coboundary over
\(\mathbb Q_p(1)\). Thus \(\kappa_p(a)\)
is rationally zero already over \(L_0\), so
some \(p\)-power of \(a\) belongs to the
integral kernel \(\mu_{p'}\) of
Claim~\ref{clm:FFstrictdetection}. Hence
\(a\in\mu_\infty\). The converse is immediate,
and the rational kernel is the same because
a power of an element that is a root of unity
can only come from a root of unity.

For (FF.33.1), \(P_M/H\) is central in
\(G_M/H\): the cyclotomic character embeds
it in \(\mathbb Z_p^\times\), and conjugation
preserves that character. Choose a central
\(s\in P_M/H\) with \(\chi_p(s)\ne1\).
The same cocycle computation (FF.33.2), now
using only centrality of \(s\), shows that
\(H^1(G_M/H,\mathbb Q_p(1))=0\).
Inflation--restriction gives injectivity,
because \(H\) acts trivially on the coefficient.
The identical argument applies to \(I_M/H\).
No triviality of the tame action is assumed.
\end{proof}


\begin{claim}{FF.34}{Proved}\label{clm:FFcyclotomicsplitting}
Let \(A\) be a semiabelian variety over a finite extension
\(M/\mathbb Q_p\), put
\(H=\ker(\chi_p|_{P_M})\), and set
\[
 S=\operatorname{Hom}_H(V_pA,\mathbb Q_p(1)).
                                                        \tag{FF.34.1}
\]
After a finite extension of the finite local model, the
inertia action on this space has a canonical decomposition
\[
 S=S_0\oplus S_1,\qquad
 I_M|_{S_0}=1,\quad I_M|_{S_1}=\chi_p.
                                                        \tag{FF.34.2}
\]
Both summands are stable under the full local Galois group;
its unramified action on either summand need not be trivial.
For an algebraic model over a finite extension \(K/U_p\),
the summand \(S_1\) vanishes on the rational Kummer classes
of all points over all finite extensions of \(K\).
The summand \(S_0\) is jointly detected by finitely many
such point evaluations. If \(j:X\to A\) is pointed and
its geometric image generates \(A\), those evaluations
may all be taken at points \(j(x)\).
\end{claim}
\begin{proof}
We use the convention in which \(\mathbb Q_p(1)\)
has Hodge--Tate weight \(1\). The Hodge--Tate comparison
for a smooth variety with a smooth normal-crossing
compactification applies to a semiabelian variety.
Precisely, \citet[Proposition in Section~3.4 and
Theorem~3.6, especially proof~(d)]{FF_Beilinson2012}
gives the filtered, Galois-compatible de Rham
comparison for this nonproper variety; taking
the associated graded gives the Hodge--Tate
decomposition.
Here \(H^1_{\mathrm{\acute et}}(A_{\overline M},\mathbb Q_p)
=\operatorname{Hom}(V_pA,\mathbb Q_p)\); dualizing
its logarithmic Hodge filtration gives the stated
sign convention for the Tate module.
It makes \(V_pA\) Hodge--Tate with weights \(0,1\).
One can also read the weights from the torus and abelian
quotient after the comparison theorem supplies the
decomposition: the torus has weight \(1\), and the
abelian quotient has weights \(0,1\).
Consequently \(\operatorname{Hom}(V_pA,\mathbb Q_p(1))\)
and its \(G_M\)-stable subrepresentation \(S\) are
Hodge--Tate with weights \(0,1\).
Here closure under subobjects, duals, and tensor
products, with exactness on the weight spaces, is
\citet[Theorem~4.2]{FF_Fontaine2004}.

The quotient \(P_M/H\) is central in \(G_M/H\), as in
Claim~\ref{clm:FFcyclotomiccohomology}. The tame quotient
\(I_M/P_M\) is procyclic of pro-prime-to-\(p\) order.
A central extension with procyclic quotient is abelian,
so \(I_M/H\) is abelian. Its pro-prime-to-\(p\)
subgroup has finite image in \(\operatorname{GL}(S)\):
a compact linear group has an open pro-\(p\) subgroup,
whose intersection with a pro-prime-to-\(p\) image is
trivial. Kill this finite image and the finite torsion
in the cyclotomic character by a finite extension.
On a sufficiently small open subgroup of inertia,
the action is therefore
\[
           \rho_S(g)=\exp\bigl(\log\chi_p(g)\,B\bigr)
              \quad\text{for some }B\in\operatorname{End}_{\mathbb Q_p}(S).
                                                        \tag{FF.34.3}
\]
The same formula on \(S\oplus\mathbb Q_p(1)\)
shows that the inertia Lie algebra is the
one-dimensional space generated by \(\operatorname{diag}(B,1)\).
By \citet[Theorem~3.1, with \(d=0\)]{FF_Ohkubo2014},
the Sen operator belongs to the scalar extension of
this inertia Lie algebra to \(\mathbb C_p\).
Its component on \(\mathbb Q_p(1)\) is \(1\), so
it is exactly \(\operatorname{diag}(B,1)\).
The Hodge--Tate decomposition makes that operator
semisimple, with eigenvalues \(0,1\) on \(S\).
This is the precise criterion in
\citet[Section~2.7, printed p.~45]{FF_Fontaine2004},
with the same sign convention.
Thus \(B\) is semisimple over \(\mathbb Q_p\),
with those eigenvalues. Its two eigenspaces give
(FF.34.2), after shrinking inertia once more if necessary.
Conjugation by \(G_M\) fixes the normalized generator
\(B\), since it fixes the cyclotomic character on
the central group \(P_M/H\). Hence the two eigenspaces
are \(G_M\)-stable. This argument makes no assertion
that the unramified factors have finite image.

All finite changes of local model used here can be
realized by finite extensions of the chosen algebraic
model, by the henselization/completion equivalence.
For fixed \(S\), the finite inertia images can even
be killed inside the fixed field of \(H\), since
\(H\) acts trivially on \(S\oplus\mathbb Q_p(1)\).
Enlarging the finite model inside \(K\) in this way
does not change \(K\). Alternatively one may make
a finite extension of \(K\), as allowed throughout
the point-function argument.

We prove the asserted invisibility before detection.
For \(\lambda\in S_1\), the formula for the natural
action on Hom,
\[
        g\lambda=\chi_p(g)\lambda\rho_A(g)^{-1},
\]
and (FF.34.2) say that, after choosing a basis of
the target, the underlying linear map
\(\lambda':V_pA\to\mathbb Q_p\) is inertia-equivariant
for the trivial target action. After splitting the
torus by a further finite extension, it kills its
Tate module: that module is a sum of copies of
\(\mathbb Q_p(1)\), and inertia has infinite
cyclotomic image. Thus \(\lambda'\) factors through
the Tate module of the abelian quotient \(B_A\).

Every point in question is algebraic over \(\mathbb Q\).
It is therefore defined over a finite extension \(L/K_p\),
after adjoining its finitely many coordinates to a
finite field of definition. Make a further finite
extension so that \(B_A\) has semistable reduction.
The finite-part construction used in
Claim~\ref{clm:FFpdetector} supplies a connected
\(p\)-divisible group \(\mathcal G^0\) over the
complete strict valuation ring, whose formal group
is the formal completion of the identity of the
connected N\'eron model. Restrict \(\lambda'\) to
\(V_p\mathcal G^0\) and multiply by a power of \(p\)
to make it integral. The resulting map
\(T_p\mathcal G^0\to\mathbb Z_p\) is equivariant
over the completion of \(L\). By
\citet[Theorem~4 and Corollary~1]{FF_Tate1967}, it
extends uniquely to a map of \(p\)-divisible groups
\[
                  \mathcal G^0\longrightarrow\mathbb Q_p/\mathbb Z_p.
\]
This map is zero: every finite connected source
has zero map to a constant \'etale finite group.
Thus \(\lambda'\) kills \(V_p\mathcal G^0\).

For any \(a\in B_A(L)\), some nonzero integer
multiple of \(a\) lies in that formal completion.
Indeed, the component group is finite, and every
point of the connected reduction over
\(\overline{\mathbb F}_p\) has finite order:
it and the variety descend to a finite field.
Choose compatible successive \(p\)-division points
inside the formal group. Their differences under
Galois lie in \(T_p\mathcal G^0\), so applying
\(\lambda'\) kills the full compatible Kummer
cocycle, not just its classes at separate finite
levels. Division by the integer multiple is
legitimate with rational coefficients. Restriction
to the group of any finite extension of \(K\)
therefore gives zero. This proves invisibility
of \(S_1\).

For \(\lambda\in S_0\setminus\{0\}\), inertia
equivariance permits application of
Claim~\ref{clm:FFpdetector} over a finite extension
of \(K_p\). It gives a point over a further
finite extension whose image under \(\lambda\)
has nonzero rational Kummer class. The restriction
from this inertia group to its
\(\ker(\chi_p|_P)\) subgroup is injective by
Claim~\ref{clm:FFcyclotomiccohomology}, so the
class remains nonzero after adjoining \(U_p\).
The pointed-generation assertion follows from
the corresponding assertion in
Claim~\ref{clm:FFpdetector}. Finally, a finite
family of evaluations suffices on the
finite-dimensional space \(S_0\): successively
choose an evaluation nonzero on the remaining
common kernel until that kernel is zero.
\end{proof}

\begin{claim}{FF.35}{Proved}\label{clm:FFcyclotomicreconstruct}
For every prime \(p\), every finite extension of
\[
 U_p=\operatorname{Frac}(\mathbb Z_{(p)}^{sh})
       \bigl(p^{1/m}:(m,p)=1,\ \mu_{p^\infty}\bigr)
                                                        \tag{FF.35.1}
\]
satisfies \'etale reconstruction. Every subfield of
such an extension also satisfies \'etale reconstruction.
In particular, the failure of the functional detector
in Claim~\ref{clm:FFcyclotomicdetector} does not prevent
reconstruction of \(U_5\).
\end{claim}
\begin{proof}
We prove the actual-point-function condition of
\citet[Theorem~4.18(5)]{FF_CHW2024}. Let \(K/U_p\)
be finite, let \(X/K\) be regular, connected and
quasiprojective with \(x_0\in X(K)\), and let
\(\varphi:X_{\mathrm{\acute et}}\to
\mathbb G_{m,K,\mathrm{\acute et}}\) be pinned.
Use the notation \(b(x),b_0,\xi,j:X\to A\) of
Claim~\ref{clm:FFstrictreconstruct}; in particular
\(j(x_0)=0\). As there, \(X\) is geometrically
integral and its rational points are Zariski dense.
The density proof uses only Hensel's lemma and
infinite valuation balls, so applies to the
nondiscrete henselian field \(K\).
The zero-dimensional case is already covered by
Claim~\ref{clm:FF1}.

We may pass to finite extensions of \(K\) during
this proof: a unit realizing the original geometric
point function over such an extension descends by
Galois invariance and equality of functions on a
reduced scheme, exactly as at the beginning of
Claim~\ref{clm:FFstrictreconstruct}. This does not
presuppose reconstruction over the larger field.
All finite fields considered below remain finite
extensions of \(U_p\).

\emph{1. The part detected by actual point values.}
Choose a number field \(D_0\subset K\) containing
all the finite model data and \(b_0\), with
\(K=D_0U_p\), and set \(M_0=(D_0)_v\).
Then
\[
 G_K=H_0=\ker(\chi_p|_{P_{M_0}})
                  \triangleleft G_{M_0}.
                                                        \tag{FF.35.2}
\]
The geometric part of \(\xi-\kappa_p(b_0)\)
is an integral functional \(\lambda\) in the
space \(S\) of Claim~\ref{clm:FFcyclotomicsplitting}.
After the finite changes allowed there, write
\(\lambda=\lambda_0+\lambda_1\), with
\(\lambda_i\in S_i\). For every finite-extension
point, the actual-value identity becomes
\[
  (\lambda_0)_*\kappa_A(j(x))
                =\kappa_p(b(x)/b_0)
               \quad\text{over its field of definition}.
                                                        \tag{FF.35.3}
\]
Indeed, \(\lambda_1\) kills that entire rational
Kummer class by Claim~\ref{clm:FFcyclotomicsplitting}.

Choose finitely many point evaluations jointly
detecting \(S_0\), over one finite extension \(E/K\).
Include the points, their actual values, and a
generator of \(E/U_p\) in a number field \(D_1\).
Then \(E=D_1U_p\), \(M_1=(D_1)_v\), and
\(H_1=\ker(\chi_p|_{P_{M_1}})\) is normal in
\(G_{M_1}\). Joint detection survives further
finite restriction, by restriction--corestriction
with rational coefficients.
For \(g\in G_{M_1}\), the difference
\(g\lambda_0-\lambda_0\) lies in \(S_0\):
both \(S\) and its weight-zero summand are
\(G_{M_0}\)-stable. The restricted point and
value cocycles in (FF.35.3) are invariant under
the conjugation action of \(G_{M_1}\), since
their points and actual values are defined over
\(M_1\). The same cocycle computation as
(FF.26.4) therefore makes every chosen evaluation
of \(g\lambda_0-\lambda_0\) zero. Joint detection
gives
\[
                \lambda_0\in
             \operatorname{Hom}_{G_{M_1}}(V_pA,\mathbb Q_p(1)).
                                                        \tag{FF.35.4}
\]
No equivariance of \(\lambda_1\) under the full
local Galois group is needed.

There is a normalized rational global class
\(\xi_0\) over \(X_{M_1}\) with geometric part
\(\lambda_0\), and with value zero at \(x_0\).
Here is an explicit construction, so no integrality
of the projection is silently assumed. The point
\(x_0\) splits the arithmetic fundamental-group
sequence. On
\(\pi_1(X_{\overline M_1})\rtimes G_{M_1}\),
define the cocycle by applying \(\lambda_0\)
to the geometric abelian pro-\(p\) component
and assigning zero to the section \(G_{M_1}\).
Equivariance (FF.35.4) is exactly the cocycle
identity. Equivalently, pull the universal rational
Kummer class on \(A\) back by \(j\), and apply
\(\lambda_0\). Its evaluations satisfy (FF.35.3).
The same construction defines the compatible
class over the algebraic field \(E\).
We henceforth replace \(K\) by \(E\), as permitted.

\emph{2. Boundary residues are rational.}
The summand \(\lambda_1\) is zero on the torus
of the generalized Albanese, as proved in
Claim~\ref{clm:FFcyclotomicsplitting}. Thus the
boundary residues of \(\xi_0\) are the residues
of the original integral class \(\xi\), and
belong to \(\mathbb Z_p\). We now prove that
they also belong to \(\mathbb Q\).

On a punctured rational line over a finite model,
write the normalized class in the form (FF.26.1),
with residues \(a_i\in\mathbb Z_p\). Choose
a uniformizer \(\pi\) of the finite local model
in its dense number subfield. Evaluate at
\(r_i+\pi^N\) and \(r_i+\pi^{N+1}\), for \(N\)
larger than the valuations of all nonzero root
differences. Include these points and their
actual values in a further finite number field,
with completion \(M_2\). Over
\(H_2=\ker(\chi_p|_{P_{M_2}})\), equality
(FF.35.3) identifies each evaluation with the
Kummer class of its actual value. Both classes
come from \(G_{M_2}\); the restriction injection
of Claim~\ref{clm:FFcyclotomiccohomology} therefore
gives that equality already over \(G_{M_2}\).
Apply the ordinary normalized local valuation
to the difference of the two equalities. The
other puncture factors have constant valuation,
so
\[
        a_i v_p(\pi)
           =v_p\bigl(b(r_i+\pi^{N+1})\bigr)
                -v_p\bigl(b(r_i+\pi^N)\bigr)
                    \in\mathbb Q.
                                                        \tag{FF.35.5}
\]
Consequently \(a_i\in\mathbb Q\cap\mathbb Z_p
=\mathbb Z_{(p)}\).
This calculation uses the finite local field,
not a valuation map on Kummer cohomology over
\(K\); the latter would lose this information
because \(v(K^\times)=\mathbb Q\).

For a smooth curve, use the finite \'etale map
to a punctured line with a chosen boundary
point as its only pole, exactly as in Step~1
of Claim~\ref{clm:FFstrictreconstruct}.
Corestriction has actual point values given by
products of field norms, and the residue at
infinity is the chosen residue. The rational
class is defined over a finite local field,
and its norm values are still algebraic over
\(\mathbb Q\), so the preceding comparison
applies to that class without change.
Transverse curves through the smooth locus
of each boundary component, with intersection
multiplicity one, give the higher-dimensional
assertion. Therefore the full residue divisor
\(D\) of \(\xi_0\) has coefficients in
\(\mathbb Z_{(p)}\).

\emph{3. An algebraic unit up to powers and roots of unity.}
Apply the Picard 1-motive calculation (FF.26.5)
to the rational invariant geometric class
\(\lambda_0\). Choose a positive integer \(q\)
such that \(q\lambda_0\) is integral and
\(qD\in B_X\). Such an integer exists:
first clear the finitely many denominators
of \(D\); then use
\[
 (B_X\otimes\mathbb Q_p)\cap
       \Bigl(\bigoplus_i\mathbb QD_i\Bigr)
                         = B_X\otimes\mathbb Q
\]
and clear its rational lattice denominators;
finally multiply by a power of \(p\) to make
the entire geometric class integral.
An invariant integral lift with lattice
coordinate exactly \(qD\) gives
\([qD]\in\bigcap_n p^nP(M_1)\), by the
finite-level division-point argument in
Step~3 of Claim~\ref{clm:FFstrictreconstruct}.
That intersection is finite prime-to-\(p\)
torsion. Hence a positive multiple \(dD\)
is principal. The principal divisor descends
to the number-field model and has a function
\(h\in\mathcal O(X)^\times\), normalized by
\(h(x_0)=b_0^d\).

The zero-residue difference between
\(d\lambda_0\) and the geometric Kummer class
of \(h\) lies in \((V_pP)^{G_{M_1}}=0\).
Normalization at \(x_0\) therefore gives the
rational equality
\[
                   d\xi_0=\kappa_p(h/b_0^d).
                                                        \tag{FF.35.6}
\]
Evaluate and use (FF.35.3). The rational
Kummer kernel in every finite extension of
\(K\) is \(\mu_\infty\), by
Claim~\ref{clm:FFcyclotomiccohomology}. Thus
\[
                b(x)^d=\zeta_xh(x),
                         \qquad \zeta_x\in\mu_\infty\subset K.
                                                        \tag{FF.35.7}
\]
We have proved this pointwise relation; we
do not identify the original class \(\xi\)
with the Kummer class of \(h\). Its invisible
summand may be nonzero, as Claim~\ref{clm:FFcyclotomicdetector}
demonstrates.

\emph{4. Closed points of arbitrarily large degree.}
Every finite extension \(L/K\) has infinite
absolute Galois group. Indeed,
Claim~\ref{clm:FFcyclotomictest} identifies it
with a nonabelian subgroup of the free wild
inertia group. An infinite profinite group
has open subgroups of arbitrarily large
index; otherwise its finite quotients would
have bounded order, forcing the group to
be finite. Finite separable extensions of
\(L\) therefore have unbounded degrees.
By the primitive-element theorem their
generators give closed points on
\(\mathbb A^1_L\) of unbounded degrees.
A finite nonconstant map from the proper
normalization of a geometrically integral
curve to \(\mathbb P^1_L\) then gives
closed points of unbounded degrees in
every nonempty open of that curve: choose
base points of larger degree than every
point in the finite image of its omitted
locus. Points above them have at least
the same degree. This replaces the
Eisenstein-polynomial construction, which
would not be valid for the divisible
value group of \(K\).

\emph{5. Recovering the actual point function.}
The remaining algebraic steps of
Claim~\ref{clm:FFstrictreconstruct}, Steps~4--7,
now apply, with \(\lambda_0\) in place of
\(\lambda\) and \(\mu_\infty\) in place of
\(\mu_{p'}\). We spell out why their exact
inputs have all been retained.

If \(\lambda_0\ne0\), the unit \(h\) in
(FF.35.6) is nonconstant. Relation (FF.35.7)
and the existence of large-degree points
imply that the actual point function is
nonconstant. On the nonempty inverse-image
open of \(T-1\ne0\), apply Steps~1--3 to
that target function. This gives
\((b(x)-1)^e=\eta_xg(x)\), with
\(\eta_x\in\mu_\infty\subset K\).
Use a smooth geometrically integral curve
on which restriction is injective on the
finite span of \(h^ig^j\), as supplied by
\citet[Theorem~1.8 and Corollary~1.9]{FF_CharlesPoonen2016}.
For a closed point of degree greater than
\(e\deg\operatorname{div}_\infty h+
d\deg\operatorname{div}_\infty g\), the
resultant
\[
        \operatorname{Res}_Z
          (Z^d-\zeta_xh,\ (Z-1)^e-\eta_xg)
\]
vanishes identically, by its uniform pole
bound. Its coefficients lie in \(K\),
since all roots of unity do. Restriction
injectivity makes it vanish in \(K(X)\).
Its common root \(z\) is rational: a
conjugate satisfies \(z'=\alpha z\) and
\(z'-1=\beta(z-1)\), with \(\alpha,\beta\)
constant roots of unity, and nonconstancy
forces \(\alpha=\beta=1\). It is a unit
because \(z^d=\zeta_xh\).
Equation (FF.35.6) then gives
\(\lambda_0=\kappa_p(z)\) geometrically
with rational coefficients. Rational
coefficients suffice for this division
by \(d\), and a rationally zero unit
Kummer class still has zero divisor.
Normalize \(z(x_0)=b_0\). Evaluation at
\(x_0\) kills the constant cohomology
difference, so (FF.35.3) becomes
\(\kappa_p(b(x))=\kappa_p(z(x))\)
at every finite-extension point. Hence
\(b(x)/z(x)\in\mu_\infty\subset K\).

If \(\lambda_0=0\), (FF.35.3) puts all
values \(b(x)/b_0\) in \(\mu_\infty\subset K\).
Choose \(\beta\) in a proper finite
extension of \(K\), and apply the same
construction to \(T-\beta\). Such an
extension exists by Step~4. If its
weight-zero class is nonzero, the preceding
resultant gives a nonconstant unit, whereas
all values of that unit lie in the fixed
finite base extension containing \(\beta\):
indeed \(b(x)\in K\), and its root-of-unity
ratio also belongs to \(K\). This holds for
all geometric points, including points defined
over arbitrarily larger finite fields: every
\(b(x)-\beta\) still lies in \(K(\beta)\).
The
large-degree-point argument in Step~5
of Claim~\ref{clm:FFstrictreconstruct}
excludes this. If its weight-zero class
is zero, comparing the two root-of-unity
ratios and using the linear independence
of \(1,\beta\) over \(K\) makes
\(b(x)=b_0\). Thus the zero case has a
constant actual point function.

Finally, for a nonconstant weight-zero
class, (FF.35.3) and its realizing unit
give the same constant-scalar ambiguity
as in Step~6 of
Claim~\ref{clm:FFstrictreconstruct}.
Apply the construction also to \(T-1\),
and use a closed point of degree larger
than the relevant pole bound and the
degree of the nonconstant unit as a map
to \(\mathbb P^1\). Linear independence
at that point forces the two units and
the actual values to have one fixed
scalar relation. Applying \(T-a\) deals
with the remaining base-field-value
points. Each nonconstant point function
remains nonconstant on its inverse-image
open: otherwise the inverse image of a
closed target point would contain a dense
open of the integral source and hence
be the whole source. In higher dimension,
the same argument is made on a smooth
geometrically integral curve through the
base point and the point being tested,
using the cited Charles--Poonen theorem.
These are precisely the algebraic
scalar-removal steps already proved there;
they require only (FF.35.3), all roots of
unity in the base, and Step~4.

The weight decomposition used throughout
is compatible with restriction to opens,
finite scalar extensions, and pullback
along the curves: it is the eigenspace
decomposition of the normalized inertia
Lie generator. In particular a geometric
Kummer functional of an algebraic unit
belongs to weight zero, and comparisons
of such units in this last step preserve
the normalized class \(\lambda_0\).
We have therefore obtained a unit whose
actual geometric point function is exactly
that of \(\varphi\). Finite descent returns
it to the original \(K\). The point-function
criterion cited at the start proves
\(\mathcal R(K)\).
Every subfield \(F\subset K\) has
\(K/F\) algebraic, since \(K/\mathbb Q\)
is algebraic. Algebraic descent, Claim~\ref{clm:FFdescent},
proves \(\mathcal R(F)\).
\end{proof}

\subsection{Maximal abelian extensions of strict henselization fields}

\begin{claim}{FF.36}{Proved}\label{clm:FFabelianHT}
Let \(M/\mathbb Q_p\) be finite, and let \(S\)
be a Hodge--Tate representation of \(G_M\) whose
inertia image is abelian. The inertia representation
on \(S\) is semisimple. In particular there is a
canonical decomposition
\[
 S=S^{I_M}\oplus
       \sum_{g\in I_M}(g-1)S.                         \tag{FF.36.1}
\]
Its projection onto the first summand commutes
with \(G_M\). For every inertia-stable subspace
\(N\), that projection sends \(N\) to
\(N\cap S^{I_M}\), and taking inertia invariants
is exact on \(0\to N\to S\to S/N\to0\).
All assertions hold after any finite extension
of \(M\).
\end{claim}
\begin{proof}
Write \(\mathfrak g\) for the Lie algebra of
the compact inertia image. It is abelian.
Let \(B\subset\operatorname{End}_{\mathbb Q_p}(S)\)
be the unital commutative algebra generated by
\(\mathfrak g\). It is a finite-dimensional
commutative algebra over the characteristic-zero
field \(\mathbb Q_p\).

We recall the elementary algebra fact used
here. Such an algebra has a unique maximal
\'etale subalgebra \(B_{\rm ss}\), and
\[
           B=B_{\rm ss}\oplus\operatorname{rad}(B)
                         \quad\text{as vector spaces}.
                                                        \tag{FF.36.2}
\]
One proof uses the additive and multiplicative
Jordan decompositions of commuting matrices:
their semisimple parts belong to the algebra,
are polynomials in the original matrices, and
add and multiply as the original matrices do
modulo the nilpotent radical. These parts
give the unique section of the finite separable
algebra \(B/\operatorname{rad}(B)\). Equivalently,
in each local Artin factor, its separable residue
field has a unique lift through the nilpotent
ideal. After extension to \(\mathbb C_p\),
each factor has the form
\(\mathbb C_p\oplus\mathfrak m\), with
\(\mathfrak m\) nilpotent. An element acting
semisimply on the faithful module has zero
\(\mathfrak m\)-component. Thus the semisimple
elements of \(B\otimes\mathbb C_p\) are exactly
\(B_{\rm ss}\otimes\mathbb C_p\).

The Sen operator \(\Theta\) of \(S\) is
semisimple, because \(S\) is Hodge--Tate
\citep[Section~2.7]{FF_Fontaine2004}.
By \citet[Theorem~3.1, \(d=0\)]{FF_Ohkubo2014},
\(\Theta\in\mathfrak g\otimes\mathbb C_p\),
and \(\mathfrak g\) is the smallest
\(\mathbb Q_p\)-subspace of the endomorphism
space with this property. The preceding
algebra fact puts \(\Theta\) in
\(B_{\rm ss}\otimes\mathbb C_p\). Minimality
therefore gives \(\mathfrak g\subset B_{\rm ss}\).
In particular every element of \(\mathfrak g\)
acts semisimply, and these elements commute.

On a sufficiently small open subgroup of
the compact inertia image, matrix logarithm
has values in \(\mathfrak g\), and matrix
exponential is its inverse. Hence every
element of this open subgroup is semisimple.
An arbitrary inertia matrix has a positive
integer power in this subgroup. A matrix
over a characteristic-zero field whose
positive power is invertible and semisimple
is itself semisimple: if that power has
distinct eigenvalues \(\alpha_j\), the
matrix is annihilated by the square-free
polynomial \(\prod_j(T^n-\alpha_j)\).
All inertia matrices commute, so they are
simultaneously diagonalizable over a finite
extension of \(\mathbb Q_p\). Their
representation is semisimple over
\(\mathbb Q_p\) as well.

In this simultaneous decomposition,
\(S^{I_M}\) is the trivial-character part,
and the sum in (FF.36.1) is exactly the
sum of the nontrivial-character parts.
This description is independent of a
coefficient extension. It also proves
the assertion about every stable \(N\)
and exactness of invariants. Conjugation
by \(G_M\) preserves inertia and its
trivial-character part, as well as the
explicit second summand in (FF.36.1).
The projection therefore commutes with
\(G_M\). Restricting to any open local
Galois subgroup retains both the
Hodge--Tate and abelian-image hypotheses.
\end{proof}

\begin{claim}{FF.37}{Proved}\label{clm:FFmaxabelianreconstruct}
Let \(K_p=\operatorname{Frac}(\mathbb Z_{(p)}^{sh})\)
with the chosen embedding in \(\overline{\mathbb Q}\),
and let \(K_p^{\rm ab}\) be its maximal
abelian extension in that algebraic closure.
Every finite extension of \(K_p^{\rm ab}\),
and every subfield of such an extension,
satisfies \'etale reconstruction.
The extension \(K_p^{\rm ab}/U_p\) is infinite,
so this is a proper enlargement of
Claim~\ref{clm:FFcyclotomicreconstruct} at the
level of the specified henselian fields.
\end{claim}
\begin{proof}
Put \(I=I_{\mathbb Q_p}=G_{K_p}\), and
let \(J=\overline{[I,I]}\) be its closed
commutator subgroup. Thus
\(G_{K_p^{\rm ab}}=J\).
The tame quotient of inertia is abelian,
so \(J\subset P_{\mathbb Q_p}\).
It contains the commutator subgroup of
the free pro-\(p\) group
\(P_{\mathbb Q_p}\), and is infinite:
two free generators have a nonidentity
commutator, of infinite order by
torsion-freeness. Every open subgroup
of \(J\) is consequently infinite.
This uses the same verified free wild
inertia theorem as Claim~\ref{clm:FFcyclotomictest}.

Both \(T_p/K_p\) and the cyclotomic
extension \(K_p(\mu_{p^\infty})/K_p\)
are abelian. Hence \(U_p\subset K_p^{\rm ab}\).
In particular these fields contain
\(\mu_\infty\), have value group
\(\mathbb Q\), and have residue field
\(\overline{\mathbb F}_p\).

\emph{1. The enlargement is infinite.}
Set \(L=K_p(\zeta_p)\), including
\(L=K_2\) when \(p=2\), and let
\(\pi=\zeta_p-1\). Then \(\pi\) is
a uniformizer of \(L\), and
\(\Delta=\operatorname{Gal}(L/K_p)\)
has order \(p-1\). For the principal
unit filtration \(U^i=1+\pi^i\mathcal O_L\),
\[
                    U^1/U^2\simeq
                 \overline{\mathbb F}_p
                \quad\text{as }\mathbb F_p\text{-spaces}.
                                                        \tag{FF.37.1}
\]
The \(\Delta\)-action on this quotient
is the mod-\(p\) cyclotomic character
\(\bar\chi_p\): its residue action is
trivial, and
\(\sigma(\pi)\equiv\bar\chi_p(\sigma)\pi
\pmod{\pi^2}\).
If \(u^p\in U^1\), then \(u\) is a
unit of residue one, because the residue
field has characteristic \(p\). Writing
\(u=1+\pi c\), the binomial expansion
and \(v_L(p)=p-1\) show that
\(u^p\in U^2\). Therefore (FF.37.1)
is a quotient of a subspace of
\(L^\times/(L^\times)^p\).
Since \(|\Delta|\) is prime to \(p\),
taking its character parts is exact.
The \(\bar\chi_p\)-part of this Kummer
space is infinite-dimensional over
\(\mathbb F_p\).

Kummer theory and the twist by
\(\mu_p^{-1}\) identify that part
with \(H^1(G_L,\mathbb F_p)^\Delta\).
Inflation--restriction for \(\Delta\)
identifies the latter with
\(H^1(I,\mathbb F_p)\), because
the positive-degree cohomology of this
prime-to-\(p\) finite group is zero.
Thus \(I^{\rm ab}\) has an
infinite-dimensional elementary abelian
pro-\(p\) quotient. By contrast
\(\operatorname{Gal}(U_p/K_p)\)
is a closed subgroup of a product
of a pro-prime-to-\(p\) procyclic
group and a subgroup of
\(\mathbb Z_p^\times\); its pro-\(p\)
Sylow embeds in the second factor.
Its elementary abelian pro-\(p\)
quotient is finite-dimensional.
A finite kernel could increase that
dimension by only a finite amount.
Hence \(K_p^{\rm ab}/U_p\) cannot be finite.

\emph{2. Exact cohomological inputs for all finite extensions.}
Let \(K/K_p^{\rm ab}\) be finite.
It has the form \(L_0K_p^{\rm ab}\)
for a finite extension \(L_0/K_p\).
Every specified element of \(K\)
can be included in a larger such
\(L_0\). The extension \(K/L_0\)
is Galois abelian and has infinite
cyclotomic image. The abelian cocycle
calculation (FF.33.2), together with
the exact Kummer kernel over \(L_0\),
therefore proves
\[
 \ker\bigl(K^\times\to H^1(G_K,\mathbb Z_p(1))\bigr)
 =\ker\bigl(K^\times\to H^1(G_K,\mathbb Q_p(1))\bigr)
 =\mu_\infty.
                                                        \tag{FF.37.2}
\]
Compatible roots may be chosen because
all roots of unity are in \(K\).
This is the proof of
Claim~\ref{clm:FFcyclotomiccohomology}
with another abelian extension in
place of \(U_p\); all its hypotheses
have just been checked.

Choose a number field \(D_0\subset K\)
containing any prescribed finite data
and generators of \(K/K_p^{\rm ab}\),
and put \(M_0=(D_0)_v\). Inside the
fixed local Galois group the absolute
group of \(K\) is
\[
            H_0=J\cap G_{M_0}
                         \triangleleft G_{M_0}.
                                                        \tag{FF.37.3}
\]
The quotient \(I_{M_0}/H_0\) is
abelian, as a subgroup of \(I/J\),
and its cyclotomic image is infinite.
On this abelian quotient, (FF.33.2)
gives \(H^1(I_{M_0}/H_0,\mathbb Q_p(1))=0\).
Its coefficient invariants are zero.
Inflation--restriction for inertia
in \(G_{M_0}/H_0\) then also gives
\(H^1(G_{M_0}/H_0,\mathbb Q_p(1))=0\).
Consequently both restriction maps
\[
 H^1(I_{M_0},\mathbb Q_p(1))\longrightarrow H^1(H_0,\mathbb Q_p(1)),
 \qquad
 H^1(G_{M_0},\mathbb Q_p(1))\longrightarrow H^1(H_0,\mathbb Q_p(1))
                                                        \tag{FF.37.4}
\]
are injective. Every finite extension
\(E/K\) equals \(D_1K_p^{\rm ab}\)
for a finite number field \(D_1/D_0\),
and has \(H_1=J\cap G_{M_1}\) with
\(M_1=(D_1)_v\). The same proof gives
(FF.37.2)--(FF.37.4) at every such stage.

\emph{3. The invisible subspace.}
Use the regular pointed source \(X\),
its generating generalized Albanese
\(j:X\to A\), the actual point function
\(b\), and its geometric functional
\(\lambda\) as in
Claim~\ref{clm:FFcyclotomicreconstruct}.
Include the data in \(D_0\), and set
\[
                   S=\operatorname{Hom}_{H_0}
                                (V_pA,\mathbb Q_p(1)).
                                                        \tag{FF.37.5}
\]
It is a \(G_{M_0}\)-stable Hodge--Tate
space, by the exact comparison and
subobject results verified in
Claim~\ref{clm:FFcyclotomicsplitting}.
Its inertia action factors through
the abelian quotient in (FF.37.3).
Claim~\ref{clm:FFabelianHT} thus applies.

Let \(N\subset S\) consist of those
functionals whose pushforward of
\(\kappa_A(a)\) is zero for every
\(a\in A(L)\) over every finite
extension \(L/K\). This is a linear
subspace. It is \(G_{M_0}\)-stable:
conjugation sends such an \(L\) to
another finite extension of \(K\),
and sends the point and its Kummer
cocycle to their conjugates. This
uses normality of \(H_0\), and
includes the coefficient action.

The same \(N\) is the common kernel
of all evaluations at \(j(x)\).
To see the nontrivial inclusion, recall
that finitely many signed sums of
\(j(X_{\bar K})\) dominate \(A_{\bar K}\).
Every point of \(A(\bar K)\) is a
difference of points in this dominant
image, since two nonempty opens related
by a translation intersect. Thus it
is a finite signed sum of \(j(x_i)\),
after a finite field extension.
A nonzero rational Kummer evaluation
cannot become zero under a finite
restriction, by restriction--corestriction.
Consequently vanishing on all \(j(x)\)
implies vanishing on every \(a\).
Finite-dimensionality of \(S/N\)
now supplies finitely many such
evaluations with common kernel
exactly \(N\), over one finite
extension \(E/K\).

For every finite local model \(M_1\)
obtained as in Step~2, we have
\[
                        N\cap S^{I_{M_1}}=0.
                                                        \tag{FF.37.6}
\]
Indeed, a nonzero inertia-invariant
functional is detected over a finite
extension of the corresponding strict
henselization by
Claim~\ref{clm:FFpdetector}. Its
nonzero class survives adjoining
\(K_p^{\rm ab}\) by (FF.37.4),
so it cannot belong to \(N\).

\emph{4. The canonical invariant lift.}
Include the finitely many detecting
points and their actual values in
\(D_1\subset E\), so that
\(E=D_1K_p^{\rm ab}\) and
\(H_1=J\cap G_{M_1}\) is normal
in \(G_{M_1}\). The actual-value
cohomology identity, followed by
conjugation, gives
\[
                  g\lambda-\lambda\in N
                         \quad(g\in G_{M_1}).
                                                        \tag{FF.37.7}
\]
Here the selected point and value
cocycles are defined over \(M_1\),
and their restricted classes are
conjugation-invariant. Their common
evaluation kernel is exactly \(N\).
No descent of the topos map is
being assumed.

Take the canonical projection
\(\Pi:S\to S^{I_{M_1}}\) from
Claim~\ref{clm:FFabelianHT}, and put
\(\lambda_0=\Pi\lambda\).
Equation (FF.37.7) makes the image
of \(\lambda\) in \(S/N\)
inertia-invariant. Exactness of
invariants and (FF.37.6) imply
\[
               \lambda-\lambda_0\in N,
           \qquad \lambda_0\in
             \operatorname{Hom}_{G_{M_1}}(V_pA,\mathbb Q_p(1)).
                                                        \tag{FF.37.8}
\]
For the second assertion, \(\Pi\)
commutes with \(G_{M_1}\), and
\(\Pi N=N\cap S^{I_{M_1}}=0\).
Thus applying \(\Pi\) to
(FF.37.7) kills every difference.
This is the unique inertia-invariant
lift of the class modulo \(N\).
It is not asserted to have been
\(I_{M_0}\)-invariant before the
finite actual-value descent.

After splitting the torus of \(A\)
over a further finite model, restriction
of a functional to that torus has
inertia-trivial target
\(\operatorname{Hom}(V_pT,\mathbb Q_p(1))\).
Every element of \(N\) has zero
restriction: an inertia-equivariant
map to a trivial representation
factors through \(\Pi\), and
\(\Pi N=0\). The original integral
boundary residues of \(\lambda\)
are therefore unchanged by replacement
with \(\lambda_0\).

\emph{5. Reconstruction from the retained point values.}
Construct the normalized rational
class \(\xi_0\) from \(\lambda_0\)
using the section at \(x_0\),
exactly as in Step~1 of
Claim~\ref{clm:FFcyclotomicreconstruct}.
Since \(\lambda-\lambda_0\in N\),
its evaluation at every finite-extension
point is \(\kappa_p(b(x)/b_0)\).
Its residues are the integral residues
of the original class, by Step~4.
All the subsequent hypotheses in
that reconstruction proof are now
available, as follows.

First, (FF.37.4) lifts an equality
of actual evaluations from \(H_2\)
to the finite local model \(G_{M_2}\).
The two nearby punctured-line values
therefore force each boundary residue
into \(\mathbb Q\cap\mathbb Z_p\),
using the ordinary finite-local-field
valuation. Norm transfer and transverse
curves give the same assertion in
every dimension. Second, the invariant
rational Picard 1-motive lift can
be multiplied to have integral Tate
coordinates and an integral lattice
coordinate. The compact local Picard
group then makes a further multiple
of that boundary divisor principal,
and its invariant zero-residue Tate
part is zero. Thus a unit \(h\)
and integer \(d>0\) satisfy
\[
                    b(x)^d=\zeta_xh(x),
                         \qquad\zeta_x\in\mu_\infty\subset K,
                                                        \tag{FF.37.9}
\]
by (FF.37.2). This is an equality
of point values, not an identification
of the original integral topos class
with a unit class.

Third, every finite extension of
\(K\) has infinite absolute Galois
group, since its group is open in
\(J\). Hence the primitive-element
and finite-map argument in Step~4
of Claim~\ref{clm:FFcyclotomicreconstruct}
supplies closed points of arbitrarily
large degree on every nonempty
curve open. The \(T,T-1\) resultant
applied to (FF.37.9) therefore makes
the retained nonzero functional the
geometric Kummer functional of a unit.
If it is zero, every actual value
lies in \(b_0\mu_\infty\subset K\),
and the \(T-\beta\) argument forces
the point function to be constant.
The \(T-a\) comparisons remove the
remaining root-of-unity scalars.

These comparisons are compatible
with the canonical projections just
used. To verify this without making
\(N\) explicit, a pullback on
Albanese Tate modules carries an
inertia-invariant realizing functional
to an inertia-invariant one, and
carries an invisible functional to
an invisible one: point Kummer
classes pull back by actual points.
Uniqueness of the invariant lift
in (FF.37.8) gives the required
compatibility, after a common finite
model extension. The same applies
to restrictions to nonempty opens
and to the curves used for testing.
Thus every step of the actual
point-function and scalar-removal
proof is preserved.

It follows that the pinned map's
geometric point function is realized
by a unit over a finite extension.
As in Claim~\ref{clm:FFstrictreconstruct},
equality on geometric points makes
that unit Galois-invariant, so it
descends to the original \(K\).
The point-function criterion
\citep[Theorem~4.18(5)]{FF_CHW2024}
proves \(\mathcal R(K)\).
Since \(K\) is algebraic over
\(\mathbb Q\), it is algebraic
over each of its subfields.
Claim~\ref{clm:FFdescent} proves
the final subfield assertion.
\end{proof}


\begin{claim}{FF.38}{Proved}\label{clm:FFnilpotentinputs}
The following two assertions hold.
\begin{enumerate}
\item Let \(\Gamma\) be a profinite group
nilpotent of some fixed finite class, and let
\(\chi:\Gamma\to\mathbb Q_p^\times\) be a
nontrivial continuous character. Then
\[
                 H^1(\Gamma,\mathbb Q_p(\chi))=0.
                                                        \tag{FF.38.1}
\]
\item Let \(S\) be a Hodge--Tate representation
of the Galois group of a finite extension
\(M/\mathbb Q_p\). If its inertia Lie algebra
is nilpotent, that Lie algebra is abelian.
After a finite extension of \(M\), the
inertia representation is semisimple and
has the canonical invariant projection
of Claim~\ref{clm:FFabelianHT}.
\end{enumerate}
\end{claim}
\begin{proof}
For the first assertion, induct on the
nilpotency class and put \(Z=Z(\Gamma)\), a closed central subgroup. If
\(\chi|_Z\ne1\), choose \(z\in Z\)
with \(\chi(z)\ne1\). For every cocycle
\(c\), centrality gives
\[
             (\chi(z)-1)c(g)=(\chi(g)-1)c(z).
\]
Thus \(c\) is a coboundary, proving
(FF.38.1). If \(\chi|_Z=1\), it descends
to a nontrivial character of \(\Gamma/Z\).
Inflation--restriction gives
\[
 0\longrightarrow H^1(\Gamma/Z,\mathbb Q_p(\chi))
 \longrightarrow H^1(\Gamma,\mathbb Q_p(\chi))
 \longrightarrow H^1(Z,\mathbb Q_p(\chi))^{\Gamma/Z}.
\]
The first group is zero by induction,
since \(\Gamma/Z\) has smaller
nilpotency class. The last group is
zero too: it is the invariant part
of \(\operatorname{Hom}_{\rm cont}(Z,\mathbb Q_p)
\otimes\mathbb Q_p(\chi)\), on which
\(\Gamma\) acts only through the
nontrivial scalar \(\chi\), because
\(Z\) is central. This proves the
first assertion, including the abelian
base case.

For the second, write \(\mathfrak g\)
for the inertia Lie algebra and
\(\Theta\in\mathfrak g\otimes\mathbb C_p\)
for the Sen operator. It is semisimple
by Hodge--Tate comparison
\citep[Section~2.7]{FF_Fontaine2004}.
Hence \(\operatorname{ad}\Theta\)
is semisimple on
\(\operatorname{End}(S)\otimes\mathbb C_p\),
and on its invariant subspace
\(\mathfrak g\otimes\mathbb C_p\).
It is nilpotent on that subspace
because \(\mathfrak g\) is nilpotent.
It must therefore be zero. Thus
\(\Theta\in Z(\mathfrak g)\otimes\mathbb C_p\).
The minimal-subspace assertion of
\citet[Theorem~3.1, \(d=0\)]{FF_Ohkubo2014}
implies
\(\mathfrak g\subset Z(\mathfrak g)\).
So \(\mathfrak g\) is abelian.

Matrix exponential identifies a
sufficiently small open subgroup
of the compact inertia image with
a lattice in this abelian Lie algebra.
That open subgroup is abelian.
A finite extension of \(M\) makes
its inertia image lie in this subgroup:
choose an open subgroup of \(G_M\)
whose intersection with inertia is
inside the given open inertia subgroup.
Claim~\ref{clm:FFabelianHT} applies
to the restricted Hodge--Tate
representation and proves the
remaining assertions. This concerns
a finite extension, not an infinite
passage that might destroy the
finite local-field hypotheses.
\end{proof}

\begin{claim}{FF.39}{Proved}\label{clm:FFnilpotentreconstruct}
Let \(I=G_{K_p}=I_{\mathbb Q_p}\),
and define its closed lower central
series by
\(\gamma_1(I)=I\) and
\(\gamma_{r+1}(I)=\overline{[I,\gamma_r(I)]}\).
For each integer \(c\ge1\), put
\[
              K_{p,c}=
                \overline{\mathbb Q}^{\,\gamma_{c+1}(I)}.
                                                        \tag{FF.39.1}
\]
Every finite extension of \(K_{p,c}\),
and every subfield of such an extension,
satisfies \'etale reconstruction.
Here \(K_{p,1}=K_p^{\rm ab}\), and
every extension \(K_{p,c+1}/K_{p,c}\)
is infinite. The nilpotency class
\(c\) is fixed and finite in the
assertion; no conclusion about the
union of these fields or a general
solvable closure is included.
\end{claim}
\begin{proof}
Write \(J_c=\gamma_{c+1}(I)\).
It is characteristic in inertia,
hence normal in \(G_{\mathbb Q_p}\).
The quotient \(I/J_c\) is nilpotent
of class at most \(c\).
The field in (FF.39.1) is exactly
the compositum of all finite Galois
extensions of \(K_p\) with Galois
group of nilpotency class at most
\(c\): the lower-central identities
pass to finite quotients, and the
intersection of the kernels of the
finite quotients of \(I/J_c\) is
trivial.

We first verify that these stages
are proper and form a strictly
increasing infinite hierarchy.
The free wild inertia group \(P\)
has \(H^2(P,\mathbb F_p)=0\).
The quotient \(I/P\) is pro-prime-to-\(p\),
whose positive-degree cohomology
on \(p\)-primary discrete modules
vanishes by averaging on finite
quotients. Hochschild--Serre gives
\[
                   H^2(I,\mathbb F_p)=0.
                                                        \tag{FF.39.2}
\]
Also \(H^1(I,\mathbb F_p)\) is
infinite-dimensional by the principal-unit
calculation in Step~1 of
Claim~\ref{clm:FFmaxabelianreconstruct}.

Every finite \(p\)-group \(E\)
is a quotient of \(I\). Here are
the embedding-problem details.
Its Frattini quotient \(E/\Phi(E)\)
is a finite-dimensional \(\mathbb F_p\)-space,
so the preceding infinite-dimensional
\(H^1\) supplies a surjection from
\(I\) to that quotient. Choose a
central chief series inside
\(\Phi(E)\): each successive
factor has order \(p\) and is
central in the corresponding quotient
of \(E\). Such a series exists
because a nontrivial normal subgroup
of a finite \(p\)-group meets
its center nontrivially, and one
continues in the quotient. Each
successive central extension by
\(\mathbb F_p\) has zero pullback
obstruction by (FF.39.2), so the
map from \(I\) lifts through all
of them. Its final image is all
of \(E\): a proper subgroup lies
in a maximal subgroup, and every
maximal subgroup contains
\(\Phi(E)\), contrary to surjectivity
on the Frattini quotient.

For each \(n\), apply this to
\(E=\operatorname{UT}_{c+2}(\mathbb F_p)^n\).
Each unitriangular factor has
nilpotency class \(c+1\), and
its \((c+1)\)-st lower-central
subgroup has order \(p\).
The nested commutator of the
successive elementary transvections
is \(1+E_{1,c+2}\), proving this
also in small characteristic.
It follows that
\(\gamma_{c+1}(I)/\gamma_{c+2}(I)\)
has quotients \((\mathbb Z/p)^n\)
for arbitrarily large \(n\).
It is infinite, proving the
asserted infinite stage extensions.

We will also need that \(J_c\)
and all its open subgroups are
infinite. Independently of the
preceding quotient argument,
\(J_c\) contains \(\gamma_{c+1}(P)\).
Map \(c+1\) free generators of
\(P\) to the successive elementary
transvections just used. Their
nested commutator is nonidentity,
so \(\gamma_{c+1}(P)\ne1\).
Torsion-freeness of \(P\) makes
it infinite. Every open subgroup
of \(J_c\) is infinite as well.

Now fix a finite extension
\(K/K_{p,c}\). It contains
\(K_p^{\rm ab}\), hence
\(U_p\) and \(\mu_\infty\).
Write \(K=L_0K_{p,c}\), with
\(L_0/K_p\) finite, enlarging
\(L_0\) to contain any specified
finite data. The Galois group
\(\operatorname{Gal}(K/L_0)\)
is a closed subgroup of the
nilpotent group \(I/J_c\),
so is nilpotent of class at most
\(c\), with infinite cyclotomic
image. Claim~\ref{clm:FFnilpotentinputs}(1)
replaces the abelian cocycle identity
in the proof of (FF.37.2).
It proves the same exact integral
and rational Kummer kernels
\(\mu_\infty\), because an
infinitely \(p\)-divisible element
has its compatible Kummer cocycle
factoring through this nilpotent
quotient. Rational triviality there
implies that a \(p\)-power of
the element belongs to the known
integral kernel over \(L_0\),
and hence the element is a
root of unity.

Choose a finite number model
\(D_0\subset K\), put
\(M_0=(D_0)_v\), and identify
\[
        G_K=H_0=J_c\cap G_{M_0}
                         \triangleleft G_{M_0}.
                                                        \tag{FF.39.3}
\]
The quotient \(I_{M_0}/H_0\)
is nilpotent of class at most
\(c\) and has nontrivial
cyclotomic character. Assertion~(1)
again gives zero \(H^1\) of this
quotient with coefficients
\(\mathbb Q_p(1)\), while its
coefficient invariants are zero.
The two inflation--restriction
arguments in (FF.37.4) prove
injectivity from both inertia
and the full local Galois group
to \(H_0\). The same statements
hold at every finite field and
finite local-model extension.

For a pointed generalized Albanese,
form the finite-dimensional space
\(S=\operatorname{Hom}_{H_0}(V_pA,\mathbb Q_p(1))\).
It is Hodge--Tate and its inertia
image is nilpotent of class at
most \(c\); its Lie algebra
is therefore nilpotent.
By Claim~\ref{clm:FFnilpotentinputs}(2),
a finite local extension makes
this inertia representation abelian
and semisimple. Realize that
extension on a finite algebraic
model, possibly passing to a
finite extension of \(K\).
Retain the original space \(S\)
under this restriction; it is
unnecessary to replace it by
a larger Hom space. The finite
extension is permissible in the
point-function descent argument.

The all-point-invisible subspace
\(N\subset S\) of Step~3 in
Claim~\ref{clm:FFmaxabelianreconstruct}
has the same definition and the
same normality and generation
properties. It is unchanged by
testing only finite extensions
of the enlarged base: any finite
extension of the original base
can be combined with that base,
and rational restriction--corestriction
preserves a nonzero evaluation.
The inertia detector of
Claim~\ref{clm:FFpdetector}, followed
by the just-proved restriction
injection, again gives
\(N\cap S^{I_{M_1}}=0\) at
every further finite local stage.
Actual point/value descent and
the canonical semisimple projection
therefore give the unique invariant
lift \(\lambda_0\) modulo \(N\),
exactly as in (FF.37.7)--(FF.37.8).
The projection remains compatible
with torus restrictions, open
restrictions and curve pullback;
in particular it retains the
original boundary residues.

All remaining inputs of the
reconstruction proof in Step~5
of Claim~\ref{clm:FFmaxabelianreconstruct}
have now been established: the
exact Kummer kernel, finite-local
restriction injection, unchanged
integral boundary residues, and
the infinite groups of every
finite extension needed for
large-degree points. Thus the
normalized rational class yields
a power of an algebraic unit,
the \(T,T-1\) resultant removes
that power, and \(T-\beta,T-a\)
remove the constant and root-of-unity
ambiguities in the actual point
function. Finite descent returns
the resulting unit to the
original \(K\), and
\citep[Theorem~4.18(5)]{FF_CHW2024}
gives \(\mathcal R(K)\).
Algebraic descent proves the
subfield assertion, since all
these fields are algebraic
over \(\mathbb Q\).
The proof uses a fixed bound
\(c\) in the cohomological
central-series induction and
the inertia Lie argument;
neither argument has been
applied to an unbounded union.
\end{proof}


\begin{claim}{FF.40}{Proved}\label{clm:FFmetabeliantest}
Let \(I=G_{K_p}\), let \(I'=\overline{[I,I]}\),
let \(I''=\overline{[I',I']}\), and put
\[
                  B_p=\overline{\mathbb Q}^{\,I''}.
                                                        \tag{FF.40.1}
\]
This maximal metabelian extension of \(K_p\)
is a proper countable algebraic subfield of
\(\overline{\mathbb Q}\), with henselian
value group \(\mathbb Q\), residue field
\(\overline{\mathbb F}_p\), and
\(\mu_\infty\subset B_p\).
For every finite extension \(E/B_p\),
\[
             C_{\operatorname{Aut}(\overline{\mathbb Q})}(G_E)=1.
                                                        \tag{FF.40.2}
\]
It belongs to none of the sufficient field
families established in this section through
Claim~\ref{clm:FFnilpotentreconstruct}, including
the subfields of finite extensions of a
fixed bounded-nilpotency stage at any prime,
and to none of the negative centralizer-ghost
families. Moreover,
\[
 K_p^{{\rm ab},\times}\subset
                  \bigcap_{n\ge1}(B_p^\times)^{p^n}.
                                                        \tag{FF.40.3}
\]
Thus the exact \(\mu_\infty\) Kummer-kernel
input used in Claims~\ref{clm:FFcyclotomicreconstruct}--
\ref{clm:FFnilpotentreconstruct} fails here.
No reconstruction failure is asserted.
\end{claim}
\begin{proof}
The field \(B_p\) contains
\(K_p^{\rm ab}\), hence \(U_p\).
The valuation, residue and roots-of-unity
assertions follow. It is countable and
of transcendence degree zero because
it is algebraic over \(\mathbb Q\).

We verify the needed nonabelian groups
explicitly. For any finite number model
\(M/\mathbb Q_p\), wild inertia
\(P_M\) is free pro-\(p\) of
countably infinite rank. Choose eight
free generators and map them to
\(1+E_{i,i+1}\), \(1\le i\le8\),
in \(\operatorname{UT}_9(\mathbb F_p)\).
The words
\[
 u=[[x_1,x_2],[x_3,x_4]],\qquad
 v=[[x_5,x_6],[x_7,x_8]]
\]
belong to \(P_M''\), and map to
\(1+E_{1,5}\) and \(1+E_{5,9}\).
Their commutator maps to
\(1+E_{1,9}\ne1\). Thus
\(P_M''\) is nonabelian for every
prime, including \(p=2\).
In particular \(I''\ne1\), proving
that \(B_p\ne\overline{\mathbb Q}\).

For a finite extension \(E/B_p\),
choose a number field \(D\subset E\)
with \(E=DB_p\), and write
\(M=D_v\). Its absolute group is
\[
                 H=G_E=I''\cap G_M
                             \triangleleft G_M.
                                                        \tag{FF.40.4}
\]
It lies in \(P_M\) and contains
\(P_M''\), so is nonabelian.
Its centralizer in \(P_M\) is
trivial. Indeed, if \(a\ne1\)
centralized \(H\), the finitely
generated closed free pro-\(p\)
subgroup generated by \(a,u,v\)
would put two noncommuting elements
in the procyclic centralizer of
\(a\), contrary to the free-group
centralizer theorem used in
Claim~\ref{clm:FFcyclotomictest}.

The subgroup \(C_{G_M}(H)\) is
normal in \(G_M\), by (FF.40.4).
Its commutator with \(P_M\) lies
in \(C_{G_M}(H)\cap P_M=1\),
so it centralizes \(P_M\).
The fixed field of \(P_M\) is
a finite extension of \(T_p\).
It reconstructs by
Claim~\ref{clm:FFtamereconstruct};
the necessary condition in
Claim~\ref{clm:FF6} therefore makes
the centralizer of \(P_M\) in
\(\operatorname{Aut}(\overline{\mathbb Q})\)
trivial. Thus \(C_{G_M}(H)=1\).

Any global automorphism centralizing
\(H\) preserves its fixed field
\(E\). The rank-one henselian
valuation of \(E\) is unique:
two distinct such valuations would
make it algebraically closed, by
the two-valuation argument proved
in Claim~\ref{clm:FFstricttest},
whereas \(H\) is nonabelian.
Its unique prolongation to
\(\overline{\mathbb Q}\) is
therefore preserved. The global
centralizing automorphism belongs
to \(G_{\mathbb Q_p}\); a positive
power belongs to the open subgroup
\(G_M\) and is then trivial.
Torsion-freeness of \(G_{\mathbb Q_p}\),
verified in Claim~\ref{clm:FFcyclotomictest},
makes the automorphism itself
trivial. This proves (FF.40.2).

We next exclude the new nilpotent
families, including abstract field
embeddings. Let \(s=1\) for
odd \(p\), and \(s=2\) for
\(p=2\). For arbitrarily large
\(n\), the finite \(p\)-group
\[
 Q_n=\langle a,b:
       a^{p^n}=b^{p^n}=1,\quad
       bab^{-1}=a^{1+p^s}\rangle
                                                        \tag{FF.40.5}
\]
is a semidirect product of two
cyclic groups of order \(p^n\).
The displayed action has order
dividing \(p^n\), so the
presentation is consistent.
It is metabelian, and for every
fixed \(c\),
\[
           \gamma_{c+1}(Q_n)
                  =\langle a^{p^{sc}}\rangle,
            \qquad
           |\gamma_{c+1}(Q_n)|=p^{n-sc}
                   \quad(n>sc).
                                                        \tag{FF.40.6}
\]
Every finite \(p\)-group is a
quotient of \(I\), by the
Frattini lifting proof in
Claim~\ref{clm:FFnilpotentreconstruct}.
Since \(Q_n''=1\), these quotients
factor through \(I/I''\).
Consequently the image of
\(\gamma_{c+1}(I)\) in
\(I/I''\) is infinite, by
(FF.40.6).
If \(B_p\) were contained in
a finite extension \(L/K_{p,c}\),
the subgroup \(G_L\) would be
open in \(\gamma_{c+1}(I)\)
and contained in \(I''\).
That would make this image
finite, a contradiction.

Suppose instead that an abstract
embedding of \(B_p\) had image
in a finite extension \(L/K_{q,c}\).
Extend it to an automorphism
\(g\) of \(\overline{\mathbb Q}\).
The transported henselian valuation
of \(gB_p\) extends to a henselian
valuation on the algebraic extension
\(L\). Its residue characteristic
is \(p\); the given valuation
of \(L\) has residue characteristic
\(q\). Since \(L\) is not
algebraically closed, uniqueness
of a rank-one henselian valuation
forces equality of the valuations
and \(p=q\). Their unique
prolongations to the algebraic
closure then give \(g\in G_{\mathbb Q_p}\).
Conjugation by this group preserves
both \(I''\) and the closed
lower central series of \(I\).
Thus \(gB_p=B_p\), reducing to
the already excluded inclusion.

For completeness, the preceding
positive families are excluded
by their already verified criteria.
The field contains \(\mu_\infty\),
which is a nontrivial divisible
multiplicative subgroup, so it
cannot be a Kummer-faithful
sub-\(q\)-adic field for any
\(q\). The values of
\(\zeta_{q^m}-1\) have unbounded
\(q\)-denominators at every
residue-characteristic-\(q\)
place, excluding the bounded-wild
families at all primes.
Its transcendence degree zero
excludes exactly the families with an explicit
positive-transcendence hypothesis in
Claim~\ref{clm:FFclasses}(3)--(4) and the
finitely generated transcendental fields
of Claim~\ref{clm:FFfunction}.
We do not assert arbitrary subfield
closure for a proper normal envelope
over a transcendental finitely generated field.

The proper-normal-envelope families
are excluded by the same complete
two-valuation argument as in
Claim~\ref{clm:FFstricttest}.
Explicitly, for any number field
\(D\) in a finite extension
\(E/B_p\), choose
\(g\in G_D\) moving the selected
valuation. Such a \(g\) exists
because the decomposition group
of a finite place of a number
field is proper; the local-square,
global-nonsquare CRT construction
in that claim proves this directly.
The fields \(E\) and \(gE\)
then give two distinct rank-one
henselian valuations on their
compositum. They remain distinct
there because otherwise their
unique prolongations to the
algebraic closure would coincide.
Hence \(EgE=\overline{\mathbb Q}\).
The normal closure of \(E/D\)
is therefore \(\overline{\mathbb Q}\).
Together with the precise normal-closure
equivalence in Claim~\ref{clm:FFclasses}(2),
this excludes the formation subfield
families as stated there. It also
excludes containment in a
proper normal algebraic extension
of any number field, after
adjoining that number field if
necessary. It also excludes being
finite over such a proper normal
field: two conjugate copies of
\(G_E\) would then be open in
its normal Galois group but have
trivial intersection, forcing
the infinite group \(G_E\) to
be finite. Finally (FF.40.2)
excludes all the negative
centralizer-ghost families.

To prove (FF.40.3), take
\(a\in K_p^{{\rm ab},\times}\).
Over \(K_p^{\rm ab}\), adjoin
all compatible \(p\)-power
roots of every conjugate of
\(a\) under \(I\). There are
only finitely many such conjugates,
since \(a\) is algebraic over
\(K_p\). All roots of unity
are already in the base, so
this is an abelian pro-\(p\)
extension of \(K_p^{\rm ab}\),
stable under \(I\). Its Galois
group over \(K_p\) is an
extension of an abelian group
by an abelian group and hence
metabelian. It lies in \(B_p\).
In particular \(a\) has
compatible \(p\)-power roots
in \(B_p\), proving (FF.40.3).
The element \(a=p\) already
shows that the Kummer kernel
is larger than \(\mu_\infty\).
\end{proof}

The field \(B_p\), for example \(B_2\),
is consequently a specified remaining
test with the smallest possible
cardinality and transcendence degree
in characteristic zero. An
inclusion-minimal field is not
claimed. Whether \(\mathcal R(B_p)\)
holds remains \textbf{Open} in
this work. The failed exact
Kummer-kernel step identifies a
precise boundary of the preceding
proofs; it is not a topos map or
a refutation of centralizer
sufficiency.


\begin{claim}{FF.41}{Proved}\label{clm:FFmetabelianbarrier}
The invariant-lift step of
Claim~\ref{clm:FFmaxabelianreconstruct}
fails for an actual elliptic curve over
the metabelian field \(B_p\), even
after every finite change of local
arithmetic model. More precisely, put
\(H=I''=G_{B_p}\), and take
\[
                  A:\quad y^2+xy=x^3-p.
                                                        \tag{FF.41.1}
\]
There is a finite local field
\(M_0/\mathbb Q_p\) of degree at most
two, whose chosen global quadratic
model lies in \(K_p^{\rm ab}\), such
that the following holds. The space
\[
 S=\operatorname{Hom}_{H}
          (V_pA,\mathbb Q_p(1))
                                                        \tag{FF.41.2}
\]
is two-dimensional. Let \(N\) be the
subspace of functionals whose pushforward
of the point Kummer class is zero for
every \(a\in A(E)\) and every finite
extension \(E/B_p\). Then for every
finite extension \(M/M_0\),
\[
       N\simeq\mathbb Q_p(1),\qquad
       S/N\simeq\mathbb Q_p,\qquad
       S^{I_M}=0,\qquad
       (S/N)^{I_M}=S/N\ne0.
                                                        \tag{FF.41.3}
\]
The first two isomorphisms describe
the actual local Galois representations.
Consequently the visible invariant
class has no inertia-invariant lift.
This is a counterexample to a proposed
replacement lemma, not a counterexample
to \(\mathcal R(B_p)\).
\end{claim}
\begin{proof}
For the displayed integral Weierstrass
equation one computes
\[
             c_4=1,\qquad
             \Delta=p(1-432p),\qquad
             v_p(j)=-1.
                                                        \tag{FF.41.4}
\]
In particular \(j\ne0,1728\).
Tate's convergent inverse to the
\(j\)-series gives \(q\in\mathbb Q_p\)
with \(v_p(q)=1\) and a Tate curve
\(E_q\) with this \(j\)-invariant.
Two characteristic-zero elliptic curves
with this \(j\)-invariant become
isomorphic over a quadratic extension:
their geometric automorphism group
is \(\{1,-1\}\), so the isomorphism
torsor is quadratic. Thus
\(A_{M_0}\simeq E_q\) for a field
\(M_0\) of degree at most two.
A square class in \(\mathbb Q_p^\times\)
has a rational representative, by
openness of squares and density of
\(\mathbb Q\). Hence \(M_0\) is
the completion of a chosen quadratic
number field \(D_0\). Since its
inertia character is abelian,
\(D_0\subset K_p^{\rm ab}\).
One could also use the split node
with tangent cone \(y(y+x)\) in
(FF.41.1); allowing the quadratic
extension avoids needing a separate
split-reduction criterion.

The exact primary uniformization
input is Tate's Theorem~1, together
with the description of finite-order
points on pages~18--19 of the author
text \citep{FF_Tate1995}. It gives
the compatible local Galois sequences
\[
 \begin{gathered}
   \varphi:M^\times/q^{\mathbb Z}
                         \xrightarrow{\sim}A(M),\qquad
   0\longrightarrow\mathbb Z_p(1)
    \longrightarrow T_pA
    \longrightarrow\mathbb Z_p\longrightarrow0,\\
   0\longrightarrow\mathbb Q_p(1)
    \longrightarrow V_pA
    \longrightarrow\mathbb Q_p\longrightarrow0.
 \end{gathered}
                                                        \tag{FF.41.5}
\]
The extension class of the last
sequence is the local Kummer class
\(\kappa_M(q)\). This use of \(q\)
is entirely local: we do not assert
that \(q\) is algebraic over
\(\mathbb Q\).

Every element of \(I'\) lies in
\(G_{M_0}\), since the quadratic
character factors through an abelian
quotient. Its cyclotomic character
is one. In a basis adapted to
(FF.41.5), the image of \(I'\)
therefore consists of unipotent
upper-triangular matrices. They
commute, so \(H=I''\) acts trivially
on \(V_pA\). It also acts trivially
on \(\mathbb Q_p(1)\). Thus (FF.41.2)
is the full two-dimensional Hom space.
Dualizing and twisting (FF.41.5) gives
\[
 0\longrightarrow\mathbb Q_p(1)
  \longrightarrow S
  \xrightarrow{\ r\ }\mathbb Q_p
  \longrightarrow0,
                                                        \tag{FF.41.6}
\]
where \(r\) is restriction to the
torus Tate line. Its extension class
is \(\kappa_M(q)\), up to sign.

For any finite \(M/M_0\), valuation
shows that \(\kappa_M(q)\) is nonzero
in \(H^1(G_M,\mathbb Q_p(1))\):
the valuation component in the
rationalized Kummer completion is
\(v_M(q)\ne0\). Restriction to
\(I_M\) is injective, by
inflation--restriction and
\(\mathbb Q_p(1)^{I_M}=0\).
Consequently (FF.41.6) does not split
on \(I_M\). Its connecting map
from \(\mathbb Q_p\) is nonzero,
and the left term has no inertia
invariants. Therefore
\[
                         S^{I_M}=0.
                                                        \tag{FF.41.7}
\]

We now determine the invisible space,
with all finite extensions included.
If \(r(\lambda)=0\), take any
finite \(E/B_p\) and \(a\in A(E)\).
Choose a number field \(D\subset E\)
containing \(D_0\) and the coordinates
of \(a\), and put \(M=D_v\).
Then \(G_E\subset G_M\), using
the chosen henselian valuation.
By (FF.41.5), write
\(a=\varphi(u)\), \(u\in M^\times\).
Choose roots \(u_n\) recursively
with \(u_{n+1}^p=u_n\) and
\(u_n^{p^n}=u\). The points
\(\varphi(u_n)\) are coherent
\(p\)-power divisions of \(a\).
For \(g\in G_M\), their Kummer
cocycle is
\[
       g\varphi(u_n)-\varphi(u_n)
           =\varphi\bigl(g(u_n)/u_n\bigr),
       \qquad g(u_n)/u_n\in\mu_{p^n}.
                                                        \tag{FF.41.8}
\]
It lies in the torus subgroup at
every level, compatibly. Hence
\(\lambda\) kills the full Tate-module
Kummer class, not merely its separate
finite-level images. Restricting to
\(G_E\) proves
\(\ker(r)\subset N\).

Conversely let \(r(\lambda)\ne0\).
We construct an actual point of
\(A(B_p)\) which detects it.
The group \(H\) is itself a closed
subgroup of wild inertia, because
\(I/P\) is abelian. It is nontrivial
by Claim~\ref{clm:FFmetabeliantest}.
Closed subgroups of a free pro-\(p\)
group are free pro-\(p\), so \(H\)
has a nontrivial continuous quotient
\(\mathbb F_p\). This uses freeness
of \(H\) itself, not merely the
presence of a free subgroup.
Kummer theory and \(\mu_p\subset B_p\)
give \(a\in B_p^\times\) which is
not a \(p\)-th power.
Since \(\mu_{p^\infty}\subset B_p\),
its compatible Kummer class is a
homomorphism
\[
       \kappa_H(a):H\longrightarrow\mathbb Z_p(1)
                                                        \tag{FF.41.9}
\]
whose reduction modulo \(p\) is
nonzero. It remains nonzero after
rationalization.

We may normalize \(a\) to an
arbitrarily small principal unit
without losing this property.
All compatible roots of \(p\)
lie in \(B_p\): prime-to-\(p\)
roots lie in its tame subfield,
and the radical tower over
\(K_p^{\rm ab}\) is metabelian
over \(K_p\), as in (FF.40.3).
For any \(s\in\mathbb Q\), a
chosen \(p^s\) has zero rational
Kummer class on \(H\), since an
integer multiple of that class
is a multiple of \(\kappa_H(p)=0\).
Divide \(a\) by \(p^{v_p(a)}\),
and divide its residue by the
matching prime-to-\(p\) root of
unity in \(K_p\). The result
\(u_0\in B_p\) has positive
valuation of \(u_0-1\) and the
same nonzero rational Kummer class.
Raising it to a sufficiently large
\(p\)-power produces
\(u=u_0^{p^m}\) as close to one
as desired, still with nonzero
rational Kummer class.

Choose a number field
\(D\subset B_p\) containing \(D_0\)
and \(u\), and put \(M=D_v\).
The auxiliary analytic point
\(a_0=\varphi(u)\in A(M)\) lies
arbitrarily close to the identity.
Equations (FF.41.8)--(FF.41.9) give
\[
        \lambda_*\kappa_H(a_0)
              =r(\lambda)\kappa_H(u)\ne0.
                                                        \tag{FF.41.10}
\]
Here \(a_0\) need not have globally
algebraic coordinates, and is only
an intermediate local point.

After multiplying \(\lambda\) by
a nonzero integer, it is integral
on \(T_pA\). Choose \(n\) such
that the integral class in
(FF.41.10) has nonzero image modulo
\(p^n\). The subgroup
\([p^n]A(M)\) is open: in a
characteristic-zero formal coordinate
at the identity, multiplication by
\(p^n\) has nonzero derivative and
is an analytic isomorphism between
sufficiently small neighborhoods.
Thus the finite-level Kummer class
is constant on an open neighborhood
of \(a_0\).

An etale coordinate at the rational
identity, defined over \(D\), shows
that \(A(D^h)\) is dense in the
small identity neighborhood of
\(A(M)\). Indeed choose its coordinate
in the dense subfield \(D\); Hensel's
lemma supplies the inverse point
over the chosen henselization
\(D^h\). Since \(B_p\) is henselian
and contains \(D\), this chosen
\(D^h\) lies in \(B_p\). We can
therefore choose
\(a'\in A(D^h)\subset A(B_p)\)
with \(a'-a_0\in[p^n]A(M)\).
The nonzero finite-level evaluation
in (FF.41.10) persists for \(a'\).
On \(H\), all the representations
are trivial, so a nonzero integral
Kummer homomorphism into the
torsion-free module \(\mathbb Z_p(1)\)
remains nonzero rationally.
Thus \(\lambda\notin N\).

We have proved \(N=\ker(r)\).
Equations (FF.41.6)--(FF.41.7)
now give all of (FF.41.3).
The accompanying standard-library program
checks the invariant polynomial identity
(FF.41.4) and the three unitriangular
commutator identities of
Claim~\ref{clm:FFmetabeliantest} over
\(\mathbb Z\), hence at every prime.
It does not compute a global Tate
parameter or certify the infinite arguments.
\end{proof}

The most direct ghost attempt from
this example would assign to a
point its multiplicative Tate
uniformization parameter. Two
requirements remain unproved.
First the analytic parameter of
a globally algebraic point need
not be known to lie in the global
field \(B_p\); the uniformization
theorem only places it in a finite
\(p\)-adic completion. Second, a
homomorphism to the geometric
Tate module is not a map of full
etale topoi. Such a map must also
give compatible maps after removing
every finite set of target points.
A chosen analytic fundamental
domain does not supply these
Zariski-open restrictions. We have
constructed neither the required
actual point function nor this
compatible family of morphisms.
The constant-point construction is
already excluded by
Claim~\ref{clm:FFstrictdetection}.


\begin{claim}{FF.42}{Proved}\label{clm:FFmetabeliankernel}
Put \(F_p=K_p^{\rm ab}\), and define
\[
 W_p=\{a\in B_p^\times:
                  a^m\in F_p^\times
                  \text{ for some integer }m\ge1\}.
                                                        \tag{FF.42.1}
\]
For every finite extension \(E/B_p\),
let \(\kappa_E\) denote the rational
\(p\)-primary Kummer map. Then
\[
 \ker\bigl(\kappa_E:E^\times
             \longrightarrow H^1(G_E,\mathbb Q_p(1))\bigr)
       =\bigcap_{n\ge1}(E^\times)^{p^n}=W_p.
                                                        \tag{FF.42.2}
\]
In particular no new element in a
finite extension enlarges this kernel.
The subgroup recovers its abelian
constant field by the equation
\[
                     W_p\cap(1-W_p)=F_p\setminus\{0,1\}.
                                                        \tag{FF.42.3}
\]
There are two precise value-detection
consequences.
\begin{enumerate}
\item If \(a,b\in E\setminus\{0,1\}\)
have the same pair
\((\kappa_E(a),\kappa_E(1-a))\),
then \(a=b\) or both \(a,b\in B_p\).
Thus the pair is injective outside
\(B_p\), with no collision between
such a value and a value in \(B_p\).
\item If \(a,b\in E\) satisfy
\[
                  \kappa_E(a-c)=\kappa_E(b-c)
                  \quad\text{for every }c\in\mathbb Q
                  \text{ with }a,b\ne c,
                                                        \tag{FF.42.4}
\]
then \(a=b\) or both \(a,b\in F_p\).
Conversely any two values in \(F_p\)
satisfy (FF.42.4). Thus all rational
target translates detect every value
outside \(F_p\) exactly.
\end{enumerate}
\end{claim}
\begin{proof}
Let \(a\in E^\times\) have zero
rational Kummer class. Choose a finite
Galois number field \(D/\mathbb Q\)
containing \(a\) and finitely many
generators for \(E/B_p\). We may
replace \(E\) temporarily by
\(DB_p\), since restriction preserves
the zero class. Its local completion
\(M=D_v\) is a finite Galois extension
of \(\mathbb Q_p\). With the fixed
embedding into the local algebraic
closure, put
\[
 \begin{gathered}
 I=G_{K_p},\qquad J=I\cap G_M,\qquad
 H=I''\cap J=G_{DB_p},\\
 \Gamma=J/H,\qquad C=(I'\cap J)/H.
 \end{gathered}
                                                        \tag{FF.42.5}
\]
The Galois choice of \(M\) makes
\(J\) and \(H\) normal in \(I\),
so every conjugation used below acts
on this same model. The group \(C\)
is abelian, the quotient \(\Gamma/C\)
is abelian, and the cyclotomic
character \(\chi\) is trivial on
\(C\) and has infinite image on
\(\Gamma/C\). The latter follows
because \(J\) is open in \(I\).

The class of \(a\) on \(J\)
inflates from \(H^1(\Gamma,\mathbb Q_p(1))\),
since its restriction to \(H\) is
zero. The abelian cocycle identity
used in Claim~\ref{clm:FFcyclotomiccohomology}
gives
\[
                 H^1(\Gamma/C,\mathbb Q_p(1))=0.
\]
Inflation--restriction therefore makes
the map
\[
 H^1(\Gamma,\mathbb Q_p(1))
       \longrightarrow H^1(C,\mathbb Q_p(1))
                                                        \tag{FF.42.6}
\]
injective. Every \(g\in I'\)
centralizes \(C\): the commutator
of \(g\) with an element of
\(I'\cap J\) lies both in
\(I''\) and in \(J\), hence
in \(H\). Also \(\chi(g)=1\).
Thus \(g\) acts trivially on the
right side of (FF.42.6), and hence
on its left side. Functoriality of
Kummer classes gives
\[
                         \kappa_J(g(a)/a)=0
                         \qquad(g\in I').
                                                        \tag{FF.42.7}
\]

All elements \(g(a)/a\) lie in
the same finite field \(M\).
Restriction
\(H^1(G_M,\mathbb Q_p(1))
\longrightarrow H^1(J,\mathbb Q_p(1))\)
is injective: \(J=I_M\) and
\(\mathbb Q_p(1)^J=0\), so
inflation--restriction has zero
kernel. The rational Kummer kernel
of \(M^\times\) is precisely its
finite group of roots of unity.
Indeed valuation detects the
uniformizer component, and on a
sufficiently small principal-unit
group the \(p\)-adic logarithm
identifies the remaining rationalized
Kummer completion with the additive
local field. It follows from
(FF.42.7) that
\[
                            g(a)/a\in\mu(M)
                            \qquad(g\in I').
                                                        \tag{FF.42.8}
\]
Choose a single integer \(m\)
annihilating the finite group
\(\mu(M)\). Then \(a^m\) is
fixed by every \(g\in I'\).
Since \(a\) is globally algebraic,
\(a^m\in\overline{\mathbb Q}^{I'}=F_p\).
Moreover \(\mu_m\subset F_p\),
so \(F_p(a)/F_p\) is cyclic.
The field \(B_p\) is the maximal
abelian extension of \(F_p\), with
groups \(G_{F_p}=I'\) and
\(G_{B_p}=I''\). Consequently
\(a\in B_p\), and hence
\(a\in W_p\).

Conversely suppose \(a^m\in F_p^\times\).
For every \(n\), a root \(a_n\)
with \(a_n^{p^n}=a\) satisfies
\(a_n^{mp^n}=a^m\). Because
\(F_p\) contains all roots of
unity, the extension
\(F_p(a_n)/F_p\) is cyclic and
therefore lies in \(B_p\).
Choosing these roots recursively
gives compatible divisions of \(a\)
in \(B_p\), so its integral and
rational Kummer classes vanish
over every \(E\).
Finally \(G_E\) acts trivially on
\(\mathbb Z_p(1)\), making the
integral Kummer class a homomorphism
to a torsion-free module. Its
rationalization vanishes exactly
when it vanishes integrally. This
is equivalent to divisibility by
every \(p^n\); compatible choices
exist since all \(p\)-power roots
of unity are already in \(E\).
This proves (FF.42.2).

For (FF.42.3), the inclusion of
\(F_p\setminus\{0,1\}\) is immediate.
Conversely let \(x\in W_p\) and
\(1-x\in W_p\). For every
\(g\in I'\), there are roots
of unity \(\zeta,\eta\in F_p\)
such that
\[
                    g(x)=\zeta x,\qquad
                    g(1-x)=\eta(1-x).
\]
If \(g(x)\ne x\), their sum gives
\((\zeta-\eta)x=1-\eta\).
The case \(\zeta=\eta\) would
force both to equal one, contrary
to the choice of \(g\). Otherwise
the equality puts \(x\) in \(F_p\),
also contradicting \(g(x)\ne x\).
Thus every \(g\in I'\) fixes
\(x\), and \(x\in F_p\).

For the first value-detection assertion,
equality of the two classes puts
\[
                 \alpha=a/b\in W_p,\qquad
                 \beta=(1-a)/(1-b)\in W_p.
\]
If \(\alpha=\beta\), summing the
two defining equalities gives
\(\alpha=1\), so \(a=b\).
Otherwise
\(b=(1-\beta)/(\alpha-\beta)\)
and \(a=\alpha b\) both belong
to \(B_p\).

For the second assertion assume
\(a\ne b\). Choose distinct
rational \(c_0,c_1\) avoiding
\(a,b\). The two ratios
\((a-c_i)/(b-c_i)\in W_p\subset B_p\)
are distinct; if equal, subtraction
would make both one and give
\(a=b\). Solving the two linear
equations therefore puts
\(a,b\in B_p\).
Choose a finite Galois number field
\(D\) containing them. For every
\(g\in I'\), consider
\[
 R_g(T)=\frac{(g(a)-T)(b-T)}{(g(b)-T)(a-T)}
                                 \in D(T).
                                                        \tag{FF.42.9}
\]
For each rational \(c\) away
from its finitely many exceptional
values, (FF.42.4) puts
\(\alpha_c=(a-c)/(b-c)\) in
\(W_p\). A power of \(\alpha_c\)
is fixed by \(g\), so
\(R_g(c)=g(\alpha_c)/\alpha_c\)
is a root of unity. It belongs
to the same number field \(D\).
There are only finitely many roots
of unity in \(D\), so one value
occurs for infinitely many \(c\).
A rational function with infinitely
many occurrences of one value is
constant. Its value at infinity
is one, hence \(R_g(T)=1\).
Comparing coefficients gives
\[
              g(a)-a=g(b)-b=\delta,\qquad
              \delta(b-a)=0.
\]
Since \(a\ne b\), both are
fixed by \(g\). This holds
for every \(g\in I'\), proving
\(a,b\in F_p\). The converse
follows directly from (FF.42.2).
\end{proof}

This supplies exact value detection
from the collection of all rational
target translates, despite the
failure of the single-functional
invariant lift in
Claim~\ref{clm:FFmetabelianbarrier}.
It does not yet prove reconstruction:
one must still derive sufficiently
coherent arithmetic control of those
translated Kummer maps from an
actual pinned topos map, and handle
the values in \(F_p\) which the
displayed family does not distinguish.
Nor does it make an analytic Tate
parameter globally algebraic. The
local version of the same argument
has the local fixed field
\(\overline{\mathbb Q}_p^{I'}\)
as its constants; that field already
contains \(\mathbb Q_p\) and the
local parameter \(q\) in
Claim~\ref{clm:FFmetabelianbarrier}.


\begin{claim}{FF.43}{Proved}\label{clm:FFmetabelianpointdetection}
Let \(E/B_p\) be finite and let
\(a,b\in E\). There is a proper
finite extension \(L/E\) and
\(\beta\in L\setminus E\) such
that equality of the two pairs
\[
 \begin{split}
 &\bigl(\kappa_L(a-\beta),
                  \kappa_L(1-a+\beta)\bigr),\\
 &\bigl(\kappa_L(b-\beta),
                  \kappa_L(1-b+\beta)\bigr)
 \end{split}
                                                        \tag{FF.43.1}
\]
implies \(a=b\). In fact every
proper finite extension and every
\(\beta\in L\setminus E\) work.

Consequently the following criterion
detects the actual point function of
a pinned map. Let \(X/B_p\) be
of finite type and let
\(f,g:X_{\acute et}\to\mathbb A^1_{B_p,\acute et}\)
be pinned morphisms over \(BG_{B_p}\).
For every finite \(L/B_p\) and
\(\beta\in L\), set
\[
 Y_\beta=\mathbb A^1_L\setminus\{\beta,\beta+1\},
 \qquad
 U_\beta=f_L^{-1}(Y_\beta)\cap g_L^{-1}(Y_\beta)
                         \subset X_L.
                                                        \tag{FF.43.2}
\]
Suppose the pullbacks by \(f_L\)
and \(g_L\) of the two Kummer
classes of \(T-\beta\) and
\(\beta+1-T\) agree after restriction
to \(U_\beta\). Then the two
actual geometric point functions
are equal, including their values
in \(K_p^{\rm ab}\).
This is a uniqueness criterion for
point functions; it neither constructs
the comparison morphism nor asserts
that equality of point functions
alone identifies arbitrary topos maps.
\end{claim}
\begin{proof}
Every finite extension of \(B_p\)
has a nontrivial absolute Galois
group, by the free-group argument
in Claim~\ref{clm:FFmetabeliantest}.
Thus a proper finite extension
\(L/E\) exists. If
\(\beta\in L\setminus E\), both
\(a-\beta\) and \(b-\beta\)
lie outside \(E\), hence outside
\(B_p\); neither equals zero or
one. The first value-detection
assertion of
Claim~\ref{clm:FFmetabeliankernel},
applied over \(L\), proves the
first conclusion.

For the criterion, take any
\(x\in X(\overline{\mathbb Q})\).
Choose a finite \(E/B_p\) over
which \(x\) is defined. Pinning
and the retained relative structure
give actual values
\(a=f(x),b=g(x)\in E\), by
\citep[Proposition~3.6]{FF_CHW2024}.
Choose \(L/E\) and \(\beta\)
as above. The point \(x_L\)
lies in both inverse-image opens
in (FF.43.2), since neither \(a\)
nor \(b\) is \(\beta\) or
\(\beta+1\). Evaluating the two
assumed cohomology equalities at
\(x_L\) gives (FF.43.1), so
\(f(x)=g(x)\). This argument
works at every geometric point
and makes no exception for values
in \(K_p^{\rm ab}\).
\end{proof}

The criterion uses genuine operations
available for pinned morphisms:
finite scalar extension, inverse images
of Zariski opens, and pullback and
evaluation of Kummer classes. Its
remaining hypothesis is substantial.
The visible nonsplit functional in
Claim~\ref{clm:FFmetabelianbarrier}
does not by itself produce compatible
equalities for these translated pairs,
or a scheme morphism with which to
compare them. Deriving those equalities
is the unresolved existence step.


\begin{claim}{FF.44}{Proved}\label{clm:FFmetabelianaffine}
Let \(E/B_p\) be finite. Every
pinned automorphism of
\(\mathbb A^1_{E,\acute et}\) over
\(BG_E\), with its relative
structural isomorphism retained,
has actual geometric point function
\[
                          x\longmapsto ux+v
                          \qquad
                          (u\in E^\times,\ v\in E).
                                                        \tag{FF.44.1}
\]
The two coefficients are unique.
This is an existence theorem for
the actual point function in the
automorphism case. It does not
assert that the entire topos
automorphism equals the one induced
by (FF.44.1), and does not classify
arbitrary pinned morphisms.
\end{claim}
\begin{proof}
Write \(f:\overline{\mathbb Q}
\to\overline{\mathbb Q}\) for
the actual point function. It is
a bijection, because the inverse
topos automorphism supplies its
inverse. It is \(G_E\)-equivariant,
and for every finite \(L/E\)
it restricts to a bijection of
\(L\), by pinned relative point
evaluation \citep[Proposition~3.6]{FF_CHW2024}.
In particular \(f(0),f(1)\in E\)
are distinct. Postcompose with
the affine scheme automorphism
\[
                   T\longmapsto
                   \frac{T-f(0)}{f(1)-f(0)}.
\]
It suffices to prove that the
new point function is the identity.
Henceforth
\[
                            f(0)=0,\qquad f(1)=1.
                                                        \tag{FF.44.2}
\]

\emph{The cohomology of each puncture.}
Let \(a\in\overline{\mathbb Q}\),
and choose a finite extension
\(L/E\) containing \(a\). Then
\(f(a)\in L\). The scalar-extended
topos automorphism restricts to an
equivalence
\[
 (\mathbb A^1_L\setminus\{a\})_{\acute et}
       \longrightarrow
 (\mathbb A^1_L\setminus\{f(a)\})_{\acute et}.
                                                        \tag{FF.44.3}
\]
Indeed its bijective geometric point
function takes precisely \(a\)
to \(f(a)\), so the inverse-image
open is precisely the displayed
source. Geometric continuous
\(p\)-adic cohomology in degree
one, with coefficients
\(\mathbb Z_p(1)\), is free of
rank one, with basis
\(\kappa(T-a)\). The equivalence
(FF.44.3) therefore acts on that
integral geometric line by a unit
\(\alpha_a\in\mathbb Z_p^\times\).
The unit is unchanged by further
finite scalar extension; equivalently
one may compute after the algebraic
scalar extension to
\(\overline{\mathbb Q}\), obtained
as the inverse limit of finite
etale base extensions.

The Hochschild--Serre sequence and
the global lift \(\kappa(T-a)\)
show that the pullback of
\(\kappa(T-f(a))\), with rational
coefficients, is
\[
                         \alpha_a\kappa(T-a)+\eta_{a,L},
           \qquad \eta_{a,L}\in H^1(G_L,\mathbb Q_p(1)).
                                                        \tag{FF.44.4}
\]
The relative structural isomorphism
identifies the coefficient system
and its base pullback here. Evaluating
at any \(x\in L\setminus\{a\}\)
gives
\[
       \kappa_L(f(x)-f(a))
          =\alpha_a\kappa_L(x-a)+\eta_{a,L}.
                                                        \tag{FF.44.5}
\]
These equations always compare points
and classes over one common finite
base \(L\).

\emph{Recovering the affine lines.}
Put \(F=K_p^{\rm ab}\). For distinct
\(a,b\) and \(x\notin\{a,b\}\)
in \(\overline{\mathbb Q}\), choose
one finite \(L/E\) containing all
three. Subtracting the evaluations
at \(x\) and \(b\) in (FF.44.5)
cancels the arithmetic term and gives
\[
 \kappa_L\!\left(\frac{f(x)-f(a)}{f(b)-f(a)}\right)
     =\alpha_a
          \kappa_L\!\left(\frac{x-a}{b-a}\right).
                                                        \tag{FF.44.6}
\]
The corresponding comparison using
the puncture \(b\) gives
\[
 \kappa_L\!\left(\frac{f(b)-f(x)}{f(b)-f(a)}\right)
     =\alpha_b
          \kappa_L\!\left(\frac{b-x}{b-a}\right).
                                                        \tag{FF.44.7}
\]
The coefficients are nonzero. By
Claim~\ref{clm:FFmetabeliankernel},
both classes on the right vanish
exactly when \((x-a)/(b-a)\in F\).
Both on the left vanish exactly
when \((f(x)-f(a))/(f(b)-f(a))\in F\).
Including the endpoints, and using
the inverse automorphism for the
reverse inclusion, we obtain
\[
                 f\bigl(a+(b-a)F\bigr)
                     =f(a)+(f(b)-f(a))F.
                                                        \tag{FF.44.8}
\]
Thus \(f\) preserves all affine
\(F\)-lines, in both directions.

\emph{An elementary affine-geometry argument.}
We give the argument also in this
infinite-dimensional setting. Regard
\(\overline{\mathbb Q}\) as a vector
space over \(F\). Its dimension
is at least two. A bijection fixing
zero and preserving lines maps any
plane spanned by independent vectors
\(x,y\) onto the plane spanned by
\(f(x),f(y)\). To see containment,
each point off the two axes is on
a line joining a point of \(Fx\)
to a point of \(Fy\); for example
\(sx+ty\) is on the line joining
\(2sx\) and \(2ty\). Use the
inverse bijection for equality.
Noncollinearity is preserved as well.
Within such a plane, parallel lines
are characterized by disjointness,
so parallelism is preserved.
The parallelogram intersections give
\[
                      f(x+y)=f(x)+f(y)
                   \quad\text{if }x,y\text{ are independent.}
                                                        \tag{FF.44.9}
\]
For dependent vectors, introduce a
vector \(z\) outside their line.
First apply (FF.44.9) to
\(z=x+(-x+z)\) and to \(-x+z\)
to obtain \(f(-x)=-f(x)\).
If \(x+y\ne0\), compute
\(f(x+y+z)\) both as
\(f(x+y)+f(z)\) and as
\(f(x)+f(y+z)\). This gives
additivity also for dependent vectors;
the zero cases follow immediately.

For each nonzero vector \(x\),
line preservation defines a bijection
\(\sigma_x:F\to F\) by
\(f(cx)=\sigma_x(c)f(x)\).
For independent \(x,y\), compare
the two additive expansions of
\(f(c(x+y))\) to see
\(\sigma_x=\sigma_y=\sigma_{x+y}\).
For dependent nonzero vectors,
use a third independent vector.
Thus a single \(\sigma\) works
for every vector. Additivity of
\(f\) makes \(\sigma\) additive,
and comparing \(f(cd x)\) in
two ways makes it multiplicative.
It is bijective and takes one
to one, hence is an automorphism
of \(F\). In view of (FF.44.2),
\[
      f(x+y)=f(x)+f(y),\qquad
      f(cx)=\sigma(c)f(x)\ (c\in F),\qquad
      f|_F=\sigma.
                                                        \tag{FF.44.10}
\]
In particular \(f(F)=F\) and
\(f\) fixes \(\mathbb Q\).
Only two-dimensional spans were
used, so there is no finite-dimension
assumption hidden in this argument.

\emph{Recovering multiplication.}
In (FF.44.5) take \(a=0\).
Evaluating at \(x=1\) and using
(FF.44.2) gives \(\eta_{0,L}=0\).
Consequently, on every common finite
extension \(L/E\),
\[
                            \kappa_L(f(x))
                                   =\alpha_0\kappa_L(x).
                                                        \tag{FF.44.11}
\]
Let \(x,y\ne0\), and choose
one such \(L\) containing both.
It also contains all their images
under \(f\). For each rational
\(c\) with \(x+c\ne0\),
multiplicativity of Kummer classes
and (FF.44.11) yield
\[
       \kappa_L(f((x+c)y))
                     =\kappa_L(f(x+c))+\kappa_L(f(y)).
\]
Using (FF.44.10), this becomes
\[
       \kappa_L\!\left(\frac{f(xy)}{f(y)}+c\right)
                     =\kappa_L(f(x)+c).
                                                        \tag{FF.44.12}
\]
The equalities hold for all rational
translates at which both sides
are defined. The final assertion
of Claim~\ref{clm:FFmetabeliankernel}
therefore gives either
\[
                        \frac{f(xy)}{f(y)}=f(x)
\]
or both displayed values lie in
\(F\). If \(x\notin F\),
then \(f(x)\notin F\), so
only the first alternative is
possible. If \(x\in F\),
semilinearity already gives
\(f(xy)=\sigma(x)f(y)=f(x)f(y)\).
Zero factors are harmless. Hence
\(f\) is multiplicative everywhere.

The normalized point function is
therefore a field automorphism of
\(\overline{\mathbb Q}\).
Its \(G_E\)-equivariance says exactly
that it centralizes \(G_E\).
Claim~\ref{clm:FFmetabeliantest}
applies to every finite \(E/B_p\)
and makes this centralizer trivial.
Thus the normalized \(f\) is
the identity. Undoing the affine
normalization gives (FF.44.1), with
\(u=f_{
m original}(1)-f_{
m original}(0)\)
and \(v=f_{
m original}(0)\).
\end{proof}

The unresolved automorphism kernel
now has a precise form: after
composition with a unique affine
scheme automorphism, an additional
pinned automorphism, if one exists,
must fix every geometric point.
The proof above does not identify
its action on the entire etale
site, and does not provide the
missing existence theorem for
noninvertible pinned maps.


\subsection{Faithfulness of generic restriction on the pointwise kernel}

\begin{claim}{FF.45}{Proved}\label{clm:FFgenerickernel}
Let \(k\) be a characteristic-zero field, let \(X/k\)
be a smooth geometrically integral curve, and put
\(F=k(X)\). A pinned automorphism
\(\varphi:X_{\mathrm{\acute et}}\to
X_{\mathrm{\acute et}}\) over \(BG_k\) induces a generic
outer automorphism
\[
 \alpha_\varphi\in
 \operatorname{Aut}_{G_k}(G_F)/
              \operatorname{Inn}(G_{\bar k(X)}).
                                                        \tag{FF.45.1}
\]
If its actual function on \(X(\bar k)\) is the identity
and (FF.45.1) is trivial, then \(\varphi\) is isomorphic
to the identity over \(BG_k\).

Consequently, for every finite \(E/B_p\), the remaining
pointwise kernel in Claim~\ref{clm:FFmetabelianaffine}
injects into the relative birational outer automorphism
group of \(G_{E(T)}\). Its image preserves the geometric
conjugacy class of every finite-point inertia subgroup,
and its induced map on the corresponding residue Galois
group is the identity. In particular, an inner generic
action cannot hide a nontrivial automorphism of the full
etale topos.
\end{claim}

\begin{proof}
Write \(j:\operatorname{Spec}F\to X\), and choose a
separable closure of \(F\) containing \(\bar k\).
We retain all structural isomorphisms over \(BG_k\).

\emph{The generic action.}
An automorphism preserves the unique generic point and
therefore permutes the nonempty open subtopoi. The filtered
category of quasicompact separated etale objects over
these opens, with objects and arrows identified after
further shrinking, is the category of finite etale
\(F\)-schemes. Indeed, every finite separable extension
of \(F\), and every morphism between finitely many such
extensions, spreads out over a nonempty open of \(X\).
Conversely, two morphisms that agree generically agree
after shrinking. Thus \(\varphi\) induces an
autoequivalence of finite continuous \(G_F\)-sets.
The specified relative structure fixes the subcategory
of sets inflated from \(G_k\), with its specified
comparison. Galois-category reconstruction gives exactly
(FF.45.1), with inner ambiguity only in the geometric
kernel.

Assume from now on that the geometric point function is
the identity and (FF.45.1) is trivial. Choose the resulting
natural identity for the generic inverse-image functor,
over the category of \(G_k\)-sets. It extends from finite
continuous \(G_F\)-sets to all continuous \(G_F\)-sets:
each such set is a union of its finite orbits.

\emph{The maximal etale extension of a generic object.}
For a finite continuous \(G_F\)-set \(S\), the sheaf
\(j_*S\) is represented by the etale locus of the finite
normalization of \(X\) in the finite etale \(F\)-algebra
corresponding to \(S\). The normalization is finite
because \(X\) is a finite-type regular curve over a
field. To check the assertion, an \(F\)-map from the
generic fibre of an etale normal \(X\)-scheme extends
uniquely to its normalization map. At every point in its
image the extension of discrete valuation rings has
ramification index one and separable residue extension,
so the image lies in the etale locus. This gives precisely
the defining sections of \(j_*S\).

Every quasicompact separated etale \(X\)-scheme \(U\)
with generic fibre \(S\) is an open subobject
\[
                           U\ \lhook\joinrel\longrightarrow j_*S.
                                                        \tag{FF.45.2}
\]
For example, this follows by applying the same
normalization argument to each component: a separated
quasi-finite etale map into a normal curve is an open
subscheme of its finite normalization. All components of
\(U\) dominate \(X\), since etale maps are open.
The complement in the etale normalization locus consists
of finitely many closed points.

For a geometric point \(\bar x\) above a closed point
\(x\), choose a prolongation of its valuation to the
separable closure of \(F\). Let \(D_x\) and \(I_x\)
be the corresponding decomposition and inertia groups.
Then
\[
               (j_*S)_{\bar x}=S^{I_x},\qquad
               D_x/I_x=G_{\kappa(x)}.
                                                        \tag{FF.45.3}
\]
The stalk \(U_{\bar x}\) in (FF.45.2) is a
\(D_x\)-stable subset of \(S^{I_x}\). This description
also follows directly by base change to the strict
henselization at \(\bar x\): its unramified sheets are
exactly the \(I_x\)-fixed generic sheets.

\emph{The relevant normalizer.}
In this setting
\[
                            N_{G_F}(I_x)=D_x.
                                                        \tag{FF.45.4}
\]
Indeed, the fixed field \((F^{\rm sep})^{I_x}\) is the
fraction field of the strict henselization at the chosen
place. It has a henselian discrete rank-one valuation,
separably closed residue field of characteristic zero,
and value group \(\mathbb Z\). It is not separably
closed: a uniformizer has no square root in it. An element
normalizing \(I_x\) transports this valuation to another
rank-one henselian valuation on the same field. If the
two valuations were distinct, they would be incomparable,
and the two-valuation argument proved in
Claim~\ref{clm:FFstricttest} would make the field separably
closed. They therefore coincide. Its unique prolongation
to \(F^{\rm sep}\) is then preserved as well, so the
normalizing element belongs to \(D_x\). The reverse
inclusion is the usual normality of inertia in the
decomposition group.

\emph{The specialization comparison is forced.}
The chosen generic identity and the adjunction
\(j^*\dashv j_*\), together with invertibility of
\(\varphi^*\), give natural identifications
\[
                             \varphi^*j_*S\simeq j_*S.
                                                        \tag{FF.45.5}
\]
Actual pointwise identity gives an isomorphism
\(\bar x^*\varphi^*\simeq\bar x^*\), by
\citet[Proposition~3.6]{FF_CHW2024}. Combining it with
(FF.45.5) yields a natural automorphism of the functor
\[
       \{\text{finite continuous }G_F\text{-sets}\}
          \longrightarrow\{\text{sets}\},\qquad
                            S\longmapsto S^{I_x}.
                                                        \tag{FF.45.6}
\]
There is no arbitrary specialization permutation here.
The functor (FF.45.6) is prorepresented by the profinite
\(G_F\)-set \(G_F/I_x\). Concretely, write it as the
inverse system of finite \(G_F\)-sets
\(G_F/(I_xN)\), where \(N\) ranges over open normal
subgroups. A continuous map from this system to a finite
\(G_F\)-set is exactly an \(I_x\)-fixed element.
Pro-Yoneda therefore identifies the natural automorphisms
of (FF.45.6) with
\[
                      N_{G_F}(I_x)/I_x=D_x/I_x.
                                                        \tag{FF.45.7}
\]
Equivalently, the comparison acts on every \(S^{I_x}\)
by one element of \(D_x/I_x\). Hence it preserves
every \(D_x\)-stable subset, in particular
\(U_{\bar x}\). If the relative point comparisons are
chosen explicitly over \(BG_k\), their action on sets
inflated from \(G_k\) is the identity; this even makes
the element in (FF.45.7) trivial, since
\(G_{\kappa(x)}\to G_k\) is injective. Preservation of
the subsets is all that is needed.

It follows from (FF.45.5)--(FF.45.7) that \(\varphi^*\)
fixes the subobject (FF.45.2) inside \(j_*S\): its
generic stalk is fixed by construction and all its closed
geometric stalks are fixed by the preceding paragraph.
Geometric points detect equality of subobjects. This
produces an isomorphism \(\varphi^*U\simeq U\) for
every quasicompact separated etale \(U\). These
isomorphisms are natural. A morphism between two such
objects is determined by its generic restriction, because
the source is reduced with all components dominant and
the target is separated. Both ways around the proposed
naturality square agree generically and hence everywhere.

The quasicompact separated etale objects generate the
etale topos, and automorphisms preserve this subcategory;
see \citet[Theorem~A.8, Corollaries~A.9 and~A.11]
{FF_CHW2024}. Thus the natural identity extends to the
whole inverse-image functor. It respects the specified
structure over \(BG_k\): for a finite \(G_k\)-set the
corresponding constant etale cover is already its own
maximal extension \(j_*S\), and the generic identity
was chosen with precisely that relative compatibility.
This proves the first assertion.

For the final assertion take \(X=\mathbb A^1_E\).
The same specialization calculation, without assuming
the generic action trivial, identifies the prorepresented
functors attached to \(I_x\) and its transform. Thus
the generic action preserves the inertia conjugacy class
at each fixed geometric point. The relative point
comparison makes the conjugator geometric and makes the
map on the residue Galois group the identity. The first
assertion proves injectivity on the pointwise kernel.
\end{proof}

\paragraph{Exact birational scope at the metabelian field.}
For every finite \(E/B_p\), Koenigsmann's
\citet[Theorem~4.1, Corollaries~4.3 and~4.5]
{FF_Koenigsmann2003} apply: \(E\) is neither separably
closed nor real closed, and has a rank-one mixed
characteristic \((0,p)\) henselian valuation with
\(p\) dividing the order of \(G_E\).
They characterize the geometric decomposition and inertia
groups of the complete curve and the fundamental-group
quotients of its opens. Therefore the generic action in
Claim~\ref{clm:FFgenerickernel}, already fixing every
finite geometric place, fixes the remaining place at
infinity as well. Corollary~4.5 deduces a birational
Isom conclusion only when a separate fundamental-group
Isom conclusion is available for an open. It does not
prove that the pointwise birational stabilizer here is
trivial.

The more recent relative and semi-absolute theorems of
\citet[Theorems~A and~E, Definition~1.1,
Proposition~2.6]{FF_Tsujimura2026} retain an abelian-variety
Kummer-faithfulness hypothesis. For \(\Sigma=\{p\}\)
that hypothesis fails here: the elliptic curve in
Claim~\ref{clm:FFmetabelianbarrier} has trivial
\(G_{B_p}\)-action on its full \(p\)-adic Tate module,
so its nonzero rational \(p\)-power torsion points have
compatible \(p\)-divisions in \(B_p\). They are not
prime-to-\(p\) torsion. In fact every nonempty
\(\Sigma\) is excluded. If \(\ell\in\Sigma\)
and \(\ell\ne p\), the pro-\(p\) image of
\(G_{B_p}\) on the Tate module of any positive-dimensional
abelian variety injects into \(\operatorname{GL}_n(\mathbb F_\ell)\):
its intersection with the pro-\(\ell\) principal
congruence kernel is trivial. The image is therefore
finite. After a finite base extension all
\(\ell\)-power torsion is rational. A nonzero point
of this divisible \(\ell\)-primary subgroup is
\(\Sigma\)-divisible, since multiplication by each
other prime is bijective on that subgroup, but is not
\(\Sigma'\)-torsion. For \(\Sigma=\{p\}\) use
the preceding Tate example. This is a failure of that
source's hypothesis; the positive fields already proved
above show why it is not a reconstruction counterexample.
No Isom theorem for \(B_p\),
or vanishing of the remaining stabilizer, is inferred
from either source.

The commensurator theorem of
\citet[Theorem~6.8]{FF_HoshiMochizukiTsujimura2025}
also assumes TKND-AVKF. Its more general
Theorem~5.15 assumes QAA; Definition~5.12(ii)--(iii)
retains conditions on proper-cover cohomology and
function fields on all covers. These are not supplied
by the genus-zero conclusions below.


\subsection{Cyclic rational covers determine every inertia exponent}

\begin{claim}{FF.46}{Proved}\label{clm:FFallprimekernel}
Let \(k\) be a characteristic-zero field and \(X/k\)
a smooth geometrically integral curve with function field
\(F\). The following data determine an automorphism of
\(X_{\mathrm{\acute et}}\) over \(BG_k\): a continuous
automorphism \(\alpha\) of \(G_F\) over \(G_k\),
and a \(G_k\)-equivariant bijection
\(\theta:X(\bar k)\to X(\bar k)\), such that
\(\alpha\) sends the geometric conjugacy class of
\(I_x\) to that of \(I_{\theta(x)}\) for every
\(x\in X(\bar k)\). The resulting actual point
function is \(\theta\), and the resulting generic
action is \(\alpha\), with the usual geometric inner
ambiguity.

Now let \(E/B_p\) be finite, and let \(\varphi\)
be a pointwise-trivial pinned automorphism of
\(\mathbb A^1_{E,\mathrm{\acute et}}\) over \(BG_E\).
For every finite extension \(L/E\) and every positive
integer \(m\), its generic action fixes
\[
                    H^1(G_{L(T)},\mu_m)
\]
pointwise. Equivalently, it acts trivially on the maximal
abelian profinite quotient of \(G_{L(T)}\), after
every finite constant base change. All geometric inertia
exponents are one, at every prime. Neither assertion
identifies the remaining nonabelian generic action with
an inner automorphism.
\end{claim}

\begin{proof}
\emph{Reconstructing the generating site.}
We first give a converse to the construction in
Claim~\ref{clm:FFgenerickernel}. Fix prolongations and
geometric residue embeddings at the closed points of
\(X\). The category of quasicompact separated etale
\(X\)-schemes has the following description:
\[
 \begin{gathered}
 \left(S,\ (B_x)_{x\in|X|_{\rm cl}}\right),
 \qquad S\text{ a finite continuous }G_F\text{-set},\\
 B_x\subset S^{I_x}\text{ a }D_x\text{-stable subset}
       \qquad(x\in|X|_{\rm cl}).
 \end{gathered}
                                                        \tag{FF.46.1}
\]
where \(B_x=S^{I_x}\) for all but finitely many closed
\(x\). Its morphisms are the \(G_F\)-maps on generic
sets that take every \(B_x\) into its counterpart.
Indeed, form the finite normalization belonging to
\(S\), take its etale locus \(j_*S\), and remove
the finitely many closed points specified by the omitted
\(D_x\)-orbits. This reconstructs the object. Conversely,
Claim~\ref{clm:FFgenerickernel} identifies every object
as precisely such an open. A generic map extends to the
finite normalizations, and its image lies in the indicated
open exactly when the conditions on the \(B_x\) hold.
This proves the asserted description of arrows as well.
Multiple generic components are allowed in \(S\).
Nonseparated objects are not needed: this separated
subcategory is the generating site by
\citet[Theorem~A.8 and Corollary~A.11]{FF_CHW2024}.

Let \(N=G_{\bar k(X)}\). For a chosen geometric
\(x\), choose \(g_x\in N\) such that
\[
                  \alpha(I_x)=g_x I_{\theta(x)}g_x^{-1}.
                                                        \tag{FF.46.2}
\]
On generic sets use pullback by \(\alpha\), so the new
action on \(S\) is \(g\cdot_{\rm new}s=\alpha(g)s\).
Replace the local subset at \(x\) by
\[
                                g_x B_{\theta(x)}.
                                                        \tag{FF.46.3}
\]
This is a subset of \((\alpha^*S)^{I_x}\). It is
\(D_x\)-stable because (FF.45.4) and (FF.46.2) imply
\(\alpha(D_x)=g_xD_{\theta(x)}g_x^{-1}\).
Another choice of \(g_x\) differs on the right by an
element of \(D_{\theta(x)}\), again by (FF.45.4), and
therefore gives the same subset. Changes of valuation
prolongations give the usual identical description up
to canonical generic-set isomorphism. The
\(G_k\)-equivariance of \(\theta\) makes the rule
well-defined on closed-point orbits. Only finitely many
full subsets are changed, since \(\theta\) bijects
those orbits.

The rule preserves arrows and has an inverse obtained
from \(\alpha^{-1},\theta^{-1}\). Finite limits are
computed on the generic and closed geometric stalks, as
are jointly surjective etale covering families. Thus it
is an equivalence of the generating sites and induces
an equivalence of their topoi. On sets inflated from
\(G_k\), the generic action is the specified identity
and \(g_x\in N\) acts trivially. Hence the equivalence
has the required structure over \(BG_k\).

For completeness, the residue-field identification is
also determined. Since \(\theta\) is an equivariant
bijection, the stabilizers of \(x\) and \(\theta(x)\)
in \(G_k\) coincide. Their embedded finite residue
fields inside \(\bar k\) are consequently the same.
The isomorphism on \(D_x/I_x\) has the identity
projection to this stabilizer, since both \(\alpha\)
and the conjugation in (FF.46.2) are relative over
\(G_k\). The stalk identifications in (FF.46.3)
therefore identify the geometric point functor at \(x\)
with the target point functor at \(\theta(x)\), with
its specified base comparison. By
\citet[Proposition~3.6]{FF_CHW2024}, the actual point
function is exactly \(\theta\). The generic action
is the prescribed one by construction. This proves the
first assertion.

\emph{The cyclic rational cover.}
Fix now \(L/E\) finite and \(a\in L\), and base
change \(\varphi\) to \(L\). Its point function is
still the identity. On
\(V=\mathbb A^1_L\setminus\{a\}\), geometric first
cohomology with \(\mu_m\)-coefficients is generated
by \(\kappa_m(T-a)\). The equivalence therefore gives
\[
 \varphi^*\kappa_m(T-a)
       =u_m\kappa_m(T-a)+\eta_m,
 \quad u_m\in(\mathbb Z/m\mathbb Z)^\times,
 \quad\eta_m\in H^1(G_L,\mu_m).
                                                        \tag{FF.46.4}
\]
The displayed geometric generator is already lifted by
the indicated global unit, so the remaining term is
constant, by inflation--restriction. Evaluation at
\(T=a+1\) gives \(\eta_m=0\).

The field \(L\) contains \(\mu_m\). Thus the
extension
\[
              L(T)\subset L(U),\qquad T=U^m+a,
                                                        \tag{FF.46.5}
\]
is cyclic Galois of degree \(m\). Its absolute Galois
subgroup is the kernel of the homomorphism
\(\kappa_m(T-a):G_{L(T)}\to\mu_m\).
By (FF.46.4), the generic automorphism \(\alpha\)
preserves this subgroup. Its restriction
\(\alpha_U\) is an automorphism of \(G_{L(U)}\)
over \(G_L\).

The exact local theory in
\citet[Theorem~4.1 and Corollaries~4.3--4.5]
{FF_Koenigsmann2003} applies to \(L\), and after every
finite constant extension, both for \(\mathbb P^1_T\)
and \(\mathbb P^1_U\). It gives the associated
\(G_L\)-equivariant permutation \(\theta_U\) of
geometric places. More explicitly, first make a chosen
geometric place rational over a finite constant field
\(L'/L\). Corollary~4.3 characterizes its arithmetic
decomposition conjugacy class. Any conjugating element
can be multiplied by the inverse of an element of the
target decomposition group with the same projection to
\(G_{L'}\), because that projection is surjective.
The resulting conjugator is geometric. Applying this
argument over all finite constants, uniqueness of the
recovered valuation classes makes the point permutations
compatible with scalar extension and with the
\(G_L\)-action. This is the geometric conjugacy
transport required in (FF.46.2).
This permutation lies above the
identity permutation on \(\mathbb P^1_T\). To see
the needed compatibility explicitly, choose a place of
the common separable closure. Its decomposition and
inertia groups for \(L(U)\) are the intersections of
the corresponding groups for \(L(T)\) with
\(G_{L(U)}\). The subgroup is preserved by
\(\alpha\), so these intersections and their place
restrictions are transported compatibly. Uniqueness in
the local characterization identifies the transported
places with those of \(\theta_U\). In particular,
\[
                           \theta_U(u)^m=u^m
                  \qquad(u\in\bar L).
                                                        \tag{FF.46.6}
\]
The unique place over \(T=\infty\) is fixed, as is
the unique place \(U=0\) over \(T=a\). Therefore
\(\alpha_U,\theta_U\) satisfy the first assertion
for the affine line \(\mathbb A^1_U\). They define
a pinned automorphism of its full etale topos over
\(BG_L\), rather than merely of its generic finite
cover category.

Claim~\ref{clm:FFmetabelianaffine} now applies to this
lifted automorphism. Its actual point function is affine.
Since it fixes zero and satisfies (FF.46.6), it has the
form
\[
                       \theta_U(u)=c u,
                       \qquad c\in\mu_m(L).
                                                        \tag{FF.46.7}
\]
On the other hand, (FF.46.4) says that \(\alpha\)
acts on the cyclic deck quotient of (FF.46.5) by the
power \(u_m\). Conjugation and transport of inertia
classes are compatible. Indexing the deck rotations by
their action on the geometric points of the cover,
this gives
\[
                    \theta_U(\zeta u)
                          =\zeta^{u_m}\theta_U(u)
                    \qquad(\zeta\in\mu_m).
                                                        \tag{FF.46.8}
\]
Combining
(FF.46.7) and (FF.46.8) gives
\(\zeta^{u_m}=\zeta\) for every \(\zeta\in\mu_m\).
Hence \(u_m=1\) in \(\mathbb Z/m\mathbb Z\).
We have proved, for every \(a\) and \(m\),
\[
                  \varphi^*\kappa_m(T-a)=\kappa_m(T-a).
                                                        \tag{FF.46.9}
\]

\emph{Descent of an arbitrary rational Kummer class.}
For \(r\in L(T)^\times\), take its nonvanishing
open \(V_r\subset\mathbb A^1_L\) and a finite
Galois constant extension \(L'/L\) splitting its
divisor. Over \(L'\), it is a constant times a product
of integer powers of the linear units in (FF.46.9).
The relative structural isomorphism fixes the constant
classes. Thus
\(\delta=\varphi^*\kappa_m(r)-\kappa_m(r)\)
restricts to zero on \((V_r)_{L'}\).
Inflation--restriction for this purely constant cover
puts \(\delta\) in the image of
\(H^1(\operatorname{Gal}(L'/L),\mu_m)\), hence in
the constant subgroup from \(H^1(G_L,\mu_m)\).
Choose \(t_0\in V_r(L)\). Actual pointwise identity
gives \(t_0^*\delta=0\), and evaluation at \(t_0\)
is a left inverse on constant classes. Therefore
\(\delta=0\) before the scalar extension as well.

Every element of \(H^1(G_{L(T)},\mu_m)\) is a
Kummer class. Moreover \(\mu_m\subset L\), so
these cohomology groups are exactly the continuous
characters with values in \(\mu_m\). Their equality
for every \(m\) is equivalent to equality on the full
abelian profinite quotient. The exponents on finite
geometric inertia are detected by (FF.46.9), and the
single peripheral relation on the projective line gives
the same exponent at infinity. All of them are one.
\end{proof}

This improves the direct \(p\)-primary calculation:
the other primes are controlled by the rational cyclic
covers and the already-proved affine point-function
rigidity, even though point Kummer classes at those
primes vanish over \(B_p\). A nonabelian generic
automorphism acting trivially on all these abelian
quotients has not been excluded.


The distinction from inner triviality is concrete even
for an abstract four-punctured line. Let \(P\) be the
free pro-\(p\) group on \(x,y,z\), with fourth
peripheral generator \((xyz)^{-1}\). The automorphism
\[
 x\longmapsto(xy)x(xy)^{-1},\qquad
 y\longmapsto(xy)y(xy)^{-1},\qquad
 z\longmapsto z
                                                        \tag{FF.46.10}
\]
has inverse given by conjugating the first two generators
by \((xy)^{-1}\). It fixes \(xy\), \(xyz\), every
peripheral conjugacy class, and abelianization. It is not
inner. In the torsion-free class-two pro-\(p\) quotient,
the independent central commutators are
\([x,y],[x,z],[y,z]\). Conjugation by an element with
abelian coordinates \((a,b,c)\) could agree with
(FF.46.10) on \(x,y\) only if \(a=b=1,c=0\).
Its commutator with \(z\) would then be
\([x,z][y,z]\ne1\), a contradiction. This quotient
exists also for \(p=2\): use coordinates in
\(\mathbb Z_p^3\times\mathbb Z_p^3\), with the
central product correction
\((a_i a'_j)_{i<j}\), so no division by two is involved.
This is an abstract obstruction to an abelianization-only
argument. No compatibility of (FF.46.10) with the actual
\(G_E\)-action, with every finite cover, or with the
full affine-line etale topos is asserted.

The class-two calculation was executed modulo 2, 3, 4, 5 and 8. The order-64 model was checked on all 4096 products. The program and saved output record these finite checks:

\path{agent_F/F_metabelian_checks.py}.


\subsection{The pointwise kernel fixes the genus-zero part of the site}

\begin{claim}{FF.47}{Proved}\label{clm:FFgenuszerokernel}
Let \(E/B_p\) be finite and let \(\varphi\) be a
pointwise-trivial pinned automorphism of
\(\mathbb A^1_{E,\mathrm{\acute et}}\) over \(BG_E\).
Its inverse-image functor is naturally isomorphic to the
identity on the full subcategory of quasicompact separated
etale objects whose smooth projective geometric components
have genus zero. The same assertion holds after every
finite constant extension.

This constraint is strictly stronger than peripheral and
abelianization control for an abstract free profinite
four-puncture group. Specifically, the word automorphism
\[
 x\longmapsto(xy)x(xy)^{-1},\qquad
 y\longmapsto(xy)y(xy)^{-1},\qquad z\longmapsto z
                                                        \tag{FF.47.1}
\]
of the free \emph{profinite} group on \(x,y,z\) moves
a degree-three genus-zero cover, despite fixing every
peripheral conjugacy class and abelianization. Hence it
cannot be the geometric action of such a \(\varphi\)
on a four-punctured open. This is a statement about the
full profinite word automorphism (FF.47.1); it is not an
assertion about every full automorphism whose pro-\(p\)
image equals the corresponding pro-\(p\) word automorphism.
\end{claim}

\begin{proof}
We first take a geometrically connected etale object
\(Y\to\mathbb A^1_E\) with projective completion
\(C\) of genus zero, and write
\(Y'=\varphi^*Y\). Pullback to the slice topoi gives
an equivalence \(Y'_{\mathrm{\acute et}}\to
Y_{\mathrm{\acute et}}\) over the specified base.
It induces an isomorphism of the generic Galois groups
over \(G_E\). The constant fields remain \(E\),
because the relative generic action preserves the
projection to \(G_E\). Let \(C'\) denote the smooth
projective completion of \(Y'\).

Koenigsmann's local theory identifies all geometric
inertia groups on \(C\) and \(C'\), not merely those
at the points of the two displayed opens. Killing their
closed normal span gives the etale fundamental group of
the complete curve, compatibly over \(G_E\). Hence
the geometric fundamental groups of \(C\) and \(C'\)
are isomorphic. In characteristic zero the geometric
abelianization has rank \(2g\), by
\citet[Expos\'e~XIII, Corollary~2.12]{FF_SGA1}.
It follows that \(C'\) also has genus zero.

The same local theory supplies a compatible
\(G_E\)-equivariant point bijection
\[
                         h:C'(\bar E)\longrightarrow C(\bar E).
                                                        \tag{FF.47.2}
\]
Over the displayed opens it is the actual point function
of the slice equivalence. Indeed, the specialization
functors \(S\mapsto S^{I_x}\) used in
Claim~\ref{clm:FFgenerickernel} identify the inertia
class of every actual point with precisely its transported
class. At a rational point the arithmetic conjugator can
be made geometric using the surjection of its decomposition
group onto the constant Galois group, as explained in
Claim~\ref{clm:FFallprimekernel}.

Choose a finite Galois extension \(L/E\) over which
\(C\) has a rational point \(c\). Then
\(c'=h^{-1}(c)\) is \(L\)-rational as well.
A smooth projective genus-zero curve with a rational
point is \(\mathbb P^1\). Choose coordinates on
\(C_L\) and \(C'_L\) taking \(c\) and \(c'\)
to infinity. The generic group isomorphism, together
with (FF.47.2), transports the inertia classes of the
two remaining affine lines. The site construction in
Claim~\ref{clm:FFallprimekernel}, applied after these
coordinate identifications, gives a pinned automorphism
of the affine-line etale topos whose actual point function
is \(h\) in these coordinates.

Claim~\ref{clm:FFmetabelianaffine} makes this function
affine. It therefore extends to an \(L\)-isomorphism
\(C'_L\to C_L\), realizing all of (FF.47.2).
It carries \(Y'_L\) onto \(Y_L\), because on these
opens the actual point function is the bijection of the
slice equivalence. This isomorphism descends to \(E\):
the point bijection is \(G_E\)-equivariant, so every
Galois conjugate of the \(L\)-isomorphism has the same
geometric point function and is therefore equal to it.
This is ordinary finite Galois descent of scheme maps.
We obtain a unique \(E\)-isomorphism
\[
                                  h_Y:Y'\xrightarrow{\sim}Y
                                                        \tag{FF.47.3}
\]
with the required actual point function.

It is an isomorphism over \(\mathbb A^1_E\).
The slice equivalence lies over \(\varphi\), whose
base point function is the identity. Thus both composites
from \(Y'\) to \(\mathbb A^1_E\) have the same
geometric point function, and the corresponding scheme
maps are equal. For any arrow \(Y\to Z\) between
two objects under consideration, the same pointwise
compatibility gives
\(h_Z\circ\varphi^*(Y\to Z)=(Y\to Z)\circ h_Y\).
Equality follows because the sources are reduced and
geometric points are dense. This proves naturality.

For a connected object with a finite constant field
extension, apply the same argument over that extension;
the relative structure identifies the extension on both
sides. For a disconnected object, use its finitely many
connected components. Equivalently, one can split all
components and choose their rational points over one
finite Galois constant extension, then descend the unique
pointwise isomorphisms. This proves the full first
assertion and its compatibility with further constants.

\emph{A finite quotient detecting the abstract word automorphism.}
Write \(w=(xyz)^{-1}\). In \(S_3\), assign
\[
           (x,y,z,w)\longmapsto
           \big((12),(23),(12),(13)\big).
                                                        \tag{FF.47.4}
\]
The four permutations have product one and generate
\(S_3\). By the characteristic-zero presentation of the
geometric fundamental group of a four-punctured line
in \citet[Expos\'e~XIII, Corollary~2.12]{FF_SGA1},
this gives a connected degree-three etale cover over
\(\bar E\). Each of the four branch permutations
is a transposition. Its smooth completion therefore has
\[
                         2g-2=-2\cdot3+4=-2,
\]
so has genus zero. The cover and its map descend to a
finite constant extension of \(E\), by finite
presentation.

The automorphism (FF.47.1) fixes \(xy\) and \(z\),
and therefore fixes \(w\). It conjugates the first
two permutations in (FF.47.4) by their nontrivial
three-cycle product, leaving the last two unchanged.
The original and transformed degree-three permutation
representations are not isomorphic. A simultaneous
conjugator would have to centralize both \((12)\) and
\((13)\), so would be the identity; but the first
permutation has changed. Thus (FF.47.1) moves this
genus-zero cover and contradicts the first assertion
if proposed as the action of \(\varphi\).

This finite calculation was executed in
\texttt{F\_metabelian\_checks.py}; its output verifies
the product, transitivity, Riemann--Hurwitz integer, and
the absence of any of the six possible simultaneous
conjugators. Since \(S_3\) is not a \(p\)-group,
the test concerns the stated free profinite word action
and does not identify all possible lifts of a pro-\(p\)
action. No arithmetic-compatible ghost is constructed.
\end{proof}

The genus-zero subcategory is not a generating family
for the whole curve etale site. For example, a nonempty
etale object with projective geometric genus at least
one cannot be covered by genus-zero etale curves: a
dominant map between their projective completions would
contradict Riemann--Hurwitz. Nor has any cofinality of
Galois closures of genus-zero covers been established.
The action on arbitrary positive-genus covers remains the missing
part of the birational kernel problem.


\subsection{Genus-zero generation of abelian covers and elliptic Galois covers}

\begin{claim}{FF.48}{Proved}\label{clm:FFlatteskernel}
Let \(E/B_p\) be finite, and let \(\varphi\) be a pinned
automorphism of \((\mathbb A^1_E)_{\acute et}\) over \(BG_E\)
whose actual geometric point function is the identity.
The natural scheme realizations of
Claim~\ref{clm:FFgenuszerokernel} extend to the following
objects, after every finite extension of constants:
finite fibre products of the genus-zero objects, their
open subobjects, and finite etale quotients by finite
groups. In particular they extend to every subcover of
the Galois closure of a connected finite etale
genus-zero cover of an open of \(\mathbb A^1\).

They consequently extend to every etale restriction of
a finite abelian cover of \(\mathbb P^1_E\), of
arbitrary genus, and to every finite Galois cover
\[
                  \pi:C\longrightarrow\mathbb P^1_E
       \quad\hbox{with }g(C)=1.
                                                        \tag{FF.48.1}
\]
Here ``extend'' means that the induced slice point
function is realized by a unique scheme isomorphism
\(\varphi^*V\to V\) over the chosen base open, and
these isomorphisms are natural in every morphism between
the stated objects. The statement includes a quadratic
cover with four geometric branch points. It also
includes every degree-three cover with genus-one
source, without a Galois assumption.

Here \(C\) is smooth and geometrically integral.
Thus, after identifying the inverse image of any
connected cyclic torsor of the stated kind with the
original one, its possible fibrewise deck rotation is
a single global deck transformation. Point-dependent
rotations on cyclic covers of arbitrary genus are
excluded. No generation or cofinality
claim for all positive-genus covers is asserted.
\end{claim}

\begin{proof}
All constructions below may be made after removing
finitely many points of the base. The resulting objects
are quasicompact separated etale objects, and their
inverse images are represented by objects of the same
kind. We keep the structural comparisons over the
constant field throughout.

\emph{Closure under the stated constructions.}
Say that an etale object \(V\to X\), for an open
\(X\subseteq\mathbb A^1_E\), is realized if its
actual slice point function is induced by an
\(X\)-isomorphism
\[
                       a_V:\varphi^*V\xrightarrow{\sim}V.
                                                        \tag{FF.48.2}
\]
Such an isomorphism is unique: two maps between these
reduced finite-type schemes, with separated target, that
agree on all geometric points agree as scheme maps.
Moreover, for any arrow \(V\to W\) between realized
objects, the two maps \(\varphi^*V\to W\) formed
using (FF.48.2) agree on geometric points by the
naturality of point evaluation. They therefore agree
as scheme maps. Realizations are consequently natural
on every full subcategory on which they exist.

If \(V_1,\ldots,V_r\) are realized, their fibre
product is realized by the product of the maps
\(a_{V_i}\). Indeed inverse image preserves finite
limits, and a geometric point of the product is
determined by its projections. If \(W\subseteq V\)
is an open subobject of a realized object, the image of
\(\varphi^*W\) under \(a_V\) is precisely \(W\):
the two opens have the same geometric points. Restriction
of \(a_V\) gives \(a_W\). This includes open-and-closed
connected components. These conclusions use the actual
point-function comparison, rather than just invariance
of the isomorphism class of \(V\).

Suppose that a finite group \(H\) acts on a realized
finite etale object \(Z\to X\). Naturality for every
action arrow identifies the transported action with the
given action on \(Z\). Inverse image preserves its
coequalizer. The finite etale quotient therefore has
an induced isomorphism
\[
             \varphi^*(Z/H)\ \xrightarrow{\sim}\ Z/H.
                                                        \tag{FF.48.3}
\]
It realizes the slice point function, because the finite
etale surjection \(Z\to Z/H\) is surjective on
geometric points and that point function is natural.

For clarity, the Galois-closure application can be made
entirely inside these constructions. If \(V\to X\)
is connected finite etale of degree \(d\), take in
\(V^d_X\) the open-and-closed subscheme of ordered
\(d\)-tuples of pairwise distinct geometric sheets.
It is the frame torsor of \(V\). A connected component
\(Z\) is its Galois closure, with group the monodromy
image in \(S_d\). Once \(V\) is realized, its frame
torsor and each specified component are realized by
the preceding paragraph. Each finite subcover of
\(Z\) is a quotient \(Z/H\), so (FF.48.3) applies.
This proves the first assertion using
Claim~\ref{clm:FFgenuszerokernel}.

\emph{All abelian covers of the rational base.}
After a finite extension of constants, every connected
cyclic cover of degree \(m\) is a Kummer cover.
Indeed all \(m\)-th roots of unity are available,
and Kummer theory applies to its cyclic function-field
extension. After a further finite extension that splits
the divisor and extracts a constant root, write it as
\[
       y^m=f(T)=c^m\prod_{i=1}^r(T-a_i)^{e_i},
           \qquad c\ne0,\quad e_i\in\mathbb Z.
                                                        \tag{FF.48.4}
\]
Remove the finitely many \(a_i\), the original affine
infinity, and any additional omitted base points.
Each cover \(u_i^m=T-a_i\) of the resulting open
has genus-zero completion and is realized by
Claim~\ref{clm:FFgenuszerokernel}. Their finite fibre
product is therefore realized, as are its connected
components and finite group quotients.

More explicitly, the classes of the distinct
\(T-a_i\) are independent in
\(\bar E(T)^\times/(\bar E(T)^\times)^m\), as is
seen from their respective valuations at \(a_i\).
Their compositum is thus a connected Galois cover
with group \(\mu_m^r\). The map
\[
       (u_1,\ldots,u_r)\longmapsto
                     y=c\prod_i u_i^{e_i}
                                                        \tag{FF.48.5}
\]
identifies (FF.48.4) with its quotient by the kernel
of the character
\((\zeta_i)\mapsto\prod_i\zeta_i^{e_i}\).
The character is surjective when the original cover
is geometrically connected of degree \(m\). If
constant splitting produces several components, apply
this argument to each component and descend their
union. This proves the cyclic assertion. A finite
abelian Galois extension is a compositum of cyclic
subextensions, by the elementary structure theorem
for its finite abelian Galois group. Its cover is a
connected component of the corresponding finite fibre
product, so every abelian cover is realized as well.

\emph{The elliptic Latt\`es construction.}
We first give the quadratic case explicitly. After a
finite Galois extension \(L/E\), a branch point of
\(\pi\) has an \(L\)-rational point above it, which
we choose as the origin of the elliptic curve \(A=C_L\).
The involution of the double cover is then \([-1]\).
Enlarge \(L\) so that all points of \(A[3]\) are
rational. Write
\[
                 \pi:A\longrightarrow A/\{\pm1\}
                                      =\mathbb P^1_L.
\]
Multiplication by three commutes with \([-1]\), and
therefore induces a rational morphism of degree nine
\[
          f_3:\mathbb P^1_L\longrightarrow\mathbb P^1_L,
                 \qquad f_3\circ\pi=\pi\circ[3].
                                                        \tag{FF.48.6}
\]
Let \(B\) be the branch-value set of \(\pi\), and
put \(X=\mathbb P^1_L\setminus(B\cup\{\infty\})\),
with \(\infty\) interpreted in the originally chosen
affine base coordinate. The maps
\[
 \begin{split}
  Z&=(\pi\circ[3])^{-1}(X)\longrightarrow X,\\
  V&=f_3^{-1}(X)\longrightarrow X
 \end{split}
                                                        \tag{FF.48.7}
\]
are finite etale: \([3]\) is etale, and the branch
values of the composite lie in \(B\). The second
cover has genus-zero completion.

The first cover in (FF.48.7) is Galois, with its
eighteen deck transformations
\[
            P\longmapsto \epsilon P+t,
             \qquad t\in A[3],\quad\epsilon\in\{1,-1\}.
                                                        \tag{FF.48.8}
\]
They are distinct and their number equals the degree
\(2\cdot9\) of the map. Its group is
\(G=A[3]\rtimes\{\pm1\}\). The intermediate
cover \(V\) is the quotient by
\(H=\{\pm1\}\). The core of \(H\) in \(G\)
is trivial: conjugation of \([-1]\) by translation
by \(t\ne0\) gives \(P\mapsto-P+2t\), which
does not belong to \(H\). Thus \(Z\) is the
Galois closure of the genus-zero cover \(V\to X\).
The first part of the proof realizes \(Z\).
Its quotient by the normal translation subgroup
\(A[3]\) is the original cover
\[
               \pi^{-1}(X)\longrightarrow X,
                                                        \tag{FF.48.9}
\]
with quotient map supplied by \([3]\). It is therefore
realized as well. The finite group and its faithful
degree-nine permutation action are also checked in
\path{agent_F/F_metabelian_checks.py}.

\emph{The other cyclic genus-one covers.}
The same argument does not require the deck group to
have order two. Over a finite extension, write a
generator of the cyclic deck group as
\(P\mapsto h(P)+t\) on an elliptic model of \(C\).
If \(h=1\), the group consists of translations and
its quotient has genus one, contrary to (FF.48.1).
Thus \(h-1\) is a nonzero endomorphism of the elliptic
curve, hence a surjective isogeny. The generator has
a fixed point; choose it as origin, after a further
finite extension. The deck group \(H\) now consists
of group automorphisms of \(A\), and
\(A/H=\mathbb P^1\).

Choose an integer \(n\ge3\) such that
\[
                 n^2>\deg(h-1)
                  \quad\text{for every }1\ne h\in H,
                                                        \tag{FF.48.10}
\]
and enlarge the finite constant field to contain
\(A[n]\). The map \([n]\) commutes with \(H\)
and induces a degree-\(n^2\) map \(f_n\) on
\(A/H=\mathbb P^1\). The composite
\(\pi\circ[n]\) is Galois with group
\(A[n]\rtimes H\). The core of \(H\) is trivial.
Indeed, an element \(h\) in that core must remain
in \(H\) after conjugation by every translation
\(t\in A[n]\), so \((h-1)t=0\) for all such
\(t\). If \(h\ne1\), this would put \(n^2\)
distinct points in the kernel of an isogeny of degree
strictly less than \(n^2\), contradicting (FF.48.10).
After removing the branch values and the affine
infinity, the preceding Galois-closure argument applies
unchanged. Quotienting by \(A[n]\) recovers
\(\pi\).

\emph{Every Galois genus-one cover.}
Let now \(\pi:A\to\mathbb P^1\) be a Galois
cover with finite deck group \(G\), and genus one
source. Work after a finite extension so that the
origins, deck transformations and torsion used below
are defined. The translations in \(G\) form a normal
subgroup \(T\). Its quotient
\[
                  \lambda:A\longrightarrow B=A/T
\]
is an isogeny of degree \(d=|T|\). The group
\(G/T\) is cyclic: the linear-part map embeds it in
the finite group of origin-preserving automorphisms
of an elliptic curve, which embeds by its action on
an invariant differential in the multiplicative
constant field. This action on differentials is
faithful in characteristic zero. The residual action
on \(B\) cannot consist only of translations,
since \(B/(G/T)=\mathbb P^1\). Its cyclic generator
therefore has a fixed point. Choose that point as
origin of \(B\), and a preimage as origin of \(A\),
so that \(\lambda\) remains a homomorphism. Put
\(H=G/T\), now acting on \(B\) by group
automorphisms, and let
\(\pi_B:B\to B/H=\mathbb P^1\).

Choose \(N\) divisible by \(d\), large enough
that \(N^2>\deg(h-1)\) for all \(1\ne h\in H\).
The Latt\`es argument shows that
\(\pi_B\circ[N]_B\), on the branch complement,
is the Galois closure of a genus-zero cover and is
therefore realized. The dual isogeny
\(\lambda^\vee:B\to A\) satisfies
\(\lambda\circ\lambda^\vee=[d]_B\). Thus
\[
 \pi\circ\lambda^\vee\circ[N/d]_B
                       =\pi_B\circ[N]_B.
                                                        \tag{FF.48.11}
\]
The original cover \(\pi\) is a subcover of this
realized finite Galois cover. It is realized by the
quotient part of the proof. This establishes the
full genus-one Galois assertion. The construction alone
does not treat every non-Galois genus-one map; the
following argument also settles all degree-three maps.

\emph{Degree-three genus-one covers, without a Galois assumption.}
The same generated category also contains every degree-three
cover with genus-one source. We give the finite monodromy
argument explicitly. It is the genus-one case of the
classical trigonal construction described in
\citet[Section~9.1, Remark~9.13]{FF_RecillasRodriguez2003};
the finite verification below makes this consequence
independent of a convention about the genus in a
moduli-space statement.

Work over \(\overline{\mathbb Q}\). A connected
degree-three cover has geometric monodromy \(C_3\)
or \(S_3\). The first case is already abelian.
In the second case let \(a\) be the number of branch
transpositions and \(b\) the number of branch
three-cycles. Riemann--Hurwitz and the product-one
relation give
\[
          a+2b=6,\qquad a>0\text{ even},\qquad
          (a,b)\in\{(6,0),(4,1),(2,2)\}.
                                                        \tag{FF.48.12}
\]
Let \(q:S_4\twoheadrightarrow S_3\) be the action
on the three partitions of four letters into two
unordered pairs. Its kernel is the Klein four group.
Every branch tuple in (FF.48.12) that generates
\(S_3\) and has product one has a lift through
\(q\) that generates \(S_4\), has product one,
and preserves each permutation index. Here the index
of a permutation on \(d\) letters is \(d\) minus
its number of cycles.

This last assertion is a finite lemma verified by
complete enumeration in the supplied program. The
five nonidentity elements of \(S_3\) give
\(5^4+5^5+5^6=19375\) tuples of the possible
lengths. Filtering by index sum six, product one,
and generation of \(S_3\) leaves respectively
72, 270, and 240 tuples. A transposition has two
index-preserving lifts through \(q\), and a
three-cycle has four. For every one of the 582
retained tuples the program enumerates all 64
possible index-preserving lifts and finds exactly
twelve with product one and generated group
\(S_4\). It verifies the quotient map by all
\(24^2\) group products and saves an explicit
lift for every retained tuple. Thus the finite
assertion checks all cases, not just selected
ramification patterns. The program and certificates
are supplied at the following paths.

\path{agent_F/F_metabelian_checks.py}

\path{agent_F/F_recillas_certificates.json}

Apply the characteristic-zero Riemann existence
theorem to a chosen lift at the actual branch
points of the original cover. Its degree-four
source has
\(2g-2=-2\cdot4+6=-2\), hence genus zero.
Its Galois closure has group \(S_4\). The quotient
through \(q\) has exactly the original
\(S_3\) branch tuple, so the original cubic
cover is the subcover associated with the inverse
image of a point stabilizer in \(S_3\).
It therefore belongs to the category already
realized above.

Here descent from the complex construction retains the
covering diagram, not merely its branch data. Let
\(X=\mathbb P^1_{\overline{\mathbb Q}}\setminus B\),
where \(B\) is the actual finite branch set, enlarged
by the original affine infinity if necessary. The
Riemann existence theorem identifies its complex finite
etale covers with finite topological covering spaces
\citep[Expose~XII, Theoreme~5.1]{FF_SGA1}.
Base change from \(\overline{\mathbb Q}\) to
\(\mathbb C\) is an equivalence on finite etale
covering categories: apply
\citet[Expose~XIII, Proposition~4.6(a)]{FF_SGA1}
with \(k=\overline{\mathbb Q}\), the smooth curve
\(X\), and \(Y=\operatorname{Spec}\mathbb C\).
The required resolution hypothesis holds for curves,
and in characteristic zero the prime-to-characteristic
quotient is the entire fundamental group. Since
\(\pi_1(Y)=1\), the resulting fundamental-group map
is an isomorphism. This equivalence is fully faithful,
so the \(S_4\)-action, quotient diagram and isomorphism
with the original cubic cover all descend to
\(\overline{\mathbb Q}\). They use finitely many
algebraic coefficients and morphisms, and therefore
descend to one finite extension \(E'/E\).
Uniqueness of the scheme realization with the specified
actual point function then descends that realization
to \(E\), as detailed below.

\emph{Descent and the fibre rotations.}
All the constructions just used occur after one finite
Galois extension \(L/E\). The resulting realization
of the original etale object is unique from its
geometric point function. That function is
\(G_E\)-equivariant, so the realization is fixed by
\(\operatorname{Gal}(L/E)\) and descends to \(E\).
If the originally specified etale restriction omits
further points of \(C\), use the open-subobject part
of the proof. Naturality follows from the uniqueness
argument at the beginning.

In any connected cyclic torsor case, compare the descended
realization with any chosen isomorphism of the
transported torsor with the original one. Their
composite is an automorphism of the connected cyclic
cover over the base, hence a single global deck
transformation. Accordingly the putative roots of unity
\(c_x\) are constant. This conclusion is about
the slice point function on this object; it does not
identify the induced equivalence on all its further
positive-genus covers.
\end{proof}


The generated subcategory in Claim~\ref{clm:FFlatteskernel}
is in fact not cofinal among all geometric covers.
By the proved fixed-genus composition-factor theorem,
there is a finite set \(\mathcal S_0\) of
nonalternating nonabelian simple groups such that every
simple composition factor of a genus-zero monodromy
group is cyclic, alternating, or belongs to
\(\mathcal S_0\)
\citep[Theorem~1.1, with characteristic zero and genus
zero]{FF_Guralnick2003}.
Every function field obtained by the operations in
Claim~\ref{clm:FFlatteskernel} is contained in a finite
compositum of Galois closures of genus-zero function-field
extensions. The Galois group of that compositum is a
subdirect product of the individual monodromy groups,
so it has no additional simple composition factors.
For completeness, this last assertion follows by induction
on the number of factors: the kernel of projection to
one factor projects normally into each remaining factor;
apply the induction hypothesis to that kernel and then
the composition series of the resulting extension.

Choose a nonalternating nonabelian finite simple group
\(S\notin\mathcal S_0\). Such groups exist, for
example among the infinitely many simple groups of
Lie type. The characteristic-zero fundamental-group
presentation and Riemann existence theorem give an
\(S\)-Galois cover of an open of
\(\mathbb P^1_{\overline{\mathbb Q}}\)
\citep[Expos\'es~XII, Th\'eor\`eme~5.1, and~XIII,
Corollaire~2.12]{FF_SGA1}.
It descends over a finite extension of the constants.
This cover cannot be dominated by an object of the
stated generated subcategory, since then \(S\)
would be a composition factor of one of the preceding
subdirect products. This is a scope restriction on the
closure argument. Noncofinality does not prove that the
remaining automorphism kernel is nontrivial.

\subsection{Vanishing of the Kummer shear on elliptic Galois slices}

\begin{claim}{FF.49}{Proved}\label{clm:FFellipticshear}
Let \(E/B_p\) be finite, let \(\varphi\) be a pinned
relative automorphism of \((\mathbb A^1_E)_{\acute et}\)
with identity actual geometric point function, and let
\[
                  \pi:C\longrightarrow\mathbb P^1_E
                                                        \tag{FF.49.1}
\]
be a finite Galois cover with smooth geometrically
integral genus-one source. On a branch-complement open,
use the unique actual scheme realization of
Claim~\ref{clm:FFlatteskernel} to regard its generic
slice action as an outer automorphism of \(G_{E(C)}\).
For every finite constant extension \(L/E\), every
\(m\ge1\), and every \(f\in L(C)^\times\), this
outer action satisfies
\[
                     \varphi_C^*\kappa_m(f)=\kappa_m(f).
                                                        \tag{FF.49.2}
\]
Equivalently it is the identity on the full abelian
profinite quotient of \(G_{L(C)}\), after every such
constant extension. This includes arbitrary translated
rational functions on the elliptic slice.

Thus every cyclic generic cover of this slice is
preserved as a covering object. Neither its actual
point-function realization nor identity on the entire
elliptic slice topos is asserted by this claim.
\end{claim}

\begin{proof}
We use integral Kummer characters throughout. Since
\(\mu_\infty\subset B_p\), the coefficient sheaves
\(\mathbb Z_\ell(1)\) are constant on every field
under discussion. The relative structural comparison
fixes these coefficients and all characters inflated
from the constant field. Compatible roots of a
nonzero rational function define its character
\(\kappa_{\ell^\infty}(f)\); reduction modulo
\(\ell^r\) is \(\kappa_{\ell^r}(f)\).
Inner changes in the generic action do not affect
these characters.

\emph{Completed-point compatibility on the elliptic slices.}
The geometric local theory used in
Claims~\ref{clm:FFgenerickernel} and
\ref{clm:FFgenuszerokernel} supplies a point
bijection on every smooth projective covering curve,
compatible with finite covering maps. For a
realized genus-one object we must check that this
bijection also agrees at the omitted boundary points
with the extension of its actual scheme realization.
Agreement on the branch complement alone does not
justify that assertion.

Use the dominating Latt\`es cover constructed in
Claim~\ref{clm:FFlatteskernel}. Thus, after finite
constants, the original genus-one Galois cover is a
quotient of
\(Z=B\to\mathbb P^1\) with deck group
\(B[N]\rtimes H\), where \(B\) is elliptic,
\(H\ne1\) consists of origin-preserving group
automorphisms, and \(B/H\) has genus zero. Choose
\(N\) divisible by the degree of the translation
quotient as in that construction, and increase it
so that
\[
              N^2>\sum_{1\ne h\in H}\deg(1-h).
                                                        \tag{FF.49.3}
\]
For every \(t\in B[N]\), put
\(H_t=\tau_tH\tau_{-t}\). The quotient
\(Z/H_t\) has genus zero. These quotients jointly
separate all geometric points of \(Z\). Indeed,
if two distinct points \(z,z'\) have the same
image in the quotient by \(H_t\), then
\[
                  z'=h(z)+(1-h)t
\]
for some \(h\in H\setminus\{1\}\). For each
fixed \(h\), the nonzero isogeny \(1-h\) has
at most \(\deg(1-h)\) solutions \(t\) to this
equation, even before restricting to \(B[N]\).
The union of these exceptional sets has size less
than \(N^2\), proving separation.

Claim~\ref{clm:FFgenuszerokernel} matches the
completed generic point bijection with the scheme
isomorphism on each \(Z/H_t\). On the branch
complement, naturality in
Claim~\ref{clm:FFlatteskernel} identifies the
quotient arrows with the scheme realization on
\(Z\); these scheme diagrams extend to their
smooth projective completions. The generic point
bijections are also compatible with every quotient
arrow. Joint separation therefore makes the
completed point bijection on \(Z\) equal to the
completion of its scheme realization. The finite
surjection from \(Z\) to the original genus-one
cover then gives the same equality there, by
surjectivity on geometric points. The construction
works for every genus-one Galois slice used below.
No assertion about separation of branches by the
projections of a normalized fibre product is used.

\emph{The possible error is geometrically unramified.}
For the moment work geometrically and choose an
origin on \(\bar C\), writing it as an elliptic
curve \(A\). Each finite collection of points,
functions and maps used below is defined after a
finite extension of constants, where the actual
scheme realizations in
Claim~\ref{clm:FFlatteskernel} are available and
natural. Their geometric cohomology groups are
unchanged by these extensions. The preceding
argument identifies their completed geometric
point functions as well.

Every geometric inertia exponent of the generic
base action is one, by
Claim~\ref{clm:FFallprimekernel}. The inertia at a
point of a finite covering curve is a finite-index
subgroup of the corresponding base inertia. After
the geometric point identification supplied by the
actual realization, the restricted exponent is
therefore also one. This holds at every point of
the smooth completion, including points over
infinity. Conjugator choices do not change the
conclusion: geometric curve inertia is procyclic,
and conjugation by its decomposition group acts
through the cyclotomic character, which is trivial
on the constant field here.

It follows that the geometric difference
\[
 b_1(f):=\varphi_C^*\kappa_{\ell^\infty}(f)
                         -\kappa_{\ell^\infty}(f)
       \ \in H^1(A,\mathbb Z_\ell(1))
                                                        \tag{FF.49.4}
\]
is unramified at all points. The indicated group
means geometric cohomology of the proper curve.
Indeed the residue of a Kummer character on point
inertia is the valuation of its function, and all
these residues of the difference vanish. The
resulting unramified character factors through the
proper curve fundamental group. It is an integral
class, not just a rational cohomology class.

\emph{Multiplication covers and their translations.}
For every positive integer \(n\), the map
\(\pi\circ[n]:A\to\mathbb P^1\), after finite
constants, is again a Galois cover with genus-one
source. To check this without assuming that all
deck transformations fix the origin, write each
deck transformation of \(\pi\) as
\(P\mapsto h(P)+t\), where \(h\) is a group
automorphism. It has exactly \(n^2\) lifts through
\([n]\), namely
\(P\mapsto h(P)+s\) with \(ns=t\). These give
\(n^2\deg(\pi)\) distinct automorphisms of the
composite, equal to its degree. All are defined over
one finite constant extension. In particular the
deck transformations include all translations
\(\tau_t\), \(t\in A[n]\).

Denote this covering object by \(A_n\), retaining
its map to the original rational base, and denote
its induced generic action by \(\varphi_n^*\).
Claim~\ref{clm:FFlatteskernel} realizes \(A_n\),
its deck arrows, and the map
\([n]:A_n\to A_1\) naturally on their branch
complements. Hence on cohomology
\[
 \varphi_n^*[n]^*=[n]^*\varphi_1^*,\qquad
 \varphi_n^*\tau_t^*=\tau_t^*\varphi_n^*.
                                                        \tag{FF.49.5}
\]
These equations are not an assumed extension to
unrelated rational subfields: all their arrows are
over the original base. The possible inner
ambiguities in the generic group comparisons are
irrelevant on \(H^1\).
Define \(b_n(h)\) as in (FF.49.4). The inertia
argument makes every such class lie in
\(H^1(A,\mathbb Z_\ell(1))\).
Translations act as the identity on this group,
and \([n]^*\) acts as multiplication by \(n\).
One may see both statements from
\(A(\mathbb C)=\mathbb C/\Lambda\): translations
are homotopic to the identity and multiplication
acts by \(n\) on the lattice. The algebraic and
profinite identifications follow from the
characteristic-zero comparison theorem
\citep[Expose~XII, Theoreme~5.1]{FF_SGA1}.
The group is free of rank two over \(\mathbb Z_\ell\).

\emph{An integral averaging identity.}
Suppose first that the principal divisor
\[
                       D=\operatorname{div}(f)
                              =\sum_i m_i[P_i]
\]
has some coefficient equal to \(1\) or \(-1\).
Choose \(Q_i\) with \(nQ_i=P_i\). Since \(D\)
is principal, \(\sum_i m_i=0\) and
\(\sum_i m_iP_i=0\) in the elliptic group.
Consequently \(s=\sum_i m_iQ_i\) belongs to
\(A[n]\). Adjust the chosen \(Q_i\) at a
coefficient \(\pm1\) by the appropriate element
of \(A[n]\) so that this sum becomes zero.
The divisor
\[
                         D_0=\sum_i m_i[Q_i]
\]
now has degree zero and zero sum in the elliptic
group, so it is principal, say
\(D_0=\operatorname{div}(h)\).
All these choices use only finitely many algebraic
constants. Divisors then give
\[
          f\circ[n]=c\prod_{t\in A[n]}\tau_t^*h
                  \qquad(c\in\overline{\mathbb Q}^{\times}).
                                                        \tag{FF.49.6}
\]
A constant has zero geometric Kummer character.
Apply the difference operator on \(A_n\) to
(FF.49.6), and use (FF.49.5). It gives the equality
in integral geometric cohomology
\[
           n\,b_1(f)
              =[n]^*b_1(f)
              =\sum_{t\in A[n]}\tau_t^*b_n(h)
              =n^2b_n(h).
                                                        \tag{FF.49.7}
\]
Torsion-freeness permits cancellation of \(n\).
Taking \(n=\ell^r\) for every \(r\ge1\) shows
that \(b_1(f)\) belongs to
\(\ell^r H^1(A,\mathbb Z_\ell(1))\) for every
\(r\), and hence is zero.

This proves the assertion for every rational
function, not just the ones with a coefficient
\(\pm1\). Given \(f\), choose a function \(q\)
with a simple zero outside the support of
\(\operatorname{div}(f)\), and with no cancellation
there. For example, take a degree-two function on
the elliptic curve and subtract a constant avoiding
its finitely many branch values and the images of
that support. Both \(q\) and \(fq\) have such
a simple zero. The proved case and additivity give
\(b_1(f)=b_1(fq)-b_1(q)=0\).
Thus the geometric restriction of every error in
(FF.49.4) vanishes, for every \(\ell\).

\emph{Removing the arithmetic error by the norm.}
Return to an arbitrary finite \(L/E\) and
\(f\in L(C)^\times\). Put
\(F=L(T)\), \(K=L(C)\), and
\(d=[K:F]\). The full integral difference
\[
       \delta=\varphi_C^*\kappa_{\ell^\infty}(f)
                         -\kappa_{\ell^\infty}(f)
          \in\operatorname{Hom}_{\mathrm{cts}}
                     (G_K,\mathbb Z_\ell(1))
                                                        \tag{FF.49.8}
\]
has zero restriction to the geometric kernel by
the preceding argument. It therefore factors as
an inflated character
\(\gamma:G_L\to\mathbb Z_\ell(1)\).
This factorization is at the full compatible
integral level; no conclusion is drawn merely
from separate finite-level classes.

Corestriction commutes with the induced generic
automorphisms, because the actual slice realization
identifies the covering map with \(\pi\). The
Kummer norm identity and
Claim~\ref{clm:FFallprimekernel} on the rational
base imply
\[
 \operatorname{cor}_{K/F}(\delta)
       =\varphi^*\kappa_{\ell^\infty}(N_{K/F}f)
                   -\kappa_{\ell^\infty}(N_{K/F}f)=0.
                                                        \tag{FF.49.9}
\]
Since \(\delta\) is the restriction of the
character inflated from \(\gamma\), the
restriction--corestriction identity identifies the
left side with \(d\gamma\), inflated to \(G_F\).
Inflation is injective, and
\(\operatorname{Hom}_{\mathrm{cts}}
(G_L,\mathbb Z_\ell(1))\) is torsion-free.
Therefore \(\gamma=0\), and \(\delta=0\).
This norm argument needs no rational point on
\(C_L\), so it also handles genus-one torsors
before choosing an origin.

Reducing the equality modulo every \(\ell^r\)
and taking the product over primes proves
(FF.49.2). Every finite Kummer character comes
from some \(f\) by Kummer theory, and
\(\mu_m\subset L\), so these equalities are
exactly identity on the full abelian profinite
quotient. This conclusion determines cyclic
covering objects; it does not supply the further
actual point comparisons required to identify
the nonabelian generic action or the whole slice
topos.
\end{proof}


% Root and geometric_routes independently approved the complete proof.
% Integrated as FF.50; retained separately as the reviewed source.
\subsection{Integral norms and a higher-genus slice}

\begin{claim}{FF.50}{Proved: Kummer rigidity for an explicit genus-two slice}
\label{clm:FFhighergenusnorm}
Retain the pointwise pinned automorphism hypotheses of
Claims~\ref{clm:FFellipticshear} and the preceding
genus-zero and Lattes results, over a finite extension
\(E/B_p\). Let \(C\) be the smooth projective curve
\[
 C:\quad y^2=(u^2-1)(u^2-2)(u^2-3),\qquad T=u^2.
                                                        \tag{FF.50.1}
\]
This is a genus-two Galois cover of \(\mathbb P^1_T\)
with group \(V_4\). Its transported actual point
function, including every completed point, is realized
by a unique scheme isomorphism over \(\mathbb P^1_T\).
Using that isomorphism to identify its generic action,
the automorphism fixes
\[
                    H^1(G_{L(C)},\mu_m)
                                                        \tag{FF.50.2}
\]
for every finite extension \(L/E\) and every positive
integer \(m\).

More generally, the same Kummer conclusion holds for a
realized smooth proper slice \(C\) whenever it has
coherently realized connected finite etale covers
\(D_s\to C\), with completed-point compatibility,
whose pushforward images in integral geometric first
homology are contained in
\(\ell^sH_1(C,\mathbb Z_\ell)\), for arbitrarily
large \(s\) and each prime \(\ell\).
\end{claim}

The completed-point hypothesis in the general criterion
is substantive. Preserving the isomorphism class of an
Abel cover does not establish it. The specific curve
(FF.50.1) meets the hypothesis by its two elliptic
quotients, as proved below. This is a positive-genus
Kummer constraint; it does not identify all further
nonabelian covers or prove reconstruction over \(B_p\).

\begin{proof}
\emph{1. The integral norm criterion.}
Work first over an algebraic closure of the constants.
For every coherently realized curve \(X\), use its
actual-point scheme isomorphism to identify the
transported generic group with \(G_{\overline E(X)}\).
The induced action fixes each geometric place and its
inertia exponent, the latter by restriction of the
base exponent one. Hence for each rational function
\(f\ne0\) its Kummer error is unramified:
\[
 b_X(f):=\alpha_X\kappa_\ell(f)-\kappa_\ell(f)
             \in H^1(X_{\overline E},\mathbb Z_\ell(1)).
                                                        \tag{FF.50.3}
\]
The localization sequence proves this integrally:
the residues at all places agree, so their difference
lies in proper first cohomology.

Let \(\rho:D\to C\) be any connected finite etale
cover in the stated coherent family. The norm
\(\overline E(D)^\times\to\overline E(C)^\times\)
is surjective. Here is the dimension proof needed for
that assertion. The function field of a curve over an
algebraically closed field is \(C_1\): for a
homogeneous form of degree \(d\) in \(N>d\) variables,
choose the variables in \(H^0(C,\mathcal O_C(mA))\),
where \(A\) is a fixed positive divisor. After clearing
the fixed coefficient poles, the equation is a system
of at most \(dm\deg A+O(1)\) homogeneous equations in
\(Nm\deg A+O(1)\) scalar variables, by Riemann--Roch.
For sufficiently large \(m\) there are fewer equations
than variables, so the projective dimension theorem
gives a nonzero common zero over the algebraically
closed constant field. Apply this to
\[
               N_{\overline E(D)/\overline E(C)}(v)-f z^d,
               \qquad d=[\overline E(D):\overline E(C)].
                                                        \tag{FF.50.4}
\]
It has degree \(d\) in \(d+1\) variables. A nonzero
solution has \(z\ne0\), since the norm of a nonzero
field element is nonzero. Dividing by \(z\) supplies
\(h\in\overline E(D)^\times\) with \(N(h)=f\).
All these data descend to a finite constant extension.
This is the usual Tsen argument, included here in the
precise form used.

Functoriality of corestriction and of the coherent
scheme realizations gives
\[
                   b_C(f)=\operatorname{Tr}_\rho b_D(h).
                                                        \tag{FF.50.5}
\]
The transfer direction matters: under integral
Poincare duality for the two compact curves,
\(\operatorname{Tr}_\rho:H^1(D,\mathbb Z_\ell(1))
\to H^1(C,\mathbb Z_\ell(1))\) corresponds to
\(\rho_*:H_1(D,\mathbb Z_\ell)\to
H_1(C,\mathbb Z_\ell)\). Thus the hypothesis on
the image of homology makes (FF.50.5) divisible by
\(\ell^s\) in the fixed free \(\mathbb Z_\ell\)-module
\(H^1(C,\mathbb Z_\ell(1))\). Varying \(s\) forces
\(b_C(f)=0\).

For the full Abel mod-\(n\) cover of \(C\), the
homology image is exactly \(nH_1(C,\mathbb Z)\):
its fundamental subgroup is the kernel of
\(\pi_1(C)\to H_1(C,\mathbb Z/n)\).
Every element of its homology image is divisible by
\(n\), and the \(n\)-th power of a loop realizing
any integral homology basis element gives the reverse
inclusion. Therefore these Abel covers would satisfy
the transfer requirement if their completed-point
compatibility were known. No such compatibility for
arbitrary curves is being assumed.

\emph{2. The three quotients of the specified curve.}
The six roots of the polynomial in (FF.50.1) are
distinct, so its double cover of \(\mathbb P^1_u\)
has genus two. The independent quadratic equations
\(u^2=T\) and \(y^2=(T-1)(T-2)(T-3)\) give
the \(V_4\)-cover over \(T\). Its three quadratic
quotients are
\[
 \begin{gathered}
 \mathbb P^1_u,\qquad A_1:\ y^2=(T-1)(T-2)(T-3),\\
 A_2:\ z^2=T(T-1)(T-2)(T-3),\qquad z=uy.
 \end{gathered}
                                                        \tag{FF.50.6}
\]
The latter two are elliptic curves, each with origin
at its branch point \(T=1\). Their maps to
\(\mathbb P^1_T\) are the quotients by \(\pm1\).
The point \(P_0=(u,y)=(1,0)\) on \(C\) maps to
both origins.

The previous genus-zero and Lattes results realize
all three quotient objects naturally, including
their completed-point maps. The three maps from
\(C\) jointly separate every geometric completed
point. For a group-theoretic verification, a
geometric point has cyclic inertia \(D\leq V_4\),
and equality of its images in a quotient by one
of the three order-two subgroups \(H_i\) means
that the relative translating element lies in
\(H_iD\). One has
\(\bigcap_iH_iD=D\): this is immediate if
\(D=1\), and if \(|D|=2\) one of the \(H_i\)
is exactly \(D\). Thus no distinct completed
points have all three images equal. Naturality
on the branch-complement opens, the unique
extension of a curve isomorphism to its smooth
completion, and this separation show that the
realization of \(C\) agrees with its entire
completed-point bijection. This step does not
assume that arbitrary normalized products
separate boundary branches.

\emph{3. Elliptic division covers provide the required
homology divisibility.}
Let \(\rho_i:C\to A_i\) be the two elliptic quotient
maps, based at \(P_0\). They induce homomorphisms
\[
 \varphi:J_C\longrightarrow A_1\times A_2,
 \qquad
 \psi:A_1\times A_2\longrightarrow J_C,
                                                        \tag{FF.50.7}
\]
where \(\varphi\) is the product of the pushforwards
and \(\psi\) the sum of the pullbacks. Their composite
is
\[
                              \psi\varphi=[2]_{J_C}.
                                                        \tag{FF.50.8}
\]
Indeed if \(\tau(u,y)=(-u,y)\) and
\(\iota(u,y)=(u,-y)\), the two degree-two norm
identities give the endomorphisms \(1+\tau\) and
\(1+\iota\tau\) on \(J_C\). The hyperelliptic
involution \(\iota\) acts as \(-1\), and their
sum is \(2\). In particular \(\varphi\) has
finite kernel killed by two and is an isogeny,
since both dimensions are two.

For \(n\geq1\), take the connected component
containing \((P_0,0,0)\) of
\[
 D_n=C\mathbin{\times_{A_1\times A_2,[n]}}
                                  (A_1\times A_2).
                                                        \tag{FF.50.9}
\]
It is a finite etale cover of \(C\). Its two elliptic
factors, considered over \(\mathbb P^1_T\), are
the Galois genus-one Lattes covers
\(A_i\xrightarrow{[n]}A_i\to\mathbb P^1_T\),
after a finite constant extension making the torsion
rational. They and all their displayed arrows have
the coherent completed-point realizations already
proved. Ordinary fiber products and the selected
component now give that same compatibility for
\(D_n\). This time there is no normalization
ambiguity: (FF.50.9) is etale over the smooth curve
\(C\), hence smooth, and its geometric points are
exactly the indicated tuples. The projection to
\(C\) and the two elliptic coordinates therefore
separate all its completed points.

Let \(\Lambda=H_1(C(\mathbb C),\mathbb Z)\) and
\(\Lambda_A=H_1((A_1\times A_2)(\mathbb C),\mathbb Z)\),
using any embedding of the finite coefficient field
into \(\mathbb C\). The covering subgroup in
(FF.50.9) consists exactly of the loops whose
\(\varphi_*\)-image lies in \(n\Lambda_A\).
Consequently each element \(v\) in the image of
\(H_1(D_n,\mathbb Z)\to\Lambda\) satisfies
\[
                 \varphi_*v\in n\Lambda_A,
                 \qquad 2v=\psi_*\varphi_*v\in n\Lambda.
                                                        \tag{FF.50.10}
\]
For \(n=\ell^s\), this puts the image in
\(\ell^s\Lambda\) if \(\ell\ne2\), and in
\(2^{s-1}\Lambda\) if \(\ell=2\). Integral
comparison gives the same assertion for etale
\(\mathbb Z_\ell\)-homology. The loss of one power
at two is uniform and harmless. Step~1 therefore
proves zero geometric Kummer error on \(C\) at
every prime, including two.

\emph{4. Arithmetic constants and finite coefficients.}
Fix a finite \(L/E\) and \(f\in L(C)^\times\).
Since \(\mu_\infty\subset L\), its compatible
\(\ell\)-adic Kummer class is a genuine continuous
character. The difference under the identified
generic automorphism is zero on the geometric
kernel by the preceding steps, hence is the
inflation of a character
\(\gamma:G_L\to\mathbb Z_\ell(1)\).
Corestriction to \(L(T)\) is zero, since rational
Kummer classes are fixed. On the inflated character
it is multiplication by \([L(C):L(T)]=4\).
The group \(\operatorname{Hom}_{\mathrm{cont}}
(G_L,\mathbb Z_\ell(1))\) is torsion-free, so
\(\gamma=0\), including when \(\ell=2\).
Reducing the integral Kummer character gives every
finite Kummer class, and combining the primes
proves (FF.50.2).
For the general criterion the identical argument uses
the positive degree \([L(C):L(T)]\) of its finite
slice in place of four; no divisibility assumption
on that degree is needed.
\end{proof}

\paragraph{Exact limitation and the next obstruction.}
The construction works because the Jacobian of this
curve admits an isogeny to elliptic quotients which
are themselves realized covers of the same rational
base. More generally the argument applies whenever
finitely many such quotient maps induce an injection
on rational first homology and their division
fiber products have completed-point compatibility;
a fixed isogeny denominator only changes the bounded
loss in (FF.50.10). For a Jacobian with an absolutely
simple factor of dimension greater than one, these
elliptic maps do not provide that injection. The
obstruction cannot be removed merely by taking a
finite cover \(D\to C\): norm composed with
pullback is multiplication by its degree on
\(J_C\), so \(J_C\) remains an isogeny factor
of \(J_D\). If the latter were built entirely
from elliptic factors, the former would be as well.
The
fixed-prime configuration theorem establishes a
completion topology, not the required actual point
maps on these higher-dimensional Abel covers.
Thus the remaining compatibility problem has not
been solved for an arbitrary positive-genus slice.
